\documentclass[10pt,a4paper]{article}

% Importação do preâmbulo customizado
% =============================================================================
% PREÂMBULO INTEGRAL CUSTOMIZADO
% =============================================================================

% --- Pacotes de Geometria e Tipografia ---
\usepackage[a4paper,lmargin=1.5cm,rmargin=1.5cm,tmargin=1.2cm,bmargin=1.5cm]{geometry}
\usepackage[brazilian]{babel}
\usepackage[T1]{fontenc}
\usepackage{microtype}
\usepackage{xspace}

% --- Macros de Tipografia (Corrigindo T1) ---
\newcommand{\f}{\textordfeminine\xspace} % Para usar ª
\newcommand{\m}{\textordmasculine\xspace} % Para usar º

% --- Matemática e Símbolos ---
\usepackage{amsmath,amsfonts,amssymb,mathtools,mathrsfs,amsthm}
\usepackage{stmaryrd}
\usepackage{latexsym}
\usepackage{esint}

% --- Tikz e Referências ---
\usepackage{csquotes} % Elimina o aviso do babel
\usepackage{caption}
\usepackage{tikz}
\usetikzlibrary{arrows.meta, positioning, shapes.geometric, decorations.pathreplacing}
\usepackage{pgfplots}
\pgfplotsset{compat=1.18}
\tikzset{
	passo/.style={
		rectangle, 
		draw, 
		align=center, % Alinhamento centralizado para evitar hifenização forçada
		text width=6cm, % Largura otimizada para caber os nomes sem estourar
		rounded corners, 
		minimum height=1.2cm, 
		font=\small
	},
	arrow/.style={-Stealth, thick}
}
% Referências Bibliográficas
\usepackage[backend=biber, style=abnt]{biblatex} % Recomendado usar biber para bibliografias modernas
%\addbibresource{referencias.bib}

% --- Elementos Visuais e Layout ---
\usepackage[most]{tcolorbox}
\usepackage{enumitem}
\setlist[itemize]{noitemsep, topsep=0pt, parsep=0pt, partopsep=0pt}
\usepackage[normalem]{ulem}
\usepackage[Smaller]{cancel}
\usepackage{xcolor}
\usepackage[nodisplayskipstretch]{setspace}
\usepackage{etoolbox} % ESSENCIAL para os interruptores de blocos

% --- Navegação e Links ---
\usepackage[colorlinks=true, linkcolor=black, urlcolor=blue, hypertexnames=false]{hyperref}
\usepackage{bookmark}

% --- Configurações Globais de Estilo ---
\setstretch{1.3}
\setlength{\parindent}{0pt}
\relpenalty=10000
\binoppenalty=10000
\hbadness=10000
\everymath{\displaystyle}
\pagestyle{empty}
\newcommand{\linhaestilizada}{ % Linha centralizada verticalmente com diamond
	\par\noindent\makebox[\textwidth]{\rule[0.5ex]{0.47\textwidth}{0.4pt}\hfill$\diamond$\hfill\rule[0.5ex]{0.47\textwidth}{0.4pt}}\par
}
\setlength{\overfullrule}{5pt} % Desenha uma barra de 5pt no final de hboxes para erros de estrapolação de linha horizontal

% =============================================================================
% SISTEMA DE CONTROLE DE VISIBILIDADE (INTERRUPTORES)
% =============================================================================
% Altere 'true' para 'false' para omitir o bloco no PDF final sem apagar o código.

\newbool{showanalise}
\setbool{showanalise}{true}

\newbool{showsintese}
\setbool{showsintese}{true}

\newbool{showaprofundamento} 
\setbool{showaprofundamento}{true}

\newbool{showtheorem}
\setbool{showtheorem}{true}

\newbool{showproof}
\setbool{showproof}{true}

\newbool{showtextoauxiliar}
\setbool{showtextoauxiliar}{true} 

% =============================================================================
% COMANDOS ESTRUTURAIS DO PROJETO
% =============================================================================

% Valor padrão seguro para compilação via pacote subfiles
\providecommand{\listid}{L0}

% Define o ambiente theorem* como não numerado
% --- Ambientes Condicionais de Teoremas e Provas ---
\makeatletter
\@ifundefined{theorem*}{%
	\theoremstyle{definition} 
	\newtheorem*{theorem*}{Teorema}%
}{}
\makeatother

\newenvironment{teorema}{%
	\ifbool{showtheorem}{%
		\begin{theorem*}%
		}{%
			\setbox0=\vbox\bgroup
		}%
	}{%
		\ifbool{showtheorem}{%
		\end{theorem*}%
	}{%
		\egroup
	}%
}

\newenvironment{prova}{%
	\ifbool{showproof}{%
		\begin{proof}[\normalfont\underline{Demonstração}:] 
		}{%
			\setbox0=\vbox\bgroup
		}%
	}{%
		\ifbool{showproof}{%
		\end{proof}%
	}{%
		\egroup
	}%
}

% --- Comando Customizado: Texto Auxiliar Solto Condicional ---
\newcommand{\textoauxiliar}[1]{%
	\ifbool{showtextoauxiliar}{%
		\par\vspace{3pt}%
		#1% Renderiza apenas o texto solto, sem cabeçalho
		\par\vspace{3pt}%
	}{}%
}

\newcounter{numquestao}

% =============================================================================
% SISTEMA INTELIGENTE DE SELEÇÃO DE QUESTÕES
% =============================================================================

% 1. Controle de Estados
\newbool{showcurrentQ}
\setbool{showcurrentQ}{true} 

\newbool{modoisolado}
\setbool{modoisolado}{false} 

% 2. Comandos de Seleção Otimizados para Listas (Separados por vírgula)
\newcommand{\ocultarquestao}[1]{%
	\foreach \q in {#1} {%
		\csgdef{hideQ\q}{true}% Usa csgdef global para escapar o escopo do foreach do tikz
	}%
}

\newcommand{\mostrarapenasquestao}[1]{%
	\setbool{modoisolado}{true}%
	\foreach \q in {#1} {%
		\csgdef{showOnlyQ\q}{true}%
	}%
}

% 3. Sobrescrita do Comando Questão
\newcommand{\questao}[1]{%
	\stepcounter{numquestao}% Avança o contador SEMPRE, mantendo a numeração original 
	%
	% Lógica de Visibilidade Dinâmica
	\setbool{showcurrentQ}{true}%
	\ifbool{modoisolado}{%
		\ifcsdef{showOnlyQ\arabic{numquestao}}{}{\setbool{showcurrentQ}{false}}%
	}{%
		\ifcsdef{hideQ\arabic{numquestao}}{\setbool{showcurrentQ}{false}}{}%
	}%
	%
	% Renderização Condicional
	\ifbool{showcurrentQ}{%
		\hypertarget{dest.\listid.\arabic{numquestao}}{}%
		\bookmark[level=1, dest={dest.\listid.\arabic{numquestao}}]{Questão \arabic{numquestao}}%
		\begin{tcolorbox}[
			enhanced, sharp corners, boxrule=0.5pt,
			colback=gray!6!white, colframe=black,         
			left=1pt, right=1pt, top=1pt, bottom=1pt,
			width=1\textwidth, before skip=3pt, after skip=3pt,
			fontupper=\normalfont
			]%
			\mbox{\textbf{Questão \arabic{numquestao}.}}\hspace{0.3em}\ignorespaces#1%
		\end{tcolorbox}%
	}{}%
}

% 4. Condicionais em cascata para os blocos da resolução 
\newcommand{\analise}[1]{%
	\ifbool{showcurrentQ}{%
		\ifbool{showanalise}{%
			\par\vspace{5pt}\textbf{\underline{Análise e Estratégia}}:\par 
			#1\par
		}{}%
	}{}%
}

\newcommand{\solucao}[1]{%
	\ifbool{showcurrentQ}{%
		\par\vspace{5pt}
		\textbf{\underline{Solução}}:\par
		#1 \hfill$\blacksquare$
		\par
	}{}%
}

\newcommand{\solucaoalt}[1]{%
	\ifbool{showcurrentQ}{%
		\ifbool{showanalise}{%
			\par\vspace{5pt}
			\textbf{\underline{Solução Alternativa}}:\par 
			#1 \hfill$\blacksquare$
			\par
		}{}%
	}{}%
}

\newcommand{\sintese}[1]{%
	\ifbool{showcurrentQ}{%
		\ifbool{showsintese}{%
			\par\vspace{5pt}\textbf{\underline{Síntese da Construção}}:\par 
			#1\par
		}{}%
	}{}%
}

\newcommand{\aprofundamento}[1]{%
	\ifbool{showcurrentQ}{%
		\ifbool{showaprofundamento}{%
			\linhaestilizada
			\par\noindent\textbf{\underline{Aprofundamento Teórico}}\par
			#1
			\par
			\linhaestilizada
		}{}%
	}{}%
}

% =============================================================================
% OPERADORES E NOTAÇÃO DE ANÁLISE
% =============================================================================
\newcommand{\R}{\mathbb{R}}
\newcommand{\norm}[1]{\|#1\|}
\newcommand{\inner}[2]{\langle #1, #2 \rangle}
\newcommand{\Sobolev}[2]{W^{#1, #2}}
\newcommand{\Hspace}[1]{H^#1}
\DeclareMathOperator{\supp}{supp}
\DeclareMathOperator{\dist}{dist}
\newcommand{\FF}{\mathcal{F}} 
\newcommand{\invFF}{\mathcal{F}^{-1}}

\begin{document}
\nocite{*}
	% =============================================================================
	% CAPA MINIMALISTA
	% =============================================================================
	\begin{titlepage}
		\begin{center}
			\vspace*{4cm}
			
			\rule{\linewidth}{0.5mm} \\[0.4cm]
			{ \huge \bfseries Caminhos Analíticos e Construção de Soluções em EDP II \\[0.4cm] }
			\rule{\linewidth}{0.5mm} \\[2cm]
			
			\large
			\textbf{Anselmo Ribeiro Lopes}\\
			
			\vfill
			
			\large
			Campina Grande - PB\\
			2026
		\end{center}
	\end{titlepage}
	
	% =============================================================================
	% PREFÁCIO REFLEXIVO E METODOLÓGICO
	% =============================================================================
	\section*{Prefácio}
	\addcontentsline{toc}{section}{Prefácio}
	
	Este documento é um registro pessoal de estudos e um processo de aprofundamento na disciplina de \textbf{Equações Diferenciais Parciais II (EDP II)} do semestre \textbf{2026.1}, no \textbf{Programa de Pós-Graduação em Matemática (PPGMAT)} da \textbf{UFCG}, ministrada pelo \textbf{Professor Marco Aurelio Souto}.
	
	Mais do que a simples resolução de exercícios, busco aqui visualizar a construção de cada solução, refletindo sobre o 'porquê' e o 'para que' de cada etapa. Compreendo que cada questão, para além do resultado final, carrega fatos matemáticos essenciais e provoca reflexões que são fundamentais para a construção do pensamento matemático. A estruturação deste trabalho segue rigorosamente o fluxo das \textbf{listas de exercícios da disciplina}, fundamentando-se nas discussões em sala de aula e nas referências bibliográficas de \textbf{Lawrence C. Evans, Haim Brezis e Gilbarg \& Trudinger}. Para facilitar a navegação e a compreensão conceitual, optei por organizar o conteúdo em divisões temáticas. Essas divisões não rompem a sequência lógica dos exercícios, mas servem como marcos que agrupam problemas com naturezas e ferramentas semelhantes, permitindo uma visão mais clara da evolução dos temas da disciplina.
	
	Para que este método de estudo seja claro, adotei a seguinte organização para cada problema. O fluxo de pensamento é desenhado como uma escada de complexidade, onde cada passo é fundamental para o próximo:
	
	\begin{itemize}
		\item \textbf{Enunciado da Questão:} A transcrição literal do problema, servindo como o ponto de partida e a definição do desafio.
		\item \textbf{Análise e Fundamentação Teórica:} Identificação dos fatos matemáticos, a base teórica necessária e a reflexão inicial.
		\item \textbf{Estratégia para Solução:} O roteiro lógico e o plano de ataque para a resolução.
		\item \textbf{Solução:} O desenvolvimento formal e exaustivo, sem omissões, utilizando linguagem natural e referências cruzadas.
		\item \textbf{Síntese da Construção:} Um fechamento reflexivo que apresenta a lógica da solução e a lição extraída.
	\end{itemize}
	% --- DIAGRAMA TIKZ: A ESCADA DA CONSTRUÇÃO ANALÍTICA (LARGURA OTIMIZADA) ---
	\begin{center}
			\begin{tikzpicture}[
				scale=0.7,
				node distance=1.2cm,
				arrow/.style={-Stealth, thick}
				]
				% Definição dos passos com xshift reduzido para diminuir a largura total
				\node (enun) [passo, font=\bfseries, fill=blue!5] {Enunciado da Questão};
				\node (anal) [passo, font=\bfseries, fill=blue!15, xshift=1.5cm, yshift=1.2cm] {Análise e Fundamentação Teórica};
				\node (estr) [passo, font=\bfseries, fill=blue!25, xshift=3cm, yshift=2.4cm] {Estratégia para Solução};
				\node (solu) [passo, font=\bfseries, fill=blue!35, xshift=4.5cm, yshift=3.6cm] {Solução};
				\node (sint) [passo, font=\bfseries, fill=blue!45, xshift=6cm, yshift=4.8cm] {Síntese da Construção};
				
				% 1. SETA ROXA DE SUBIDA (Lado Esquerdo)
				% Afastada 1cm para a esquerda para manter a harmonia com a nova largura
				\draw [arrow, blue!45, line width=2pt] 
				([xshift=-1.7cm]enun.west) -- ([xshift=-1.7cm]sint.west) 
				node[midway, above left, font=\small\bfseries, blue!45] {Caminho da Construção};
				
				% 2. SETA DE RETORNO (Lado Direito)
				% Reta, afastada 1cm para a direita
				\draw [arrow, gray, line width=1pt] 
				([xshift=1.7cm]sint.east) -- ([xshift=1.7cm]enun.east)
				node[midway, below right, font=\small\bfseries, gray] {Reflexão Analítica};
				
			\end{tikzpicture}
	\end{center}
	A intenção é explicitar o método e a lógica empregada, transformando a resolução em um exercício de reflexão analítica. Este material não pretende esgotar o assunto nem estabelecer verdades definitivas, mas sim servir como um diário de bordo --- um caminho de pensamento para compreender a estrutura das EDPs e a beleza da sua construção lógica.
	
	\vspace{1cm}
	\begin{center}
		\small
		Este material é distribuído sob a \textbf{Licença CC BY-NC-SA 4.0}. \\
		Sugestões, críticas ou correções são muito bem-vindas via: \href{https://forms.gle/UcBrKWL7YxdXCX7y9}{\textbf{\color{blue!70!black}Formulário de Observações}}
	\end{center}
	\begin{flushright}
		\textit{Campina Grande, \today} \\
		\textbf{Anselmo Ribeiro Lopes}
	\end{flushright}
	
	\newpage
	
	% =============================================================================
	% SUMÁRIO
	% =============================================================================
	\tableofcontents
	\newpage
	
	% =============================================================================
	% CORPO DO TRABALHO (FLUXO CONTÍNUO)
	% =============================================================================
	
	\section{Equação do Transporte}
	\def\listid{L1} 
	\setcounter{numquestao}{0} 
	\documentclass[cs_edp_2026.1.tex]{subfiles} % Aponta para o pai

\begin{document}
% CONFIGURAÇÃO PARA COMPILAÇÃO INDIVIDUAL
\ifSubfilesClassLoaded{
	\def\listid{L1}
	\setcounter{numquestao}{0}
	\section*{Equação do Transporte}
}{}

\questao{Resolva a seguinte equação linear de primeira ordem
	$$ \begin{cases} u_t + a \cdot \nabla u = u, & \text{em } \mathbb{R}^n \times \mathbb{R} \\ u(x,0) = g(x), & \text{quando } x \in \mathbb{R}^n. \end{cases} $$
	Que condições devemos impôr à função $g$ para obtermos soluções clássicas?}

\solucao{
	\textbf{Análise do Enunciado:} 
	O problema apresenta uma EDP linear de primeira ordem com um termo de fonte dependente de $u$. Buscamos a solução $u(x,t)$ dado um perfil inicial $g(x)$. O objetivo é determinar a solução e a regularidade necessária de $g$ para que $u$ seja de classe $C^1$.
	
	\textbf{Fundamentação Teórica:} 
	Utilizaremos o Método das Características (Evans, Cap. 2.2). Para a equação $u_t + \sum a_i u_{x_i} = u$, as características são as curvas $x(t)$ tais que $\frac{dx}{dt} = a$.
	
	\textbf{Desenvolvimento Formal:} 
	Definimos o sistema de equações diferenciais ordinárias (EDOs) para as características:
	$$ \frac{dx}{dt} = a \implies x(t) = at + x_0 $$
	$$ \frac{du}{dt} = u $$
	Resolvendo a segunda EDO:
	$$ \ln(u) = t + C \implies u(t) = u(0)e^t $$
	Utilizando a condição inicial $u(x_0, 0) = g(x_0)$, temos:
	$$ u(x(t), t) = g(x_0)e^t $$
	Como $x_0 = x - at$, substituímos para obter a solução geral:
	$$ u(x, t) = g(x - at)e^t $$
	Para que a solução seja clássica ($u \in C^1$), devemos exigir que a função $g$ seja de classe $C^1(\mathbb{R}^n)$, visto que as derivadas parciais de $u$ envolverão as derivadas de $g$ via regra da cadeia.
}

\end{document}
	
	\section{Equação de Laplace}
	\def\listid{L2} 
	\setcounter{numquestao}{0} 
	\questao{Se $U$ é um domínio limitado que satisfaz o teorema do divergente, mostre que existe no máximo uma solução para o problema de Dirichlet:
	$$ \begin{cases} \Delta u = f(x), & x \in U \\ u = g(x), & x \in \partial U \\ u \text{ em } C^2(U) \cap C^1(\overline{U}). \end{cases} $$}

\solucao{
	\textbf{Análise do Enunciado:} 
	O objetivo é provar a unicidade da solução para o problema de Dirichlet. Devemos mostrar que, se existirem duas soluções $u_1$ e $u_2$, então $u_1 \equiv u_2$ em todo o domínio $U$.
	
	\textbf{Fundamentação Teórica:} 
	Utilizaremos o método da energia baseado na Primeira Identidade de Green (Evans, Cap. 2.2). A estratégia consiste em analisar a função diferença $w = u_1 - u_2$.
	
	\textbf{Desenvolvimento Formal:} 
	Sejam $u_1$ e $u_2$ duas soluções do problema. Definimos $w = u_1 - u_2$. Pela linearidade do operador Laplaciano, temos:
	$$ \Delta w = \Delta u_1 - \Delta u_2 = f(x) - f(x) = 0 \quad \text{em } U $$
	E nas condições de contorno:
	$$ w = u_1 - u_2 = g(x) - g(x) = 0 \quad \text{em } \partial U $$
	Consideramos a integral da energia:
	$$ \int_U w \Delta w \, dx = 0 $$
	Aplicando a Primeira Identidade de Green:
	$$ \int_U w \Delta w \, dx = \int_{\partial U} w \frac{\partial w}{\partial \nu} \, d\sigma - \int_U |\nabla w|^2 \, dx $$
	Como $w = 0$ em $\partial U$, a integral de superfície anula-se, resultando em:
	$$ 0 = 0 - \int_U |\nabla w|^2 \, dx \implies \int_U |\nabla w|^2 \, dx = 0 $$
	Como $|\nabla w|^2 \ge 0$ e $w \in C^1(\overline{U})$, a integral nula implica que $\nabla w \equiv 0$ em $U$. Logo, $w$ é constante em cada componente conexa de $U$. Como $w = 0$ na fronteira $\partial U$, concluímos que $w \equiv 0$ em todo $U$, ou seja, $u_1 = u_2$.
}
	
	\section{Equação do Calor}
	\def\listid{L3} 
	\setcounter{numquestao}{0} 
	\documentclass[cs_edp_2026.1.tex]{subfiles} % Aponta para o pai
\begin{document}

% CONFIGURAÇÃO PARA COMPILAÇÃO INDIVIDUAL
\ifSubfilesClassLoaded{
	\def\listid{L3}
	\setcounter{numquestao}{54}
	\section*{Equação do Calor}
}{}

\questao{Suponha que $u = u(x,t)$ é a solução da equação do calor:
	$$ \begin{cases} u_t - \Delta u = 0, & \text{em } \mathbb{R}^n \times (0, \infty) \\ u(x,0) = g(x), & \text{para } x \in \mathbb{R}^n. \end{cases} $$
	Mostre que:
	(a) $\hat{u}_t(y,t) = -|y|^2 \hat{u}(y,t)$ e $\hat{u}(y,0) = \hat{g}(y)$;
	(b) Conclua que $\hat{u}(y,t) = e^{-|y|^2 t} \hat{g}(y)$.}

\solucao{
	\textbf{Análise do Enunciado:} 
	O problema solicita a resolução da equação do calor no espaço infinito utilizando a Transformada de Fourier no espaço espacial $x$. O objetivo é converter a EDP em uma EDO no tempo $t$ para cada frequência $y$.
	
	\textbf{Fundamentação Teórica:} 
	Utilizaremos as propriedades da Transformada de Fourier $\mathcal{F}\{u\}(y,t) = \hat{u}(y,t)$. A propriedade fundamental é que $\mathcal{F}\{\Delta u\} = -|y|^2 \hat{u}$ (Evans, Cap. 2.3).
	
	\textbf{Desenvolvimento Formal:} 
	Aplicamos a Transformada de Fourier em ambos os lados da equação $u_t - \Delta u = 0$ em relação a $x$:
	$$ \mathcal{F}\{u_t\} - \mathcal{F}\{\Delta u\} = 0 $$
	Assumindo que a derivada em $t$ pode ser trocada com a integral da transformada:
	$$ \frac{\partial}{\partial t} \hat{u}(y,t) - (-|y|^2 \hat{u}(y,t)) = 0 \implies \hat{u}_t(y,t) = -|y|^2 \hat{u}(y,t) $$
	Para a condição inicial, temos:
	$$ \hat{u}(y,0) = \mathcal{F}\{u(x,0)\} = \mathcal{F}\{g(x)\} = \hat{g}(y) $$
	Isto prova o item (a). Para o item (b), observamos que para cada $y$ fixo, temos uma EDO linear de primeira ordem em $t$:
	$$ \frac{d\hat{u}}{dt} = -|y|^2 \hat{u} $$
	A solução geral desta EDO é:
	$$ \hat{u}(y,t) = C(y) e^{-|y|^2 t} $$
	Aplicando a condição inicial $\hat{u}(y,0) = \hat{g}(y)$, obtemos $C(y) = \hat{g}(y)$. Portanto:
	$$ \hat{u}(y,t) = e^{-|y|^2 t} \hat{g}(y) $$
}

\end{document}
	
	\section{Equação da Onda}
	\def\listid{L4} 
	\setcounter{numquestao}{74} 
	\documentclass[cs_edp_2026.1.tex]{subfiles} % Aponta para o pai
\begin{document}
	
	% CONFIGURAÇÃO PARA COMPILAÇÃO INDIVIDUAL
	\ifSubfilesClassLoaded{
		\def\listid{L4}
		\setcounter{numquestao}{74}
		\section*{Equação da Onda} % Adiciona o título apenas na compilação isolada
	}{}
	
	%%%% QUESTÃO 75 %%%%
	\questao{(De volta a equação de transporte)
		\begin{enumerate}[label=(\alph*)]
			\item Suponha que se $v = v(x, t)$ satisfaz a equação de transporte homogênea $v_t + v_x = 0$ em $\mathbb{R} \times [0, \infty)$; $v(x, 0) = g(x); x \in \mathbb{R}$; então $v = g(x - t)$;
			\item Se $w = w(x, t)$ satisfaz a equação não homogênea de transporte $w_t - w_x = f(x, t)$ em $\mathbb{R} \times [0, \infty)$; $w(x, 0) = 0; x \in \mathbb{R}$; então 
			$$ w(x, t) = \int_0^t f(x + t - s, s) ds; $$
			(aqui tem o princípio de Duhamel)
			\item Se $f(x, t) = g(x - t)$ em (b), então 
			$$ w(x, t) = \int_0^t g(x + t - 2s) ds = \frac{1}{2} \int_{x-t}^{x+t} g(s) ds = \frac{1}{2} [G(x + t) - G(x - t)]; $$
			onde $G$ é uma primitiva de $g$.
			\item Considere $u$ uma solução de classe $C^2$ da equação da onda homogênea: $u_{tt} - u_{xx} = 0$ em $\mathbb{R} \times [0, \infty)$; $u(x, 0) = g(x); u_t(x, 0) = h(x); x \in \mathbb{R}$. Mostre que $v(x, t) = u_t - u_x$ satisfaz a equação de transporte do item (a), com $v(x, 0) = g'(x) + h(x)$, e conclua que $v(x, t) = g'(x - t) + h(x - t)$;
			\item Veja então que $u$ satisfaz a equação não homogênea $u_t - u_x = g'(x - t) + h(x - t)$ em $\mathbb{R} \times [0, \infty)$; $u(x, 0) = g(x)$;
			\item Conclua que 
			$$ u(x, t) = \frac{1}{2} \int_{x-t}^{x+t} h(s) ds + \frac{1}{2} [g(x + t) + g(x - t)]. $$
	\end{enumerate}}
	
	\analise{
		\textbf{Análise do Enunciado:} Esta questão foi estruturada como um roteiro de construção analítica. O objetivo não é apenas encontrar a solução final, mas derivar a fórmula de d'Alembert para a equação da onda através de etapas intermediárias. A lógica consiste em utilizar a teoria de equações de transporte (de 1\f ordem) para resolver a equação da onda (de 2\m ordem). Assim, cada item, de (a) a (e), atua como um lema necessário para a conclusão final no item (f), estabelecendo a ponte entre a propagação de ondas unidirecionais e a solução geral da equação da onda. As hipóteses centrais são a regularidade $u \in C^2$ no domínio $\mathbb{R} \times [0, \infty)$ e as condições iniciais clássicas $u(x,0) = g(x)$ e $u_t(x,0) = h(x)$.
		
		\par\textbf{Fundamentação Teórica:} Utilizamos o Método das Características (Evans, Cap. 2.2), onde a solução de $u_t + c u_x = 0$ é constante ao longo das retas $x - ct = \text{constante}$. Para a parte não homogênea, aplicamos o Princípio de Duhamel (Evans, Cap. 2.3), que trata a fonte externa como uma sucessão de impulsos que geram novas condições iniciais ao longo do tempo. A base para a redução da equação da onda é a fatoração do operador $\partial_{tt} - \partial_{xx}$ no produto $(\partial_t - \partial_x)(\partial_t + \partial_x)$.
		
		\par\textbf{Estratégias:} Começamos resolvendo a equação de transporte homogênea e a não homogênea via características. Em seguida, definimos a variável auxiliar $v = u_t - u_x$ para reduzir a ordem da equação da onda. Por fim, integramos a solução via Princípio da Superposição, somando a parte homogênea e a particular.
	}
	
	\solucao{
		\textbf{(a)} Começamos com a equação $v_t + v_x = 0$. Definimos a curva característica por $\tfrac{dx}{dt} = 1$. Integrando, temos $x(t) = t + x_0$, ou $x_0 = x - t$. Ao longo dessa curva, a variação de $v$ é:
		$$ \frac{d}{dt} v(x(t), t) = \frac{\partial v}{\partial x} \frac{dx}{dt} + \frac{\partial v}{\partial t} = v_x(1) + v_t = v_t + v_x $$
		Como a equação nos diz que $v_t + v_x = 0$, temos $\frac{d}{dt} v(x(t), t) = 0$, o que significa que $v$ é constante ao longo da característica. Assim, $v(x, t) = v(x_0, 0)$. Substituindo $x_0$ e a condição inicial $v(x,0) = g(x)$, chegamos a:
		$$ \boxed{v(x, t) = g(x - t) \quad (\ast)} $$
		
		\textbf{(b)} Agora analisamos $w_t - w_x = f(x, t)$ com $w(x,0) = 0$. As características são $\tfrac{dx}{ds} = -1$, logo $x(s) = x + t - s$. Calculando a derivada de $w$ ao longo dessa curva:
		$$ \frac{d}{ds} w(x + t - s, s) = -w_x + w_s = f(x + t - s, s) $$
		Integrando de $s=0$ até $s=t$:
		$$ \int_0^t \frac{d}{ds} w(x + t - s, s) ds = \int_0^t f(x + t - s, s) ds $$
		$$ w(x, t) - w(x+t, 0) = \int_0^t f(x + t - s, s) ds $$
		Como $w(x,0) = 0$, a solução é:
		$$ \boxed{w(x, t) = \int_0^t f(x + t - s, s) ds \quad (\ast\ast)} $$
		
		\textbf{(c)} Usamos o resultado $(\ast\ast)$ com a fonte $f(x, t) = g(x - t)$:
		$$ w(x, t) = \int_0^t g((x + t - s) - s) ds = \int_0^t g(x + t - 2s) ds $$
		Fazemos a mudança de variável $\sigma = x + t - 2s$, com $d\sigma = -2 ds$. Os limites mudam para $s=0 \to \sigma=x+t$ e $s=t \to \sigma=x-t$. A integral fica:
		$$ w(x, t) = \int_{x+t}^{x-t} g(\sigma) \left( -\tfrac{1}{2} \right) d\sigma = \tfrac{1}{2} \int_{x-t}^{x+t} g(\sigma) d\sigma = \tfrac{1}{2} [G(x + t) - G(x - t)] $$
		
		\textbf{(d)} Seja $v = u_t - u_x$. Calculamos as derivadas parciais:
		$$ v_t = u_{tt} - u_{tx} \quad \text{e} \quad v_x = u_{tx} - u_{xx} $$
		Somando as duas, temos $v_t + v_x = u_{tt} - u_{xx}$. Como $u$ resolve a equação da onda homogênea, $u_{tt} - u_{xx} = 0$, logo $v_t + v_x = 0$. Para a condição inicial:
		$$ v(x, 0) = u_t(x, 0) - u_x(x, 0) = h(x) - g'(x) $$
		\textit{Nota: O PDF indica $g'(x)+h(x)$, mas a dedução rigorosa dá $h(x)-g'(x)$. Usaremos o sinal correto para a consistência de (f).}
		
		\vspace{0.2cm}
		Observamos que a equação $v_t + v_x = 0$ com dado inicial $v(x,0) = h(x) - g'(x)$ é idêntica à estrutura resolvida no item (a). Portanto, aplicando a solução $(\ast)$, temos:
		$$ v(x, t) = h(x - t) - g'(x - t) $$
		
		\textbf{(e)} Pela definição de $v$ no item (d), temos:
		$$ u_t - u_x = v(x, t) = h(x - t) - g'(x - t) $$
		com $u(x, 0) = g(x)$. Esta é a equação de transporte não homogênea do item (b), cuja solução geral é dada por $(\ast\ast)$.
		
		\textbf{(f)} Para achar $u$, usamos o Princípio da Superposição, somando a solução homogênea $u_h$ e a particular $u_p$ ($u = u_h + u_p$).
		Para $u_h$, resolvemos $u_t - u_x = 0$ com $u(x,0) = g(x)$. Pelas características $\tfrac{dx}{dt} = -1$, temos $u_h(x, t) = g(x + t)$.
		Para $u_p$, resolvemos $u_t - u_x = f$ com $u(x,0) = 0$. Usando o resultado $(\ast\ast)$ com $f(x, s) = h(x - s) - g'(x - s)$:
		$$ u_p = \int_0^t [h(x + t - 2s) - g'(x + t - 2s)] ds $$
		Fazendo $\sigma = x + t - 2s$ ($ds = -\tfrac{1}{2} d\sigma$), temos:
		$$ u_p = \tfrac{1}{2} \int_{x-t}^{x+t} h(\sigma) d\sigma - \tfrac{1}{2} \int_{x-t}^{x+t} g'(\sigma) d\sigma = \tfrac{1}{2} \int_{x-t}^{x+t} h(\sigma) d\sigma - \tfrac{1}{2} [g(x + t) - g(x - t)] $$
		Somando $u_h$ e $u_p$:
		$$ u(x, t) = g(x + t) + \tfrac{1}{2} \int_{x-t}^{x+t} h(s) ds - \tfrac{1}{2} g(x + t) + \tfrac{1}{2} g(x - t) $$
		Simplificando, chegamos à fórmula de d'Alembert:
		$$ \boxed{u(x, t) = \tfrac{1}{2} [g(x + t) + g(x - t)] + \tfrac{1}{2} \int_{x-t}^{x+t} h(s) ds \quad (\ast\ast\ast)} $$
	}
	
	\sintese{
		A resolução desta questão demonstra que a equação da onda é a superposição de dois fluxos de transporte em direções opostas. A \textbf{Construção Analítica} do raciocínio pode ser resumida no seguinte fluxo:
		
		\begin{center}
			$\text{Item (a)} \longrightarrow \text{Item (d)} \longrightarrow \text{Item (e)} \longrightarrow \text{Item (f)}$ \\
			$\text{Item (b)} \longrightarrow \text{Item (f)}$
		\end{center}
		
		O resultado do Item (a) fornece a técnica de solução para a direita, usada no Item (d) para descobrir a variável auxiliar $v$. O resultado do Item (b) fornece a solução para o transporte não homogêneo, usada no Item (f) para reconstruir a solução completa de d'Alembert.
	}
	
\end{document}
	
	\section{Espaços de Sobolev}
	\def\listid{L5} 
	\setcounter{numquestao}{0} 
	\documentclass[cs_edp_2026.1.tex]{subfiles} % Aponta para o pai
\begin{document}
	
% CONFIGURAÇÃO PARA COMPILAÇÃO INDIVIDUAL
\ifSubfilesClassLoaded{
	\def\listid{L5}
	\setcounter{numquestao}{84}
	\section*{Espaços de Sobolev}
}{}

%%%% QUESTÃO 85 %%%%
\questao{Seja $f : \mathbb{R} \to \mathbb{R}$ em $f \in C^{0,\alpha}(\mathbb{R})$, isto quer dizer que $f$ é limitado ($|f(t)| \le \|f\|_\infty$, para todo $t \in \mathbb{R}$) e
	$$[f]_\alpha = \sup_{t \neq s} \frac{|f(t) - f(s)|}{|s - t|^\alpha} < \infty.$$
	Suponha que $U$ é um aberto e $u \in C^{0,\beta}(\overline{U})$. Mostre que $f \circ u$ está em $C^{0,\alpha\beta}(\overline{U})$ e
	$$\|f \circ u\|_{C^{0,\alpha\beta}(\overline{U})} \le \|f\|_\infty + [f]_\alpha [u]_{\beta, U}^\alpha.$$}

\analise{
	\textbf{Análise do Enunciado:}
	A questão solicita a prova de que a composição de duas funções com regularidade de \textbf{Hölder} resulta em uma função cuja regularidade é dada pelo produto dos expoentes $\alpha$ e $\beta$. A tese central é a obtenção de uma estimativa a priori para a norma do espaço $C^{0, \alpha\beta}(\overline{U})$.
	
	\textbf{Fundamentação Teórica:}
	A fundamentação baseia-se na definição de \textbf{Espaços de Hölder} e suas propriedades algébricas. Consultamos Gilbarg \& Trudinger (Capítulo 4), onde a norma de Hölder é definida como a soma da norma do supremo ($\|\cdot\|_\infty$) e a seminorma de Hölder ($[\cdot]_\gamma$).
	
	\textbf{Estratégias:}
	\begin{enumerate}
		\item \textbf{Estimativa da Norma do Supremo:} Demonstrar que a composição $f \circ u$ permanece limitada em $\overline{U}$, utilizando a hipótese de que $f$ é globalmente limitada em $\mathbb{R}$.
		\item \textbf{Análise da Variação Incremental:} Aplicar a definição de continuidade de \textbf{Hölder} para $f$ sobre a imagem de $u$, e subsequentemente a continuidade de \textbf{Hölder} de $u$ sobre o domínio $\overline{U}$.
		\item \textbf{Dedução do Expoente:} Mostrar que a combinação das desigualdades resulta no expoente $\alpha\beta$ para a distância $|x-y|$.
		\item \textbf{Fechamento da Norma:} Somar a estimativa da norma do supremo com a seminorma obtida para validar a desigualdade final.
	\end{enumerate}
}

\solucao{
	Iniciamos a análise controlando a norma do supremo da função composta $v = f \circ u$. Dado que $u(x) \in \mathbb{R}$ para todo $x \in \overline{U}$ e que a função $f$ é globalmente limitada em $\mathbb{R}$ por $\|f\|_\infty$, temos:
	$$|v(x)| = |f(u(x))| \le \|f\|_\infty, \quad \forall x \in \overline{U}.$$
	Dessa forma, estabelecemos a primeira estimativa a priori:
	$$\|f \circ u\|_\infty \le \|f\|_\infty \quad (\ast_1)$$
	
	Para a estimativa da seminorma de \textbf{Hölder}, fixamos $x, y \in \overline{U}$ com $x \neq y$. Como, por hipótese $f \in C^{0,\alpha}(\mathbb{R})$, então para variação de $f$, obtemos que:
	$$|f(u(x)) - f(u(y))| \le [f]_\alpha |u(x) - u(y)|^\alpha \quad (\ast_2)$$
	
	Além disso, a regularidade de $u \in C^{0,\beta}(\overline{U})$ implica para a variação de $u$, que:
	$$|u(x) - u(y)| \le [u]_{\beta, U} |x - y|^\beta \quad (\ast_3)$$
	
	Substituindo a desigualdade $(\ast_3)$ em $(\ast_2)$, obtemos:
	\begin{align*}
		|f(u(x)) - f(u(y))| &\le [f]_\alpha \left( [u]_{\beta, U} |x - y|^\beta \right)^\alpha \le [f]_\alpha [u]_{\beta, U}^\alpha |x - y|^{\alpha\beta}.
	\end{align*}
	
	Dividindo ambos os lados por $|x - y|^{\alpha\beta}$ e tomando o supremo sobre $x, y \in \overline{U}$ com $x \neq y$, isolamos a seminorma de \textbf{Hölder} de ordem $\alpha\beta$:
	$$\boxed{[f \circ u]_{\alpha\beta} \le [f]_\alpha [u]_{\beta, U}^\alpha \quad (\ast_4)}$$
	
	Finalmente, utilizando a definição da norma total no espaço $C^{0, \alpha\beta}(\overline{U})$, que é a soma da norma do supremo com a seminorma de \textbf{Hölder}, temos:
	$$\|f \circ u\|_{C^{0,\alpha\beta}(\overline{U})} = \|f \circ u\|_\infty + [f \circ u]_{\alpha\beta}.$$
	
	Aplicando as estimativas $(\ast_1)$ e $(\ast_4)$ nesta soma, chegamos ao resultado final:
	$$\boxed{\|f \circ u\|_{C^{0,\alpha\beta}(\overline{U})} \le \|f\|_\infty + [f]_\alpha [u]_{\beta, U}^\alpha.}$$
}

\sintese{
	A prova demonstra que a regularidade de \textbf{Hölder} é preservada sob composição, embora ocorra uma \textbf{perda de regularidade} no sentido de que o expoente resultante é o produto dos expoentes originais ($\alpha\beta$), sendo este estritamente menor que qualquer um dos expoentes iniciais para $\alpha, \beta \in (0,1)$. Logicamente, o fluxo foi:
	$$\text{Limitada} \to \text{Variação de } f \text{ em } u \to \text{Variação de } u \text{ em } x \to \text{Soma de Normas}.$$
}

\aprofundamento{
	Esta propriedade é crucial para a teoria de \textbf{Regularidade Elíptica}. Quando lidamos com equações do tipo $-\Delta u = f(u)$, a regularidade de $u$ (frequentemente em $C^{1,\beta}$ ou $W^{2,p}$) implica que $f(u)$ herda a regularidade de $f$ composta com $u$. 
	
	Se $f$ fosse apenas localmente de \textbf{Hölder}, a limitação global $\|f\|_\infty$ não seria dada, mas a prova permaneceria válida desde que $u$ fosse limitada (o que é garantido por $\overline{U}$ ser compacto). Historicamente, a transição de $C^k$ para $C^{k, \alpha}$ permitiu a Schauder resolver a equação de Poisson com dados contínuos, superando a limitação de que $u \in C^2$ não é garantido apenas por $f \in C^0$.
}

%%%% QUESTÃO 86 %%%%
\questao{Mostre que a derivada fraca $\frac{\partial u}{\partial x_j}$ de uma função $u \in L_{loc}^1(U)$ quando existe é única.}

\analise{
	\textbf{Análise do Enunciado:}
	A questão exige a demonstração da unicidade do elemento $v \in L_{loc}^1(U)$ que satisfaz a definição de derivada fraca de uma função $u$ em relação à variável $x_j$. A tese é que, se dois elementos satisfazem a mesma condição integral para todas as funções de teste, eles devem coincidir quase sempre em $U$.
	
	\textbf{Fundamentação Teórica:}
	A base teórica reside na definição de \textbf{Derivada Fraca} e no \textbf{Lema Fundamental do Cálculo de Variações}. Consultamos Brezis (Capítulo 9), onde a derivada fraca é introduzida via integração por partes formal, e Evans (Capítulo 5), que discute a densidade de $C_c^\infty(U)$ em espaços $L^p$.
	
	\textbf{Estratégias:}
	\begin{enumerate}
		\item \textbf{Definição de Hipóteses:} Supor a existência de duas funções $v_1, v_2 \in L_{loc}^1(U)$ que atuam como derivadas fracas de $u$ em relação a $x_j$.
		\item \textbf{Construção da Diferença:} Subtrair as identidades integrais correspondentes a $v_1$ e $v_2$ para isolar a diferença $(v_1 - v_2)$.
		\item \textbf{Aplicação do Lema Fundamental:} Utilizar o fato de que se a integral de uma função $L_{loc}^1$ contra qualquer função de teste $\phi \in C_c^\infty(U)$ é nula, então a função é nula quase sempre.
		\item \textbf{Conclusão:} Estabelecer a igualdade $v_1 = v_2$ q.s. em $U$.
	\end{enumerate}
}

\solucao{
	Seja $U \subset \mathbb{R}^n$ um aberto e $u \in L_{loc}^1(U)$. Por definição, dizemos que $v \in L_{loc}^1(U)$ é a derivada fraca de $u$ em relação a $x_j$, denotada por $\tfrac{\partial u}{\partial x_j}$, se para toda função de teste $\phi \in C_c^\infty(U)$ for satisfeita a identidade:
	$$\int_U u \frac{\partial \phi}{\partial x_j} dx = - \int_U v \phi dx \quad (\ast_1)$$
	
	Suponhamos que existam duas funções $v_1, v_2 \in L_{loc}^1(U)$ que satisfaçam a condição de derivada fraca de $u$. Assim, para qualquer $\phi \in C_c^\infty(U)$, temos simultaneamente:
	\begin{align*}
		\int_U u \frac{\partial \phi}{\partial x_j} dx &= - \int_U v_1 \phi dx \quad (\ast_2) \ \ \text{e } \ \int_U u \frac{\partial \phi}{\partial x_j} dx &= - \int_U v_2 \phi dx \quad (\ast_3)
	\end{align*}
	
	Subtraindo a equação $(\ast_3)$ da equação $(\ast_2)$, obtemos:
	\begin{align*}
		0 &= - \int_U v_1 \phi dx - \left( - \int_U v_2 \phi dx \right) = \int_U (v_2 - v_1) \phi dx
	\end{align*}
	
	Portanto, a diferença $w = v_2 - v_1$ satisfaz a seguinte condição para todo $\phi \in C_c^\infty(U)$:
	$$\boxed{\int_U w \phi dx = 0}$$
	
	Sendo $w \in L_{loc}^1(U)$, aplicamos o \textbf{Lema Fundamental do Cálculo de Variações} (ou a propriedade de que $C_c^\infty(U)$ é denso em $L^p(K)$ para qualquer compacto $K \subset U$). Este lema estabelece que, se uma função localmente integrável anula a integral contra toda função de teste com suporte compacto, então essa função deve ser nula quase sempre em $U$. 
	
	Logo:
	$$w(x) = 0 \quad \text{q.s. em } U \implies v_2(x) = v_1(x) \quad \text{q.s. em } U.$$
	
	Concluímos, portanto, que a derivada fraca, quando existe, é única a menos de um conjunto de medida nula.
}

\sintese{
	A unicidade da derivada fraca é uma consequência direta da densidade das funções de teste $C_c^\infty(U)$ no espaço de funções integráveis. O fluxo lógico foi:
	$$\text{Hipótese de Duplicidade} \to \text{Linearidade da Integral} \to \text{Ortogonalidade a } C_c^\infty \to \text{Nulidade q.s.}$$
}

\aprofundamento{
	A unicidade da derivada fraca é o que permite a definição consistente de operadores diferenciais em espaços de funções menos regulares. Do ponto de vista da \textbf{Teoria das Distribuições}, a derivada fraca de $u$ é precisamente a distribuição $\partial_j u$. Como a aplicação de uma distribuição a uma função de teste é única, a representação dessa distribuição por uma função $v \in L_{loc}^1$ (quando possível) deve ser única. 
	
	Se $u$ possuísse derivadas fracas distintas, a noção de $\nabla u$ \textbf{fraco} estaria comprometida, tornando impossível a definição de normas em espaços de Sobolev como $W^{1,p}(U)$, onde a norma depende explicitamente da unicidade de $\tfrac{\partial u}{\partial x_j}$.
}

%%%% QUESTÃO 87 %%%%
\questao{Seja $u_m \in W^{1,p}(U)$ uma sequência tal que $u_m \to u$ em $L^p(U)$ e $\tfrac{\partial u_m}{\partial x_j} \to g$ em $L^p(U)$. Mostre que existe a derivada fraca $\tfrac{\partial u}{\partial x_j}$ e que esta derivada fraca é igual a $g$.}

\analise{
	\textbf{Análise do Enunciado:}
	A questão solicita a prova de que o limite de uma sequência de funções em $W^{1,p}(U)$, onde tanto as funções quanto suas derivadas convergem em $L^p(U)$, preserva a relação de derivação fraca no limite. A tese é a identificação de $g$ como a derivada fraca de $u$ em relação a $x_j$.
	
	\textbf{Fundamentação Teórica:}
	A prova fundamenta-se na definição de \textbf{Derivada Fraca} e nas propriedades de convergência em espaços $L^p$. A ferramenta principal é a \textbf{Desigualdade de Hölder}, que garante que a convergência em norma $L^p$ implica a convergência das integrais contra funções de teste em $L^{p'}$, onde $\tfrac{1}{p} + \tfrac{1}{p'} = 1$. Consultamos Brezis (Capítulo 9) para a definição de convergência em espaços de Sobolev.
	
	\textbf{Estratégias:}
	\begin{enumerate}
		\item \textbf{Identidade de Integração:} Partir da definição de derivada fraca para cada elemento $u_m$ da sequência.
		\item \textbf{Passagem ao Limite no Lado Esquerdo:} Demonstrar que $\int_U u_m \frac{\partial \phi}{\partial x_j} dx$ converge para $\int_U u \frac{\partial \phi}{\partial x_j} dx$ utilizando a convergência $u_m \to u$ em $L^p$.
		\item \textbf{Passagem ao Limite no Lado Direito:} Demonstrar que $\int_U \frac{\partial u_m}{\partial x_j} \phi dx$ converge para $\int_U g \phi dx$ utilizando a convergência $\frac{\partial u_m}{\partial x_j} \to g$ em $L^p$.
		\item \textbf{Identificação Final:} Mostrar que a identidade resultante no limite coincide exatamente com a definição de derivada fraca para $u$, identificando $g$ como tal.
	\end{enumerate}
}

\solucao{
	Como $u_m \in W^{1,p}(U)$, cada $u_m$ possui uma derivada fraca $\frac{\partial u_m}{\partial x_j} \in L^p(U)$. Portanto, por definição, para qualquer função de teste $\phi \in C_c^\infty(U)$, é válida a identidade:
	$$\int_U u_m \frac{\partial \phi}{\partial x_j} dx = - \int_U \frac{\partial u_m}{\partial x_j} \phi dx \quad (\ast_1)$$
	
	Analisamos agora o comportamento do lado esquerdo da igualdade $(\ast_1)$ quando $m \to \infty$. Consideramos a diferença:
	$$\left| \int_U u_m \frac{\partial \phi}{\partial x_j} dx - \int_U u \frac{\partial \phi}{\partial x_j} dx \right| = \left| \int_U (u_m - u) \frac{\partial \phi}{\partial x_j} dx \right|.$$
	
	Aplicando a \textbf{Desigualdade de Hölder}, onde $p$ é o expoente da sequência e $p'$ é o expoente conjugado $\left(\tfrac{1}{p} + \tfrac{1}{p'} = 1\right)$, temos:
	$$\left| \int_U (u_m - u) \frac{\partial \phi}{\partial x_j} dx \right| \le \|u_m - u\|_{L^p(U)} \left\| \frac{\partial \phi}{\partial x_j} \right\|_{L^{p'}(U)}.$$
		
	Visto que $\phi \in C_c^\infty(U)$, a derivada $\tfrac{\partial \phi}{\partial x_j}$ é contínua com suporte compacto, logo pertence a $L^{p'}(U)$. Como $u_m \to u$ em $L^p(U)$, a norma $\|u_m - u\|_{L^p(U)} \to 0$. Portanto:
	$$\lim_{m \to \infty} \int_U u_m \frac{\partial \phi}{\partial x_j} dx = \int_U u \frac{\partial \phi}{\partial x_j} dx \quad (\ast_2)$$
	
	Analogamente, analisamos o lado direito de $(\ast_1)$. Denotamos $v_m = \tfrac{\partial u_m}{\partial x_j}$ e sabemos que $v_m \to g$ em $L^p(U)$. Novamente, via \textbf{Desigualdade de Hölder}:
	$$\left| \int_U v_m \phi dx - \int_U g \phi dx \right| = \left| \int_U (v_m - g) \phi dx \right| \le \|v_m - g\|_{L^p(U)} \|\phi\|_{L^{p'}(U)}.$$
	
	Como $v_m \to g$ em $L^p(U)$, a norma $\|v_m - g\|_{L^p(U)} \to 0$. Assim:
	$$\lim_{m \to \infty} \left( - \int_U \frac{\partial u_m}{\partial x_j} \phi dx \right) = - \int_U g \phi dx \quad (\ast_3)$$
	
	Tendo convergência em ambos os membros da identidade $(\ast_1)$, podemos passar ao limite no sinal de igualdade:
	$$\lim_{m \to \infty} \int_U u_m \frac{\partial \phi}{\partial x_j} dx = \lim_{m \to \infty} \left( - \int_U \frac{\partial u_m}{\partial x_j} \phi dx \right).$$
	
	Substituindo os resultados $(\ast_2)$ e $(\ast_3)$, obtemos a relação:
	$$\boxed{\int_U u \frac{\partial \phi}{\partial x_j} dx = - \int_U g \phi dx \quad (\ast_4)}$$
	
	A identidade $(\ast_4)$ é precisamente a definição de derivada fraca para a função $u$ em relação à variável $x_j$, com $g$ atuando como a derivada. Como $g \in L^p(U) \subset L_{loc}^1(U)$, a condição de integrabilidade local está satisfeita.
	
	Concluímos que a derivada fraca $\tfrac{\partial u}{\partial x_j}$ existe e que:
	$$\frac{\partial u}{\partial x_j} = g \quad \text{q.s. em } U.$$
}
	
\sintese{
	A prova demonstra a estabilidade do operador de derivação fraca sob a topologia de $L^p$. O raciocínio foi fundamentado pela linearidade da integral e pelo controle de erro via norma de Hölder, seguindo o fluxo:
	$$\text{Identidade de } u_m \to \text{Convergência de } L^p \to \text{Limite de Integrais} \to \text{Definição de } \partial_j u.$$
}

\aprofundamento{
	Este resultado prova que o espaço de Sobolev $W^{1,p}(U)$ é um \textbf{Espaço de Banach}. Se $u_m$ for uma sequência de Cauchy em $W^{1,p}$, então $u_m \to u$ em $L^p$ e $\nabla u_m \to g$ em $L^p$. A Questão 87 garante que $u$ possui derivadas fracas e que $\nabla u = g$, logo $u \in W^{1,p}(U)$, provando que o espaço é completo.
	
	Sob a ótica da \textbf{Teoria das Distribuições}, isso equivale a dizer que a operação de derivação $\partial_j : \mathcal{D}'(U) \to \mathcal{D}'(U)$ é contínua em relação à topologia de distribuições. Quando a distribuição $\partial_j u$ pode ser representada por uma função $L^p$, dizemos que a função possui regularidade de Sobolev.
}

%%%% QUESTÃO 88 %%%%
\questao{Seja $U$ conexo e $u \in W^{1,p}(U)$ tal que $$\int_U u \frac{\partial \phi}{\partial x_j} dx = 0 \ \ \text{para todos } j = 1, 2, \dots, n \ \ \text{e } \phi \in C_c^\infty(U).$$ Mostre que $u$ é constante, q.s. em $U$. Se $|U| = \infty$ e $p < \infty$ devemos ter $u$ constante igual a $0$.}

\analise{
	\textbf{Análise do Enunciado:}
	A hipótese dada é que a derivada fraca de $u$ em relação a cada variável $x_j$ é nula, ou seja, o \textbf{Gradiente Fraco} $\nabla u = 0$ quase sempre em $U$. A tese divide-se em duas partes: primeiro, provar que $u$ é constante quase sempre, dado que $U$ é conexo; segundo, provar que, sob a condição de medida infinita do domínio e pertinência a $L^p$ ($p < \infty$), a única constante possível é zero.
	
	\textbf{Fundamentação Teórica:}
	A prova fundamenta-se na técnica de \textbf{Regularização por Convolução}, cujas propriedades de convergência e preservação de derivadas foram construídas e validadas nas \textbf{Questões 4 e 5}. O argumento central utiliza o fato de que funções suaves com gradiente nulo em domínios conexos são constantes. Referências principais: Brezis (Capítulo 9) e Evans (Capítulo 5.3).
	
	\textbf{Estratégias:}
	\begin{enumerate}
		\item \textbf{Regularização:} Definir a sequência de funções suavizadas $u_\varepsilon = u * \eta_\varepsilon$ e demonstrar que suas derivadas clássicas são nulas.
		\item \textbf{Uso da Conexão:} Aplicar o resultado do cálculo clássico para concluir que $u_\varepsilon$ é constante em cada componente conexo.
		\item \textbf{Passagem ao Limite:} Utilizar a convergência $u_\varepsilon \to u$ em $L^p_{loc}(U)$ para transferir a propriedade de ser constante para $u$.
		\item \textbf{Análise de Medida:} Para o caso $|U| = \infty$, utilizar a norma de $L^p$ para forçar a constante a ser zero.
	\end{enumerate}
}

\solucao{
	A hipótese $\int_U u \tfrac{\partial \phi}{\partial x_j} dx = 0$ para todo $\phi \in C_c^\infty(U)$ implica, por definição, que a derivada fraca $\frac{\partial u}{\partial x_j} = 0$ quase sempre em $U$ para todo $j=1, \dots, n$. Portanto, $\nabla u = 0$ q.s. em $U$.
	
	Consideremos a função de regularização $\eta \in C_c^\infty(\mathbb{R}^n)$ com $\text{supp } \eta \subset B_1(0)$ e $\int \eta = 1$. Para $\varepsilon > 0$, definimos a função suave $u_\varepsilon = u * \eta_\varepsilon$ em $U_\varepsilon = \{x \in U : \dist(x, \partial U) > \varepsilon\}$. Pela propriedade da convolução, $u_\varepsilon \in C^\infty(U_\varepsilon)$ e a derivada clássica de $u_\varepsilon$ coincide com a convolução da derivada fraca de $u$ com $\eta_\varepsilon$:
	$$\frac{\partial u_\varepsilon}{\partial x_j}(x) = \int_U \frac{\partial u}{\partial x_j}(y) \eta_\varepsilon(x-y) dy$$
	
	Como $\frac{\partial u}{\partial x_j} = 0$ q.s. em $U$, temos que $\frac{\partial u_\varepsilon}{\partial x_j} = 0$ para todo $x \in U_\varepsilon$ e todo $j=1, \dots, n$. Assim:
	$$\boxed{\nabla u_\varepsilon = 0 \text{ em } U_\varepsilon \quad (\ast_1)}$$
	
	Visto que $U$ é conexo, $U_\varepsilon$ também é conexo para $\varepsilon$ suficientemente pequeno. Portanto:
	$$u_\varepsilon(x) = C_\varepsilon \quad \text{para todo } x \in U_\varepsilon \quad (\ast_2)$$
	
	Além disso, $u_\varepsilon \to u$ em $L^p_{loc}(U)$ quando $\varepsilon \to 0$. Assim, para qualquer $V \subset \subset U$ limitado, temos:
	$$\|u - C_\varepsilon\|_{L^p(V)} \to 0$$
	
	Isso implica que $u$ é igual a uma constante $C$ quase sempre em $U$, ou seja, 
	$$\boxed{u(x) = C \quad \text{q.s. em } U}.$$
	
	Agora, consideremos a hipótese $|U| = \infty$ com $p < \infty$. Dado que $u \in W^{1,p}(U)$, obrigatoriamente $u \in L^p(U)$. Calculamos a norma de $u$:
	$$\|u\|_{L^p(U)}^p = \int_U |u(x)|^p dx = \int_U |C|^p dx = |C|^p \int_U 1 dx = |C|^p |U|$$
	
	Se $C \neq 0$, então $\|u\|_{L^p(U)}^p = |C|^p \cdot \infty = \infty$, o que contradiz a hipótese de que $u \in L^p(U)$. Para que a integral seja finita com $|U| = \infty$, a única possibilidade é:
	$$\boxed{C = 0 }.$$
}

\sintese{
	A prova articula a transição entre a noção de derivada fraca e a derivada clássica via regularização. A conexão do domínio é a peça chave que permite a propagação da nulidade do gradiente para a constância da função. O fluxo lógico foi:
	$$\nabla u = 0 \text{ (fraco)} \to \nabla u_\varepsilon = 0 \text{ (clássico)} \to u_\varepsilon = C_\varepsilon \to u = C \text{ q.s.} \to L^p \text{ com } |U|=\infty \implies C=0.$$
}

\aprofundamento{
	Este resultado é a base para a definição de espaços de Sobolev com valores médios nulos e para a prova da \textbf{Desigualdade de Poincaré}. A exigência de que $U$ seja conexo é indispensável; se $U$ fosse a união de dois abertos disjuntos $U_1 \cup U_2$, $u$ poderia assumir constantes distintas $C_1$ e $C_2$ em cada componente, mantendo $\nabla u = 0$.
	
	Do ponto de vista da análise funcional, este teorema afirma que o núcleo ($\text{ker}$) do operador gradiente $$\nabla : W^{1,p}(U) \to L^p(U; \mathbb{R}^n),$$ consiste precisamente das funções constantes. Quando $|U| = \infty$ e $p < \infty$, a única constante que pertence ao espaço $L^p$ é a função nula, o que torna o operador $\nabla$ injetivo neste caso específico.
}

%%%% QUESTÃO 89 %%%%
\questao{
	\begin{enumerate}[label*=(\alph*)]
		\item Sejam $u \in W^{1,p}(U), v \in W^{1,q}(U)$. Mostre que $uv \in W^{1,1}(U)$ e que $\nabla(uv) = u\nabla v + v\nabla u$;
		\item Se $v \in W^{1,\infty}(U)$ então $uv \in W^{1,p}(U)$.
	\end{enumerate}
}	

\analise{
	\textbf{Análise do Enunciado:}
	O objetivo é estender a regra do produto do cálculo clássico para o contexto de derivadas fracas. No item (a), deve-se mostrar que o produto de duas funções de Sobolev pertence a $W^{1,1}(U)$, assumindo-se que os expoentes $p, q$ satisfaçam a condição de integrabilidade $\tfrac{1}{p} + \tfrac{1}{q} \le 1$. No item (b), a regularidade de $v$ é elevada para $L^\infty$, o que permite que o produto $uv$ permaneça no espaço $W^{1,p}(U)$.
	
	\textbf{Fundamentação Teórica:}
	A prova baseia-se na técnica de \textbf{Regularização por Convolução}** (estudada nas Questões 4 e 5). A estratégia consiste em aplicar a regra do produto a funções suaves (onde a validade é trivial) e passar ao limite utilizando a \textbf{Desigualdade de Hölder}.
	
	\textbf{Estratégias:}
	\begin{enumerate}
		\item \textbf{Aproximação:} Definir sequências de regularização $u_\varepsilon = u * \eta_\varepsilon$ e $v_\varepsilon = v * \eta_\varepsilon$.
		\item \textbf{Identidade Clássica:} Utilizar a regra do produto para as funções suaves $u_\varepsilon$ e $v_\varepsilon$.
		\item \textbf{Passagem ao Limite:} Provar que os termos $u_\varepsilon \nabla v_\varepsilon$ e $v_\varepsilon \nabla u_\varepsilon$ convergem em $L^1(U)$ para $u \nabla v$ e $v \nabla u$, respectivamente.
		\item \textbf{Fechamento em $W^{1,p}$:} Para o item (b), utilizar as normas de $L^\infty$ para mostrar que a resultante pertence a $L^p(U)$.
	\end{enumerate}
}

\solucao{
	\textbf{(a):}
	Sejam $u_\varepsilon$ e $v_\varepsilon$ as regularizações de $u$ e $v$ via convolução com $\eta_\varepsilon$. Como $u_\varepsilon, v_\varepsilon \in C^\infty(U_\varepsilon)$, a regra do produto clássica é válida:
	$$\nabla(u_\varepsilon v_\varepsilon) = u_\varepsilon \nabla v_\varepsilon + v_\varepsilon \nabla u_\varepsilon \quad (\ast_1)$$
	
	Para qualquer função de teste $\phi \in C_c^\infty(U)$, temos a identidade:
	$$\int_U (u_\varepsilon v_\varepsilon) \nabla \phi \, dx = - \int_U (u_\varepsilon \nabla v_\varepsilon + v_\varepsilon \nabla u_\varepsilon) \phi \, dx \quad (\ast_2)$$
	
	Analisamos a convergência de cada termo quando $\varepsilon \to 0$. Pela \textbf{Desigualdade de Hölder}, e assumindo $\tfrac{1}{p} + \tfrac{1}{q} \le 1$:
	$$\int_U |u_\varepsilon \nabla v_\varepsilon - u \nabla v| \, dx \le \|u_\varepsilon - u\|_{L^p} \|\nabla v_\varepsilon\|_{L^q} + \|u\|_{L^p} \|\nabla v_\varepsilon - \nabla v\|_{L^q}$$
	
	Visto que $u_\varepsilon \to u$ em $L^p$ e $\nabla v_\varepsilon \to \nabla v$ em $L^q$, o lado direito converge a zero. Analogamente, $v_\varepsilon \nabla u_\varepsilon \to v \nabla u$ em $L^1(U)$. Passando ao limite na relação $(\ast_2)$, obtemos:
	$$\int_U (uv) \nabla \phi \, dx = - \int_U (u \nabla v + v \nabla u) \phi \, dx \quad (\ast_3)$$
	
	A identidade $(\ast_3)$ prova que $uv$ possui derivada fraca e que:
	$$\boxed{\nabla(uv) = u \nabla v + v \nabla u \quad (\ast_4)}$$
	Como $u \nabla v \in L^1$ e $v \nabla u \in L^1$ (via Hölder), temos $uv \in W^{1,1}(U)$.
	
	\vspace{10pt}
	\textbf{Item (b):}
	Agora, suponha $v \in W^{1,\infty}(U)$. Isso implica que $v \in L^\infty(U)$ e $\nabla v \in L^\infty(U; \mathbb{R}^n)$.
	Analisamos a norma $L^p$ do produto $uv$:
	$$\|uv\|_{L^p(U)} \le \|v\|_{L^\infty(U)} \|u\|_{L^p(U)} < \infty \quad (\ast_5)$$
	
	Para o gradiente $\nabla(uv)$, utilizamos a regra do produto obtida em $(\ast_4)$:
	$$\|\nabla(uv)\|_{L^p(U)} = \|u \nabla v + v \nabla u\|_{L^p(U)} \le \|u \nabla v\|_{L^p(U)} + \|v \nabla u\|_{L^p(U)}$$
	
	Aplicando a desigualdade $\|fg\|_{L^p} \le \|f\|_{L^p} \|g\|_{L^\infty}$:
	$$\|\nabla(uv)\|_{L^p(U)} \le \|u\|_{L^p} \|\nabla v\|_{L^\infty} + \|v\|_{L^\infty} \|\nabla u\|_{L^p} < \infty \quad (\ast_6)$$
	
	As estimativas $(\ast_5)$ e $(\ast_6)$ garantem que $uv \in L^p(U)$ e $\nabla(uv) \in L^p(U)$, logo:
	$$\boxed{uv \in W^{1,p}(U)}.$$
}

\sintese{
	A prova estende a linearidade da derivada para produtos de funções de Sobolev. O ponto crítico é a utilização da regularização para justificar a regra do produto e a aplicação da desigualdade de Hölder para garantir que os termos resultantes permaneçam integráveis. O fluxo lógico foi:
	$$\text{Regularização} \to \text{Regra Clássica} \to \text{Limite } L^p \to \text{Regra Fraca} \to \text{Estabilidade em } W^{1,p}.$$
}

\aprofundamento{
	Este resultado demonstra que $W^{1, \infty}(U)$ atua como um \textbf{multiplicador} para o espaço $W^{1,p}(U)$. De forma mais geral, se $u \in W^{k,p}(U)$ e $v \in W^{k, \infty}(U)$, então $uv \in W^{k,p}(U)$. 
	
	É importante notar que, se $p < \infty$ e $q < \infty$ com $\tfrac{1}{p} + \tfrac{1}{q} > 1$, o produto $uv$ pode não pertencer a $L^1(U)$, e a derivada fraca não estaria bem definida no sentido de $L^1$. Por isso, a condição de $v \in W^{1, \infty}$ no item (b) é fundamental para \textbf{preservar} o expoente $p$ da função $u$. Esse comportamento é análogo ao fato de que $L^\infty$ é a única álgebra entre os espaços $L^p$ para $p < \infty$.
}

%%%% QUESTÃO 90 %%%%
\questao{Seja $u \in W^{2,p}(U)$. Mostre que $\frac{\partial^2 u}{\partial x_i \partial x_j} = \frac{\partial^2 u}{\partial x_j \partial x_i}$.}

\analise{
	\textbf{Análise do Enunciado:}
	A questão solicita a demonstração de que as derivadas fracas de segunda ordem comutam. A hipótese $u \in W^{2,p}(U)$ garante que $u$ e todas as suas derivadas fracas de até segunda ordem pertencem a $L^p(U)$. A tese é a igualdade quase sempre (q.s.) das funções que representam as derivadas mistas $\partial_{x_i}\partial_{x_j}u$ e $\partial_{x_j}\partial_{x_i}u$.
	
	\textbf{Fundamentação Teórica:}
	A prova baseia-se na definição de \textbf{Derivada Fraca} aplicada sucessivamente. O ponto central é a utilização de funções de teste $\phi \in C_c^\infty(U)$, para as quais a comutatividade das derivadas clássicas é garantida pelo Teorema de Clairaut. Consultamos Brezis (Capítulo 9) para a definição de $W^{2,p}$.
	
	\textbf{Estratégias:}
	\begin{enumerate}
		\item \textbf{Aplicação da Definição:} Utilizar a definição de derivada fraca para a segunda derivada $\frac{\partial^2 u}{\partial x_i \partial x_j}$ contra uma função de teste $\phi$.
		\item \textbf{Comutatividade Clássica:} Inverter a ordem de derivação na função de teste $\phi$, que é de classe $C^\infty$.
		\item \textbf{Reversão da Definição:} Reconhecer que a integral resultante corresponde à definição de derivada fraca para a ordem $\frac{\partial^2 u}{\partial x_j \partial x_i}$.
		\item \textbf{Unicidade:} Aplicar a unicidade da derivada fraca (provada na Questão 86) para concluir a igualdade q.s.
	\end{enumerate}
}

\solucao{
	Seja $u \in W^{2,p}(U)$. Por definição, a derivada fraca $\frac{\partial^2 u}{\partial x_i \partial x_j}$ é a única função $v_{ij} \in L^p(U)$ tal que, para toda função de teste $\phi \in C_c^\infty(U)$, temos:
	$$\int_U v_{ij} \phi \, dx = \int_U u \frac{\partial^2 \phi}{\partial x_i \partial x_j} dx \quad (\ast_1)$$
	
	Analogamente, a derivada fraca $\frac{\partial^2 u}{\partial x_j \partial x_i}$ é a única função $v_{ji} \in L^p(U)$ tal que, para toda $\phi \in C_c^\infty(U)$:
	$$\int_U v_{ji} \phi \, dx = \int_U u \frac{\partial^2 \phi}{\partial x_j \partial x_i} dx \quad (\ast_2).$$
	
	Como $\phi \in C_c^\infty(U)$, de $(\ast_1)$ $(\ast_2)$, obtemos:
	\begin{align*}
		\int_U (v_{ij} - v_{ji}) \phi \, dx &= \int_U v_{ij} \phi \, dx-\int_U v_{ji} \phi \, dx = \int_U u \frac{\partial^2 \phi}{\partial x_i \partial x_j}  dx - \int_U u \frac{\partial^2 \phi}{\partial x_j \partial x_i} dx \\
		&= \int_U u \left(\frac{\partial^2 \phi}{\partial x_i \partial x_j}- \frac{\partial^2 \phi}{\partial x_j \partial x_i}\right) dx
		= \int_U u \left(\frac{\partial^2 \phi}{\partial x_i \partial x_j}-\frac{\partial^2 \phi}{\partial x_i \partial x_j}\right) dx \\
		&= \int_U u.0 dx=0,
	\end{align*}
	onde usamos o teorema de comutatividade das derivadas para funções de classe $C^\infty$.
	
	Isso implica que:
	$$\int_U (v_{ij} - v_{ji}) \phi \, dx = 0, \quad \forall \phi \in C_c^\infty(U).$$
	
	Pelo \textbf{Lema Fundamental do Cálculo de Variações} (ou pela unicidade da derivada fraca provada anteriormente (\textbf{Questão 86})), temos que:
	$$v_{ij} - v_{ji} = 0 \quad \text{q.s. em } U$$
	
	Portanto, provamos que:
	$$\boxed{\frac{\partial^2 u}{\partial x_i \partial x_j} = \frac{\partial^2 u}{\partial x_j \partial x_i} \quad \text{q.s. em } U}.$$
}

\sintese{
	A prova transfere a propriedade de comutatividade das derivadas clássicas para o regime fraco através da função de teste. A lógica foi linear:
	$$\text{Definição de } \partial_{ij} u \ \text{e } \partial_{ji} u \to \text{Comutatividade de } \partial_{ij} \phi \to \text{Unicidade q.s.}$$
}

\aprofundamento{
	Este resultado é fundamental para a definição do \textbf{Hessiano Fraco} de uma função em $W^{2,p}(U)$. A simetria do Hessiano é a propriedade que permite a aplicação de técnicas de integração por partes em problemas de minimização de energia, como nos funcionais de Dirichlet. 
	
	Note que esta propriedade não depende de $p$, desde que a função pertença a $W^{2,p}$. No entanto, a interpretação geométrica da simetria das derivadas mistas está intimamente ligada ao fato de que o operador Laplaciano $\Delta = \sum \partial_{ii}$ é a traço do Hessiano. A comutatividade garante que a estrutura de segunda ordem de Sobolev seja análoga à estrutura de funções $C^2$, permitindo a generalização de teoremas de regularidade elíptica.
}

%%%% QUESTÃO 91 %%%%
% --- Bloco Introdutório: Teorema Auxiliar ---
Para a resolução da Questão 91, utilizaremos o seguinte resultado fundamental da teoria de espaços de Sobolev, cuja base técnica de regularização foi construída na \textbf{Questão 4} desta lista.

\begin{teorema}[\textbf{Teorema da Coincidência}]
	Seja $U \subset \mathbb{R}^n$ um aberto e $u \in W^{1,p}(U)$. Então, a derivada fraca $\tfrac{\partial u}{\partial x_j}$ \textbf{coincide} quase sempre com a derivada clássica $\tfrac{\partial u}{\partial x_j}$ em todo ponto $x \in U$ onde esta última existe.
\end{teorema}

\begin{prova}
	Seja $\eta \in C_c^\infty(\mathbb{R}^n)$ uma \textbf{função regularizante} tal que $\supp \eta \subset B_1(0)$, $\eta \ge 0$ e $\int_{B_1(0)} \eta(x) dx = 1$. Para $\varepsilon > 0$, definimos a função regularizante $\eta_\varepsilon(x) = \varepsilon^{-n} \eta(x/\varepsilon)$, a qual possui suporte em $B_\varepsilon(0)$ e satisfaz $\int_{B_\varepsilon(0)} \eta_\varepsilon(x) dx = 1$. Definimos então o domínio:
	$$U_\varepsilon = \{ x \in U : \dist(x, \partial U) > \varepsilon \}$$
	A função suavizada $u_\varepsilon = u * \eta_\varepsilon$ é definida por:
	$$u_\varepsilon(x) = \int_{B_\varepsilon(0)} u(x-y) \eta_\varepsilon(y) dy, \quad x \in U_\varepsilon$$
	Sabemos que $u_\varepsilon \in C^\infty(U_\varepsilon)$. Uma propriedade fundamental da convolução é que a derivada clássica de $u_\varepsilon$ coincide com a convolução da derivada fraca de $u$ com $\eta_\varepsilon$:
	$$\frac{\partial u_\varepsilon}{\partial x_j}(x) = \left( \frac{\partial u}{\partial x_j} * \eta_\varepsilon \right)(x) = \int_{B_\varepsilon(0)} \frac{\partial u}{\partial x_j}(x-y) \eta_\varepsilon(y) dy$$
	
	Como $\frac{\partial u}{\partial x_j} \in L^p(U)$, analisamos a convergência de $\frac{\partial u_\varepsilon}{\partial x_j}$ para $\frac{\partial u}{\partial x_j}$ em $L^p(U_\varepsilon)$. Pela definição de convolução e utilizando o fato de que $\int_{B_\varepsilon(0)} \eta_\varepsilon(y) dy = 1$, temos:
	\begin{align*}
		\left\| \frac{\partial u_\varepsilon}{\partial x_j} - \frac{\partial u}{\partial x_j} \right\|_{L^p(U_\varepsilon)}^p &= \int_{U_\varepsilon} \left| \int_{B_\varepsilon(0)} \eta_\varepsilon(y) \left( \frac{\partial u}{\partial x_j}(x-y) - \frac{\partial u}{\partial x_j}(x) \right) dy \right|^p dx \\
		&\le \int_{U_\varepsilon} \left( \int_{B_\varepsilon(0)} \eta_\varepsilon(y) \left| \frac{\partial u}{\partial x_j}(x-y) - \frac{\partial u}{\partial x_j}(x) \right|^p dy \right) dx
	\end{align*}
	Onde aplicamos a \textbf{Desigualdade de Jensen}, visto que $\eta_\varepsilon$ é uma função não-negativa com integral unitária. Trocando a ordem de integração via Teorema de Fubini, obtemos:
	\begin{align*}
		\left\| \frac{\partial u_\varepsilon}{\partial x_j} - \frac{\partial u}{\partial x_j} \right\|_{L^p(U_\varepsilon)}^p &\le \int_{B_\varepsilon(0)} \eta_\varepsilon(y) \left( \int_{U_\varepsilon} \left| \frac{\partial u}{\partial x_j}(x-y) - \frac{\partial u}{\partial x_j}(x) \right|^p dx \right) dy
	\end{align*}
	A integral interna representa a norma da translação da função $\frac{\partial u}{\partial x_j}$ em $L^p$. Pela propriedade de \textbf{continuidade da translação no $L^p$} ($\lim_{y \to 0} \| f(\cdot - y) - f(\cdot) \|_{L^p} = 0$), a integral interna converge para zero uniformemente para todo $y \in B_\varepsilon(0)$ quando $\varepsilon \to 0$. Logo:
	$$\left\| \frac{\partial u_\varepsilon}{\partial x_j} - \frac{\partial u}{\partial x_j} \right\|_{L^p(U_\varepsilon)} \to 0, \quad \text{quando } \varepsilon \to 0.$$
	
	Como toda sequência convergente em $L^p(U)$ admite uma subsequência que converge quase sempre, existe $\varepsilon_k \to 0$ tal que:
	$$\frac{\partial u_{\varepsilon_k}}{\partial x_j}(x) \to \frac{\partial u}{\partial x_j}(x) \quad \text{q.s. em } U.$$
	
	Por outro lado, suponha que $u$ seja diferenciável no sentido clássico em $x_0 \in U$. Utilizando a definição de $u_\varepsilon$ e a mudança de variável $y = \varepsilon z$, podemos escrever:
	$$u_\varepsilon(x_0) = \int_{B_1(0)} u(x_0 - \varepsilon z) \eta(z) dz$$
	Como $u$ é diferenciável em $x_0$, podemos aplicar a regra de \textbf{derivação sob o sinal da integral} (como discutido na Questão 22) para obter:
	$$\frac{\partial u_\varepsilon}{\partial x_j}(x_0) = \int_{B_1(0)} \frac{\partial u}{\partial x_j}(x_0 - \varepsilon z) \eta(z) dz$$
	Visto que $u$ é diferenciável em $x_0$, a função $\frac{\partial u}{\partial x_j}$ é contínua em $x_0$. Como $\eta$ é limitada e possui suporte compacto em $B_1(0)$, podemos passar ao limite $\varepsilon \to 0$ dentro da integral:
	$$\lim_{\varepsilon \to 0} \frac{\partial u_\varepsilon}{\partial x_j}(x_0) = \int_{B_1(0)} \frac{\partial u}{\partial x_j}(x_0) \eta(z) dz = \frac{\partial u}{\partial x_j}(x_0) \underbrace{\int_{B_1(0)} \eta(z) dz}_{=1}=\underbrace{\frac{\partial u}{\partial x_j}(x_0)}_{\text{clássica}}.$$
	
	Como a sequência $\tfrac{\partial u_{\varepsilon_k}}{\partial x_j}$ converge simultaneamente para a derivada fraca (q.s.) e para a derivada clássica (pontualmente), concluímos que a derivada fraca e a clássica coincidem quase sempre em $U$.
\end{prova}
% --- Início da Questão 91 ---

\questao{Seja $U = B_2(0)$ e:
	\begin{enumerate}[label*=(\alph*)]
		\item definindo
		$$v(x) = \begin{cases} 2 - |x|, & \text{se } |x| \le 1 \\ 1, & \text{se } 1 < |x| \le 2 \end{cases}$$
		mostre que $v \in W^{1,p}(U)$, para todo $1 \le p \le \infty$. Calcule $\nabla v$;
		\item Mostre que $v \notin W^{2,p}(U)$
\end{enumerate}}

\analise{
	\textbf{Análise do Enunciado:}
	A questão exige a prova de que $v \in W^{1,p}(U)$ e $v \notin W^{2,p}(U)$. O ponto crítico acontece na \textbf{fronteira interna (hipersuperfície)} $\{x : |x|=1\}$. No item (a), devemos validar a derivada fraca lidando com a singularidade na origem. No item (b), provaremos que $v \notin W^{2,p}(U)$ utilizando o Teorema de Coincidência para identificar o único candidato a derivada fraca e demonstrando que ele falha na identidade integral para uma função de teste $\phi \equiv 1$ em $B_1(0)$.
	
	\textbf{Fundamentação Teórica:}
	Utilizaremos a definição de \textbf{Derivada Fraca} e o \textbf{Teorema da Divergência}. Para a letra (a), aplicaremos a integração sobre domínios perfurados $\Omega_\varepsilon = B_1(0) \setminus \overline{B_\varepsilon(0)}$ para anular a contribuição da origem. Para a letra (b), usaremos o fato de que, se $v \in W^{2,p}$, sua derivada fraca deve coincidir com a clássica q.s.
	
	\textbf{Estratégias:}
	\begin{enumerate}
		\item \textbf{Validação do Gradiente:} Verificar a continuidade de $v$, definir $\nabla v$ via \textbf{cases} e validar a identidade fraca usando o Teorema da Divergência em $\Omega_\varepsilon$ para tratar a singularidade em $x=0$.
		\item \textbf{Pertinência a $W^{1,p}$:} Justificar via limitação de $v$ e $\nabla v$ em domínio de medida finita.
		\item \textbf{Contradição para $W^{2,p}$:} Identificar o único candidato $w_{ij}$, escolher $\phi \equiv 1$ em $B_1(0)$ e mostrar que a integral de volume do candidato não anula, enquanto o lado esquerdo da identidade fraca é zero.
	\end{enumerate}
}

\solucao{
	\textbf{(a)} Notemos primeiramente que $v$ é contínua, já que em $\partial B_1(0)$ calculando os limites laterais, obtemos:
	$$\lim_{|x| \to 1^-} v(x) = \lim_{|x| \to 1^-} (2 - |x|) = 1 \quad \text{e} \quad \lim_{|x| \to 1^+} v(x) = \lim_{|x| \to 1^+} 1 = 1.$$
	Como os limites coincidem, $v$ é contínua em $B_2(0)$.
	
	A derivada clássica para $v$ é dada por:
	$$\frac{\partial v}{\partial x_j} = \begin{cases} -\frac{x_j}{|x|}, & \text{se } |x| < 1 \\ 0, & \text{se } 1 < |x| < 2 \end{cases}$$
	Assim, propomos que o gradiente fraco seja $\nabla v = -\frac{x}{|x|} \chi_{B_1(0)}$. Para validar, tomamos $\phi \in C_c^\infty(B_2(0))$. Visto que $v$ não é diferenciável em $x=0$, definimos o domínio perfurado $\Omega_\varepsilon = B_1(0) \setminus \overline{B_\varepsilon(0)}$ para $\varepsilon > 0$. A fronteira $\partial \Omega_\varepsilon$ consiste em $\partial B_1(0)$ (normal $\nu$) e $\partial B_\varepsilon(0)$ (normal $-\nu$). Aplicamos o \textbf{Teorema da Divergência}:
	\begin{align*}
		\int_{B_2(0)} v \frac{\partial \phi}{\partial x_j} dx &= \int_{\Omega_\varepsilon} v \frac{\partial \phi}{\partial x_j} dx + \int_{B_2(0) \setminus B_1(0)} v \frac{\partial \phi}{\partial x_j} dx + \int_{B_\varepsilon(0)} v \frac{\partial \phi}{\partial x_j} dx \\
		&= \int_{\partial B_1(0)} v \phi \nu_j d\sigma + \int_{\partial B_\varepsilon(0)} v \phi (-\nu_j) d\sigma - \int_{\Omega_\varepsilon} \frac{\partial v}{\partial x_j} \phi dx + \int_{B_2(0) \setminus B_1(0)} \underbrace{v}_{=1} \frac{\partial \phi}{\partial x_j} dx
	\end{align*}
	Sendo $\phi=0$ em $\partial B_2(0)$, a última integral é $-\int_{\partial B_1(0)} \phi \nu_j d\sigma$. Além disso, como $v$ e $\phi$ são limitados e $\text{Area}(\partial B_\varepsilon(0)) \to 0$ quando $\varepsilon \to 0$, o termo $\int_{\partial B_\varepsilon(0)} v \phi (-\nu_j) d\sigma$ tende a zero. Como $v=1$ em $\partial B_1(0)$, os termos de superfície se anulam:
	$$\int_{B_2(0)} v \frac{\partial \phi}{\partial x_j} dx = \int_{\partial B_1(0)} \phi \nu_j d\sigma - \int_{B_1(0)} \left( -\frac{x_j}{|x|} \right) \phi dx - \int_{\partial B_1(0)} \phi \nu_j d\sigma = \int_{B_1(0)} \frac{x_j}{|x|} \phi dx$$
	Isso prova que 
	$$\boxed{\nabla v = -\frac{x}{|x|} \chi_{B_1(0)} \quad (\ast_1)}.$$ 
	Além disso, como $U=B_2(0)$ é limitado, $|v| \le 2$ e $|\nabla v| \le 1$ q.s., então $v, \nabla v \in L^\infty(B_2(0)) \subset L^p(B_2(0))$. Logo, $$\boxed{v \in W^{1,p}(B_2(0))}.$$
	
	\textbf{(b)} Suponhamos, por absurdo, que $v \in W^{2,p}(B_2(0))$. Pelo \textbf{Teorema de Coincidência} demonstrado anteriormente, a derivada fraca $\tfrac{\partial^2 v}{\partial x_i \partial x_j}$ deve coincidir quase sempre com a derivada clássica. Assim, o único candidato a derivada fraca é:
	$$w_{ij}(x) = \begin{cases} \frac{\partial}{\partial x_i} \left( -\frac{x_j}{|x|} \right) = \frac{\delta_{ij}|x|^2 - x_i x_j}{|x|^3}, & \text{se } |x| < 1 \\ 0, & \text{se } 1 < |x| < 2 \end{cases}$$
	Para que $v \in W^{2,p}(B_2(0))$, a identidade 
	$$\int_{B_2(0)} \frac{\partial v}{\partial x_j} \frac{\partial \phi}{\partial x_i} dx = - \int_{B_2(0)} w_{ij} \phi dx,$$ 
	deve valer para toda $\phi \in C_c^\infty(U)$. Consideremos, em particular, uma $\phi \in C_c^\infty(B_2(0))$ tal que $\phi \equiv 1$ em $B_1(0)$. Para $i=j=1$, por um lado, temos que:
	\begin{align*}
		\int_{B_2(0)} \frac{\partial v}{\partial x_j} \frac{\partial \phi}{\partial x_i} dx &= \int_{B_1(0)} \left( -\frac{x_1}{|x|} \right) \underbrace{\frac{\partial \phi}{\partial x_1}}_{=0} dx + \int_{B_2 \setminus B_1} \underbrace{\frac{\partial v}{\partial x_1}}_{=0} \frac{\partial \phi}{\partial x_1} dx = 0.
	\end{align*}
	Por outro lado, utilizando o candidato $w_{11}$ e a escolha de $\phi$, obtemos:
	$$- \int_U w_{11} \phi dx = - \int_{B_1(0)} \left(\frac{x_1^2 - |x|^2}{|x|^3}\right) \cdot 1 \, dx - \int_{B_2 \setminus B_1} 0 \cdot \phi dx.$$
	Calculando a integral em $B_1(0)$ via coordenadas esféricas ($x = r\xi$ com $\xi \in \partial B_1(0)$) e supondo $n \ge 2$:
	\begin{align*}
		\int_{B_1(0)} \frac{x_1^2 - |x|^2}{|x|^3} dx &= \int_0^1 \int_{\partial B_1(0)} \frac{r^2 \xi_1^2 - r^2}{r^3} r^{n-1} d\sigma_\xi dr 
		= \int_0^1 r^{n-2} \left( \int_{\partial B_1(0)} (\xi_1^2 - 1) d\sigma_\xi \right) dr \\
		&= \left( \int_0^1 r^{n-2} dr \right) \left( \frac{1}{n}\omega_n - \omega_n \right) = \frac{1}{n-1} \omega_n \left( \frac{1-n}{n} \right) = -\frac{\omega_n}{n} \neq 0
	\end{align*}
	Temos, portanto, $0 = \frac{\omega_n}{n}$, o que é um absurdo. Logo, $\boxed{v \notin W^{2,p}(B_2(0))}$ para $n \ge 2$. (Para $n=1$, o argumento falharia pois a integral daria zero, mas a função falha em pertencer a $W^{2,p}$ de qualquer forma devido à singularidade da derivada na fronteira $|x|=1$).
}

\sintese{
	A prova demonstra que a continuidade de $v$ anula os termos de fronteira na primeira derivada, enquanto a singularidade na origem é removível. Contudo, a descontinuidade do gradiente $\nabla v$ na \textbf{fronteira interna} $|x|=1$ gera uma contradição fundamental: o único candidato a derivada fraca de segunda ordem não satisfaz a identidade integral para funções de teste constantes no interior. O fluxo lógico foi:
	$$\text{Limites Laterais} \to \text{Teorema da Divergência em } \Omega_\varepsilon \to \nabla v \in L^p \to \text{Candidato } w_{ij} \to \text{Teste } \phi \equiv 1 \to \text{Contradição}.$$
}

\aprofundamento{
	Este exemplo é a ilustração perfeita do conceito de \textbf{Traço} em Espaços de Sobolev. Para que uma função definida por partes pertença a $W^{2,p}(U)$, não basta que a função seja contínua ($\llbracket v \rrbracket = 0$), é necessário que o traço do gradiente também seja contínuo através da \textbf{fronteira interna} ($\llbracket \nabla v \rrbracket = 0$).
	
	No caso presente, o salto do gradiente $\llbracket \nabla v \rrbracket = 0 - (- \nu) = \nu$ sobre a esfera $\partial B_1(0)$, razão pela qual $v$ falha em pertencer a $W^{2,p}$, ressaltando que a regularidade de Sobolev exige a compatibilidade das derivadas nas junções de domínios.
}

%%%% QUESTÃO 92 %%%%
\questao{Seja $I = (a, b)$, $c \in I$, $v \in L^1(I)$. Defina
	$$w(x) = \int_a^x v(t) dt.$$
	Mostre que:
	\begin{enumerate}[label=(\alph*)]
		\item $w$ é contínua em $I$;
		\item Se $\phi \in C_c^1(I)$ então
		\begin{align*}
			\int_a^b w(t)\phi(t)dt &= \int_a^b \int_a^b \chi_{[a,t]}(s)\phi(t)v(s)dtds \\
			&= \int_a^b v(s) \left( \int_a^b \chi_{[a,t]}(s)\phi(t)dt \right) ds \\
			&= \int_a^b v(s) \left( \int_a^b \chi_{[s,b]}(t)\phi(t)dt \right) ds \\
			&= \int_a^b v(s) \left( \int_s^b \phi(t)dt \right) ds;
		\end{align*}
		\item $v$ é a derivada fraca da função contínua $w$.
	\end{enumerate}
	(Aqui $\chi_E$ é a função característica do conjunto $E$: $\chi_E(t) = 0$, se $t \notin E$ e $\chi_E(t) = 1$ se $t \in E$).}

\analise{
	\textbf{Análise do Enunciado:}
	A questão traz a demonstração de um dos resultados fundamentais da teoria de Sobolev em dimensão $n=1$: a integral de uma função integrável resulta numa função contínua (na verdade, absolutamente contínua) cuja derivada fraca é a própria função integranda.
	
	\textbf{Fundamentação Teórica:}
	Para justificar a passagem do limite para dentro da integral sem assumir que $v$ é contínua, usaremos o \textbf{Teorema da Convergência Dominada de Lebesgue}. Para trocar a ordem das integrais no item (b), aplicaremos o \textbf{Teorema de Fubini-Tonelli}. Por fim, a definição de derivada fraca exige verificar a identidade integral contra uma função de teste $\psi \in C_c^\infty(I)$.
	
	\textbf{Estratégias:}
	\begin{enumerate}
		\item \textbf{Continuidade:} Usar a definição sequencial de continuidade ($x_n \to x$), reescrever a integral no domínio $[a, b]$ via funções características e aplicar o Teorema da Convergência Dominada.
		\item \textbf{Troca da ordem de integração:} Inserir a função característica, aplicar Fubini e observar a equivalência geométrica no domínio de integração: as condições $a \le s \le t$ e $t \le b$ significam que $t \in [s, b]$.
		\item \textbf{Derivada Fraca:} Escolher a função de teste do item (b) como a derivada de uma função suave com suporte compacto ($\phi = \psi'$) e aplicar o Teorema Fundamental do Cálculo clássico na integral interna.
	\end{enumerate}
}

\solucao{
	\textbf{(a) Continuidade de $w$ em $I$:}
	
	Seja $x \in I$ e considere uma sequência $(x_n)_{n \in \mathbb{N}} \subset I$ tal que $x_n \to x$. Podemos reescrever $w(x_n)$ integrando em todo o intervalo $[a, b]$ com o auxílio da função característica:
	$$ w(x_n) = \int_a^{x_n} v(t) dt = \int_a^b \chi_{[a, x_n]}(t) v(t) dt $$
	Como $x_n \to x$, temos pontualmente que $\chi_{[a, x_n]}(t) \to \chi_{[a, x]}(t)$ para todo $t \neq x$. Por outro lado, o ponto $\{x\}$ tem medida de Lebesgue nula, e portanto, $\chi_{[a,x_n]}(t)v(t) \to \chi_{[a, x]}(t) v(t)$ quase sempre (q.s.) em $I$.
	
	Além disso, para todo $n \in \mathbb{N}$, vale a estimativa:
	$$|\chi_{[a, x_n]}(t) v(t)| \le |v(t)| \quad \text{q.s. em } I $$
	Por hipótese, sabemos que $v \in L^1(I)$, e então a função $|v|$ atua como um limitante integrável. Assim, usando o \textbf{Teorema da Convergência Dominada de Lebesgue}, podemos passar o limite para dentro da integral, e obtemos:
	$$ \lim_{n \to \infty} w(x_n) = \int_a^b \left( \lim_{n \to \infty} \chi_{[a, x_n]}(t) v(t) \right) dt = \int_a^b \chi_{[a, x]}(t) v(t) dt = \int_a^x v(t) dt = w(x) $$
	Como $x_n \to x$ implica $w(x_n) \to w(x)$, concluímos que $\boxed{w \in C(I)}$.
	
	\textbf{(b) Identidades Integrais via Fubini:}
	
	Seja $\phi \in C_c^1(I)$. Partimos do lado esquerdo da identidade e escrevemos a integral interna no intervalo $[a,b]$ usando a função característica:
	$$ \int_a^b w(t) \phi(t) dt = \int_a^b \left( \int_a^t v(s) ds \right) \phi(t) dt = \int_a^b \left( \int_a^b \chi_{[a,t]}(s) v(s) ds \right) \phi(t) dt $$
	Como $\phi(t)$ não depende de $s$, podemos colocá-la dentro da integral interna:
	$$ = \int_a^b \left( \int_a^b \chi_{[a,t]}(s) \phi(t) v(s) ds \right) dt $$
	A função integranda $(s,t) \mapsto \chi_{[a,t]}(s) \phi(t) v(s)$ é absolutamente integrável em $I \times I$, pois $\phi$ é limitada (tem suporte compacto) e $v \in L^1(I)$. Pelo \textbf{Teorema de Fubini}, podemos trocar a ordem de integração:
	$$ \int_a^b \left( \int_a^b \chi_{[a,t]}(s) \phi(t) v(s) ds \right) dt = \int_a^b \left( \int_a^b \chi_{[a,t]}(s) \phi(t) v(s) dt \right) ds = \int_a^b v(s) \left( \int_a^b \chi_{[a,t]}(s) \phi(t) dt \right) ds $$
	Agora, analisamos o suporte da função característica. Ter $\chi_{[a,t]}(s) = 1$ significa que $a \le s \le t \le b$. 
	
	Olhando essa mesma desigualdade com foco na variável $t$, vemos que a restrição sobre $t$ é $s \le t \le b$, ou seja, $t \in [s, b]$. Logo, $\chi_{[a,t]}(s) = \chi_{[s,b]}(t)$. Substituindo na integral, ficamos com:
	$$ \int_a^b v(s) \left( \int_a^b \chi_{[a,t]}(s) \phi(t) dt \right) ds = \int_a^b v(s) \left( \int_a^b \chi_{[s,b]}(t) \phi(t) dt \right) ds $$
	Como $\chi_{[s,b]}(t) = 1$ apenas no intervalo $[s, b]$ e $0$ no resto, a integral interna se reduz a:
	$$ = \int_a^b v(s) \left( \int_s^b \phi(t) dt \right) ds $$
	O que conclui as igualdades pedidas.
	
	\textbf{(c) Derivada Fraca:}
	
	Para mostrar que $v$ é a derivada fraca de $w$, precisamos verificar que para qualquer $\psi \in C_c^\infty(I)$, vale:
	$$ \int_a^b w(t) \psi'(t) dt = - \int_a^b v(t) \psi(t) dt $$
	
	Tomemos uma $\psi \in C_c^\infty(I)$ arbitrária. Em particular, $\psi' \in C_c^\infty(I) \subset C_c^1(I)$. Podemos então aplicar o resultado do item (b), substituindo $\phi(t)$ por $\psi'(t)$:
	$$ \int_a^b w(t) \psi'(t) dt = \int_a^b v(s) \left( \int_s^b \psi'(t) dt \right) ds $$
	Como $\psi$ é de classe $C^\infty$, aplicamos o \textbf{Teorema Fundamental do Cálculo} na integral interna:
	$$ \int_s^b \psi'(t) dt = \psi(b) - \psi(s) $$
	Como o suporte de $\psi$ é compacto e contido no interior de $(a,b)$, temos que $\psi(b) = 0$. Logo:
	$$ \int_s^b \psi'(t) dt = 0 - \psi(s) = -\psi(s) $$
	Substituindo na integral iterada:
	$$ \int_a^b w(t) \psi'(t) dt = \int_a^b v(s) \bigl( -\psi(s) \bigr) ds = - \int_a^b v(s) \psi(s) ds = - \int_I v(t) \psi(t) dt$$
	Portanto, a derivada fraca de $w$ é exatamente $v$.
}

\sintese{
	A demonstração conecta de forma direta as noções clássicas e fracas do cálculo unidimensional. O raciocínio consistiu em transferir a derivada para a função de teste usando integrais iteradas, seguindo o esquema lógico:
	$$ w(t) \xrightarrow{\text{Função Característica}} \chi_{[a,t]}(s) \xrightarrow{\text{Fubini}} \chi_{[s,b]}(t) \xrightarrow{\text{Teorema Fundamental do Cálculo}} \text{Derivada Fraca}. $$
}

\aprofundamento{
	Este exercício estabelece um resultado central para o estudo de espaços de Sobolev em dimensão $n=1$: o espaço $W^{1,1}(I)$ coincide com o espaço das funções absolutamente contínuas, cujas derivadas fracas pertencem a $L^1(I)$.
	
	Enquanto em dimensões superiores ($n > 1$) as funções de Sobolev podem ser descontínuas em pontos ou superfícies (dependendo do expoente $p$), em $n=1$, qualquer função com derivada fracamente integrável admite um representante contínuo. Este fato é a base matemática por trás da imersão contínua $W^{1,p}(I) \hookrightarrow C^{0, 1 - 1/p}(\bar{I})$ dada pelo Teorema de Morrey.
}

%%%% QUESTÃO 93 %%%%
\questao{
	(Teoria da Medida e Integração - o operador translação) Sejam $\Omega, U$ abertos tais que $\Omega$ é limitado e $\overline{\Omega} \subset U$ ($\Omega \subset\subset U$). Para cada $h \in \mathbb{R}^n$, $|h| < \dist(\Omega, \partial U)$, defina o operador translação $\tau_h : L^p(U) \to L^p(\Omega)$ como a extensão única do operador $\tau_h : C_0^\infty(U) \subset L^p(U) \to L^p(\Omega)$ dado por $\tau_h u(x) = u(x + h)$ e mostre que
	\begin{enumerate}[label=(\alph*)]
		\item $\tau_h$ é linear, contínua e $\|\tau_h\| = 1$;
		\item Se $u \in L^p(U)$, $v \in L^\infty(U)$ tal que $\supp v \subset\subset \Omega$, então
		$$ \int_U u(\tau_h v - v)dx = \int_U v(\tau_{-h} u - u)dx. $$
		\item Para cada $u \in L^p(U)$:
		$$ \lim_{h \to 0} \|\tau_h u - u\|_{L^p(\Omega)} = 0, \, \forall \Omega \subset\subset U $$
	\end{enumerate}
}

\analise{
	\textbf{Análise do Enunciado:}
	O problema pede para verificar as propriedades básicas do operador de translação $\tau_h$ em espaços $L^p$. A condição geométrica $\Omega \subset\subset U$ garante um \textbf{espaço seguro} para mover as funções sem sair do domínio maior $U$.
	
	\textbf{Fundamentação Teórica:}
	Nos itens (a) e (b), a ferramenta principal é a \textbf{mudança de variáveis na integral}, que mostra que a medida de Lebesgue não muda quando deslocamos uma função. No item (b), usamos a \textbf{Desigualdade de Hölder} para garantir que as integrais existem; como $\Omega$ é limitado e $v \in L^\infty(U)$ tem suporte compacto, a função $v$ pertence a qualquer espaço $L^q$ necessário. No item (c), como não podemos usar a continuidade ponto a ponto diretamente em $L^p$, usamos o fato de que as funções suaves de $C_0^\infty(U)$ são \textbf{densas} em $L^p(U)$ (para $1 \le p < \infty$), junto com o fato de que funções com suporte compacto são uniformemente contínuas.
	
	\textbf{Estratégias:}
	\begin{enumerate}
		\item \textbf{Continuidade e Norma:} Aplicar a mudança $y = x+h$ na integral da norma. A condição do tamanho de $h$ garante que o domínio deslocado continua dentro de $U$, permitindo limitar a integral pela norma em $L^p(U)$.
		\item \textbf{Identidade de Integração:} Fazer a mudança de variáveis na integral do produto. O suporte compacto de $v$ zera a função fora de $\Omega$, permitindo passar o deslocamento de $v$ para $u$ trocando o sinal de $h$.
		\item \textbf{Continuidade Forte:} Construir um argumento padrão de $\varepsilon/3$. Aproximando $u$ por uma função suave $w$, controlamos a translação de $w$ pela continuidade uniforme e obtemos a estimativa usando o item (a).
	\end{enumerate}
}

\solucao{
	\textbf{(a) $\tau_h$ é linear, contínua e $\|\tau_h\| = 1$:}
	
	\textit{Linearidade:} Para quaisquer $u, w \in L^p(U)$ e constantes $\alpha, \beta$:
	$$ \tau_h(\alpha u + \beta w)(x) = (\alpha u + \beta w)(x+h) = \alpha u(x+h) + \beta w(x+h) = \alpha \tau_h u(x) + \beta \tau_h w(x). $$
	
	\textit{Continuidade e Limitação:} Calculamos a norma em $L^p(\Omega)$:
	$$ \|\tau_h u\|_{L^p(\Omega)}^p = \int_\Omega |u(x+h)|^p dx $$
	Fazendo a mudança de variáveis $y = x+h$, temos que $dx = dy$. Como $x \in \Omega$ e $|h| < \dist(\Omega, \partial U)$, o conjunto deslocado continua no interior do domínio maior: $y \in \Omega+h \subset U$. Logo:
	$$ \int_{\Omega+h} |u(y)|^p dy \le \int_U |u(y)|^p dy = \|u\|_{L^p(U)}^p $$
	Portanto, $\|\tau_h u\|_{L^p(\Omega)} \le \|u\|_{L^p(U)}$. Isso mostra que o operador $\tau_h$ é contínuo e sua norma é menor ou igual a 1.
	
	Para provar que $\|\tau_h\| = 1$, precisamos encontrar uma situação onde a desigualdade anterior se torne uma igualdade. Para isso, consideremos $u\neq 0$ tal que $\supp u\subset\Omega+h$. Logo:
	$$ \|u\|_{L^p(U)}^p = \int_U |u(y)|^p dy = \int_{\Omega+h} |u(y)|^p dy = \int_\Omega |u(x+h)|^p dx = \|\tau_h u\|_{L^p(\Omega)}^p $$
	Portanto, 
	$$ \|\tau_h\| = \sup_{u \neq 0} \frac{\|\tau_h u\|_{L^p(\Omega)}}{\|u\|_{L^p(U)}} = 1. $$
	
	\textbf{(b) Identidade Integral:}
	
	Fazendo a mudança de variáveis $y = x+h$ ($dx = dy$) e sabendo que o suporte de $v$ está contido em $\Omega$, a função $\tau_h v(x) = v(x+h)$ só é diferente de zero se $x+h \in \Omega$, ou seja, $x \in \Omega - h$. Pela restrição de $h$, sabemos que $\Omega - h \subset U$. Assim:
	$$ \int_U u(x) \tau_h v(x) dx = \int_{\Omega-h} u(x) v(x+h) dx = \int_\Omega u(y-h) v(y) dy $$
	Como $v\equiv 0$ em $\Omega^{c}$, podemos estender a integração para todo o conjunto $U$ sem alterar o valor, ou seja:
	$$ \int_\Omega u(y-h) v(y) dy = \int_U u(y-h) v(y) dy = \int_U \tau_{-h} u(y) v(y) dy $$
	Subtraindo o termo $\int_U u(x)v(x) dx$ dos dois lados da equação, chegamos ao resultado:
	$$ \int_U u(\tau_h v - v) dx = \int_U v(\tau_{-h} u - u) dx. $$
	
	\textbf{(c) Continuidade forte das translações em $L^p$ ($1 \le p < \infty$):}
	
	Seja $\varepsilon > 0$. Como $\overline{C^{\infty}_0(U)} = L^{p}(U)$, existe uma função suave $w \in C^{\infty}_0(U)$ tal que $\|u - w\|_{L^p(U)} < \tfrac{\varepsilon}{3}$. Além disso, pela desigualdade triangular, temos que:
	$$ \|\tau_h u - u\|_{L^p(\Omega)} \le \|\tau_h u - \tau_h w\|_{L^p(\Omega)} + \|\tau_h w - w\|_{L^p(\Omega)} + \|w - u\|_{L^p(\Omega)} $$
	
	Estudando cada parte:
	\begin{enumerate}
		\item Pela continuidade do item (a), $\|\tau_h(u-w)\|_{L^p(\Omega)} \le \|u-w\|_{L^p(U)} < \tfrac{\varepsilon}{3}$.
		\item Pela escolha de $w$, temos $\|w-u\|_{L^p(\Omega)} \le \|w-u\|_{L^p(U)} < \tfrac{\varepsilon}{3}$.
		\item Como $w$ é suave e tem suporte compacto, ela é uniformemente contínua em $U$ (\textbf{Teorema de Heine-Cantor}). Assim, existe um $\delta > 0$ tal que, se $|h| < \delta$, então $|w(x+h) - w(x)| < \tfrac{\varepsilon}{3|\Omega|^{1/p}}$ para todo $x$. Integrando no domínio limitado $\Omega$, temos:
		$$ \|\tau_h w - w\|_{L^p(\Omega)}^p = \int_\Omega |w(x+h) - w(x)|^p dx < \int_\Omega \frac{\varepsilon^p}{3^p |\Omega|} dx = \frac{\varepsilon^p}{3^p} \implies \|\tau_h w - w\|_{L^p(\Omega)} < \frac{\varepsilon}{3}. $$
	\end{enumerate}
	
	Somando as três partes, concluímos que:
	$$ \boxed{|h|<\delta\implies \|\tau_h u - u\|_{L^p(\Omega)} < \frac{\varepsilon}{3} + \frac{\varepsilon}{3} + \frac{\varepsilon}{3} = \varepsilon},$$
	ou seja,
	$$\boxed{\lim_{h \to 0} \|\tau_h u - u\|_{L^p(\Omega)} = 0}.$$
}

\sintese{
	As propriedades do operador translação vêm diretamente do fato de que a medida de Lebesgue não muda por deslocamentos. A demonstração passa por transferir esse deslocamento da função teste para a função principal (item b) e usar a densidade das funções suaves (item c), conectando a continuidade clássica com a estrutura dos espaços $L^p$. O fluxo lógico foi:
	$$ \text{Invariância da Medida} \implies \text{Identidade Integral} \implies \text{Densidade Suave} \implies \text{Estrutura de } L^p $$
}

\aprofundamento{
	Esta sequência de passos é a base para construir os Espaços de Sobolev e as funções regularizadoras. A identidade do item (b) funciona como uma \textbf{integração por partes discreta}, sendo a ferramenta principal para mostrar que o controle de translações em $L^p$ equivale à existência de derivadas fracas.
	
	Além disso, vale destacar um detalhe importante: a convergência do item (c) falha completamente se $p = \infty$. No espaço $L^\infty(U)$, as funções contínuas não são densas. Podemos mostrar isso de forma exata usando a função característica de um intervalo. 
	
	Por exemplo, considerando $U = (-2, 2)$, $\Omega = (-1, 1)$ e  $u \in L^\infty(U)$ dada por $u(x) = 1$ se $x > 0$ e $u(x) = 0$ se $x \le 0$, ao aplicarmos uma pequena translação $h > 0$, a diferença da função $\tau_h u(x) = u(x+h)$ por $u(x)$ é dada por:
	$$ \tau_h u(x) - u(x) = \begin{cases} 1, & \text{se } -h < x \le 0 \\ 0, & \text{caso contrário} \end{cases} $$
	Independentemente de quão pequeno seja $h$, temos que:
	$$ \|\tau_h u - u\|_{L^\infty(\Omega)} = 1 $$
	Portanto, $\lim_{h \to 0} \|\tau_h u - u\|_{L^\infty(\Omega)} = 1 \neq 0$, provando que o resultado não é válido para $p=\infty$.
}

%%%% QUESTÃO 94 %%%%
\questao{
	O operador $\tau_h : L^p(U) \to L^p(\Omega)$ definido no exercício (93), de fato está bem definido de $\tau_h : W^{1,p}(U) \to W^{1,p}(\Omega)$ e
	$$ \frac{\partial(\tau_h u)}{\partial x_j} = \tau_h \left( \frac{\partial u}{\partial x_j} \right). $$
	Conclua que $\lim_{h \to 0} \|\tau_h u - u\|_{W^{1,p}(\Omega)} = 0$.
}

\analise{
	\textbf{Análise do Enunciado:}
	O problema nos pede para elevar o operador de translação $\tau_h$ de espaços $L^p$ para os espaços de Sobolev $W^{1,p}$. Para isso, precisamos provar duas propriedades estruturais: a comutação entre a translação e o operador de derivação fraca, e a continuidade forte das translações na norma de Sobolev.
	
	\textbf{Fundamentação Teórica:}
	A ferramenta para demonstrar a boa definição e a identidade operacional é a \textbf{definição de derivada fraca por dualidade} combinada com a mudança de variáveis. Como a condição geométrica $\Omega \subset\subset U$ garante que o suporte transladado das funções teste de $\Omega$ permanece confinado dentro do aberto maior $U$, podemos transferir a translação para as funções teste suaves. Para a segunda parte, utilizaremos a própria definição da norma de $W^{1,p}$, o que nos permitirá \textbf{reduzir o problema ao caso $L^p$} já demonstrado no item (c) da \textbf{Questão 93}.
	
	\textbf{Estratégias:}
	\begin{enumerate}
		\item \textbf{Identidade de Comutação:} Aplicar a definição de derivada fraca para a função transladada $\tau_h u$ usando uma função teste $\phi \in C_0^\infty(\Omega)$. Efetuar uma mudança de variáveis na integral para deslocar o argumento e identificar a derivada fraca de $u$ agindo em $U$.
		\item \textbf{Convergência Forte:} Decompor a norma de Sobolev $\|\tau_h u - u\|_{W^{1,p}(\Omega)}$ na soma das normas $L^p$ da função e de suas respectivas derivadas fracas. Aplicar o limite $h \to 0$ componente a componente usando o resultado de continuidade forte obtido na questão anterior.
	\end{enumerate}
}

\solucao{
	\textbf{(a) Boa definição e Comutação da Derivada Fraca:}
	
	Seja $u \in W^{1,p}(U)$. \textbf{Pelo item (a) da Questão 93}, já sabemos que $u \in L^p(U) \implies \tau_h u \in L^p(\Omega)$. Para mostrar que $\tau_h u \in W^{1,p}(\Omega)$, precisamos verificar a existência de suas derivadas fracas em $\Omega$.
	
	Seja $\phi \in C_0^\infty(\Omega)$ uma função teste arbitrária. Pela definição de translação, temos que:
	$$ \int_\Omega \tau_h u(x) \frac{\partial \phi}{\partial x_j}(x) dx = \int_\Omega u(x+h) \frac{\partial \phi}{\partial x_j}(x) dx = \int_{\Omega+h} u(y) \frac{\partial \phi}{\partial x_j}(y-h) dy,$$
	onde fizemos a mudança de variáveis linear $y = x+h$ ($dx = dy$), e com isso o domínio de integração passa a ser o conjunto transladado $\Omega+h$.
	
	Agora, definamos a função $\psi(y) = \phi(y-h)$. Como $\supp \phi \subset\subset \Omega$ e $|h| < \dist(\Omega, \partial U)$, o suporte de $\psi$ satisfaz 
	$$\supp \psi = \supp \phi + h \subset\subset \Omega+h \subset U.$$ 
	Portanto, $\psi \in C_0^\infty(U)$. Além disso, pela regra da cadeia clássica, $\tfrac{\partial \psi}{\partial y_j}(y) = \tfrac{\partial \phi}{\partial x_j}(y-h)$. Substituindo na integral e estendendo o domínio para todo o aberto $U$ (visto que $\psi \equiv 0$ fora de $\Omega+h$), obtemos:
	$$ \int_{\Omega+h} u(y) \frac{\partial \phi}{\partial x_j}(y-h) dy = \int_U u(y) \frac{\partial \psi}{\partial y_j}(y) dy. $$
	Como $u \in W^{1,p}(U)$, usamos a definição de sua derivada fraca $\tfrac{\partial u}{\partial y_j}$ em $U$:
	$$ \int_U u(y) \frac{\partial \psi}{\partial y_j}(y) dy = - \int_U \frac{\partial u}{\partial y_j}(y) \psi(y) dy = - \int_{\Omega+h} \frac{\partial u}{\partial y_j}(y) \phi(y-h) dy. $$
	Retornando a nossa mudança de variáveis ($y = x+h$, $dy = dx$), temos que:
	$$ - \int_{\Omega+h} \frac{\partial u}{\partial y_j}(y) \phi(y-h) dy = - \int_\Omega \frac{\partial u}{\partial x_j}(x+h) \phi(x) dx = - \int_\Omega \tau_h\left(\frac{\partial u}{\partial x_j}\right)(x) \phi(x) dx. $$
	Portanto:
	$$ \int_\Omega \tau_h u(x) \frac{\partial \phi}{\partial x_j}(x) dx = - \int_\Omega \tau_h\left(\frac{\partial u}{\partial x_j}\right)(x) \phi(x) dx, \quad \forall \phi \in C_0^\infty(\Omega). $$
	Isso prova exatamente que a derivada fraca de $\tau_h u$ em relação a $x_j$ existe em $\Omega$ e coincide com a translação da derivada fraca de $u$, ou seja:
	$$ \boxed{\frac{\partial(\tau_h u)}{\partial x_j} = \tau_h \left( \frac{\partial u}{\partial x_j} \right)}. $$
	Como $u \in W^{1,p}(U)$, temos $\tfrac{\partial u}{\partial x_j} \in L^p(U)$, o que garante pelo \textbf{item (a) da Questão 93} que $\tau_h \left( \tfrac{\partial u}{\partial x_j} \right) \in L^p(\Omega)$. Concluímos que $\tau_h : W^{1,p}(U) \to W^{1,p}(\Omega)$ está bem definido.
	
	\textbf{(b) Conclusão do Limite Forte:}
	
	Para $1 \le p < \infty$, sabemos que:
	$$ \|\tau_h u - u\|_{W^{1,p}(\Omega)}^p = \|\tau_h u - u\|_{L^p(\Omega)}^p + \sum_{j=1}^n \left\| \frac{\partial(\tau_h u)}{\partial x_j} - \frac{\partial u}{\partial x_j} \right\|_{L^p(\Omega)}^p =\|\tau_h u - u\|_{L^p(\Omega)}^p + \sum_{j=1}^n \left\| \tau_h\left(\frac{\partial u}{\partial x_j}\right) - \frac{\partial u}{\partial x_j} \right\|_{L^p(\Omega)}^p\quad (\ast_1)$$
	onde usamos o que provamos no item (a) acima.
	
	Por outro lado, como $u \in L^p(U)$ e cada $\tfrac{\partial u}{\partial x_j} \in L^p(U)$, pelo \textbf{item (c) da Questão 93}, temos que:
	$$ \lim_{h \to 0} \|\tau_h u - u\|_{L^p(\Omega)} = 0 \quad \text{e} \quad \lim_{h \to 0} \left\| \tau_h\left(\frac{\partial u}{\partial x_j}\right) - \frac{\partial u}{\partial x_j} \right\|_{L^p(\Omega)} = 0, \quad \forall j=1, \dots, n. $$
	Disto e da igualdade em $(\ast_1)$, segue que:
	$$ \boxed{\lim_{h \to 0} \|\tau_h u - u\|_{W^{1,p}(\Omega)} = 0}. $$
}

\sintese{
	A estrutura linear de diferenciação fraca comuta com isometrias geométricas como a translação. O fluxo lógico estabelecido foi:
	$$ \text{Dualidade em } \Omega \implies \text{Suporte Compacto em } U \implies \text{Comutação } \partial_j \tau_h = \tau_h \partial_j \implies \text{Convergência via } L^p $$
}

\aprofundamento{
	É de fundamental importância salientar que a convergência forte do limite da norma de Sobolev \textbf{falha categoricamente no caso $p = \infty$}. Enquanto a imersão de Morrey garante que $u \in W^{1,\infty}(U)$ é localmente Lipschitz contínua (o que força a convergência uniforme da função $\|\tau_h u - u\|_{L^\infty(\Omega)} \to 0$), as suas derivadas fracas $\frac{\partial u}{\partial x_j}$ são apenas elementos de $L^\infty(U)$ e podem apresentar descontinuidades de salto, onde a translação falha.
	
	Para demonstrar essa quebra de convergência de forma explícita, construímos o seguinte contra-exemplo unidimensional. Considere os abertos $U = (-2,2)$ e $\Omega = (-1,1)$, com $\Omega \subset\subset U$, e a função módulo:
	$$ u(x) = |x|. $$
	Como $u$ é globalmente Lipschitz em $U$, temos $u \in W^{1,\infty}(U)$. A sua derivada fraca é dada quase sempre pela função sinal:
	$$ u'(x) = \text{sgn}(x) = \begin{cases} 1, & \text{se } x > 0 \\ -1, & \text{se } x < 0 \end{cases} $$
	Fixemos um deslocamento positivo infinitesimal $0 < h < 1$. A derivada fraca da função transladada $\tau_h u$ avaliada em $\Omega$ é expressa por $\tau_h (u')(x) = \text{sgn}(x+h)$. Vamos computar pontualmente a diferença das derivadas, $\tau_h (u')(x) - u'(x)$, dividindo o domínio $\Omega = (-1,1)$ em três subintervalos disjuntos:
	\begin{enumerate}
		\item Para $x \in (-1, -h)$: Temos $x < 0$ e $x+h < 0$. Portanto:
		$$ \tau_h (u')(x) - u'(x) = \text{sgn}(x+h) - \text{sgn}(x) = -1 - (-1) = 0. $$
		
		\item Para $x \in (-h, 0)$: Temos $x < 0$, mas $x+h > 0$. Portanto:
		$$ \tau_h (u')(x) - u'(x) = \text{sgn}(x+h) - \text{sgn}(x) = 1 - (-1) = 2. $$
		
		\item Para $x \in (0, 1)$: Temos $x > 0$ e $x+h > 0$. Portanto:
		$$ \tau_h (u')(x) - u'(x) = \text{sgn}(x+h) - \text{sgn}(x) = 1 - 1 = 0. $$
	\end{enumerate}
	
	Reunindo os três casos, a função diferença das derivadas fracas em $\Omega$ assume o comportamento:
	$$ \tau_h (u')(x) - u'(x) = 2 \cdot \chi_{(-h,0)}(x), \quad \text{q.s. em } \Omega. $$
	Como o intervalo $(-h, 0)$ possui medida de Lebesgue estritamente positiva ($|\mu| = h > 0$), o supremo essencial desta diferença não é afetado pelo comportamento em conjuntos de medida nula. Calculando a norma em $L^\infty(\Omega)$, obtemos:
	$$ \|\tau_h (u') - u'\|_{L^\infty(\Omega)} = \sup_{x \in (-1,1)} |\text{sgn}(x+h) - \text{sgn}(x)| = 2. $$
	
	Portanto, ao avaliarmos o limite da norma completa de Sobolev em $W^{1,\infty}(\Omega)$, a parcela correspondente ao gradiente impede o decaimento:
	$$ \|\tau_h u - u\|_{W^{1,\infty}(\Omega)} = \|\tau_h u - u\|_{L^\infty(\Omega)} + \|\tau_h (u') - u'\|_{L^\infty(\Omega)} \to 0 + 2 = 2 \neq 0, \quad \text{quando } h \to 0^+. $$
	Isso encerra a justificativa analítica de que a continuidade forte do operador translação fica estritamente restrita ao intervalo $1 \le p < \infty$.
}

%%%% QUESTÃO 95 %%%%
\questao{(Teoria da Medida e Integração - ainda o operador translação) Sejam $$U_0 = \{x = (z, x_n) \in \mathbb{R}^n : z \in \mathbb{R}^{n-1}, |z| < r, x_n \in \mathbb{R}\},$$ um ``cilindro vertical'' em $\mathbb{R}^n$, uma aplicação $\gamma : B \to \mathbb{R}$ contínua, onde $B$ é a bola com centro na origem de $\mathbb{R}^{n-1}$ e raio $r$. Considere $U = \{x = (z, x_n) \in U_0 : z \in B \text{ e } x_n > \gamma(z)\}$. Mostre que:
	\begin{enumerate}[label=(\alph*)]
		\item $U_0 = B \times \mathbb{R}$;
		\item Para cada $\lambda > 0$, $\tau_{\lambda e_n} : L^p(U) \to L^p(U)$ está bem definido, e como no exercício anterior, é linear, contínuo e tem norma 1;
		\item Se $h$ não é múltiplo positivo de $e_n$, $\tau_h : L^p(U) \to L^p(U)$ não está bem definido;
		\item Considere $V_\lambda = \{x = (z, x_n) \in U_0 : z \in B \text{ e } x_n > \gamma(z) - \lambda\}$. Observe que $U \subset V_\lambda$ e que $\tau_{\lambda e_n} : L^p(V_\lambda) \to L^p(U)$ e para cada $v \in L^p(V_\lambda)$: $\lim_{t \to 0} \|\tau_{t e_n} v - v\|_{L^p(U)} = 0$.
\end{enumerate}}

\analise{
	\textbf{Análise do Enunciado:}	O problema requer a caracterização geométrica de um cilindro vertical $U_0$ e a análise de um operador de translação $\tau_{\lambda e_n}$ atuando sobre o espaço $L^p(U)$, onde $U$ é o domínio situado acima do gráfico de uma função contínua $\gamma$. A tese central é provar que a estabilidade do operador de translação $\tau_h$ é estritamente dependente da direção do vetor $h$: a translação na direção positiva do eixo $x_n$ preserva a pertinência ao domínio $U$, enquanto qualquer componente lateral ou vertical não-positiva invalida a boa definição do operador. Por fim, analisa-se a continuidade forte da translação utilizando um domínio expandido $V_\lambda$.
		
	\textbf{Fundamentação Teórica:} A fundamentação baseia-se na teoria de integração de Lebesgue e nas propriedades de espaços $L^p(\Omega)$. A referência fundamental é \textbf{Brezis (Capítulo 4)}, especificamente no que tange à continuidade da translação em $L^p$ e ao teorema de mudança de variáveis. O conceito de ``bem definido'' para $\tau_h: L^p(\Omega) \to L^p(\Omega)$ implica que a condição $\chi_\Omega(x) = 1 \implies \chi_\Omega(x+h) = 1$ deve valer quase sempre. Para o \textbf{item (d)}, utilizaremos a técnica de extensão por zero para reduzir o problema ao caso clássico de $L^p(\mathbb{R}^n)$.
	
	\textbf{Estratégias:}
	\begin{enumerate}
		\item \textbf{Identificação Geométrica}: No \textbf{item (a)}, correlacionar a condição $|z| < r$ em $\mathbb{R}^{n-1}$ com a definição da bola aberta $B$ para estabelecer a identidade de produto cartesiano.
		\item \textbf{Operador de Translação Positiva}: No \textbf{item (b)}, provar que $\lambda > 0$ garante a inclusão $U + \lambda e_n \subset U$. A norma unitária será provada via mudança de variáveis e análise de funções com suporte afastado da fronteira.
		\item \textbf{Análise de Instabilidade}: No \textbf{item (c)}, dividir a prova em dois casos (componente lateral $h' \neq 0$ e componente vertical $h_n \le 0$), demonstrando que, em ambos, a translação mapeia conjuntos de medida positiva para fora de $U$.
		\item \textbf{Continuidade Forte via Extensão}: No \textbf{item (d)}, provar que $x \in U \implies x + \lambda e_n \in V_\lambda$ e, posteriormente, utilizar a extensão por zero de $v \in L^p(V_\lambda)$ para $\mathbb{R}^n$ a fim de aplicar o teorema de continuidade da translação em $L^p(\mathbb{R}^n)$.
	\end{enumerate}
}
	
\solucao{
	\textbf{(a)} Seja $U_0 = \{x = (z, x_n) \in \mathbb{R}^n : z \in \mathbb{R}^{n-1}, |z| < r, x_n \in \mathbb{R}\}$. 
	Por definição, a bola aberta $B$ em $\mathbb{R}^{n-1}$ com centro na origem e raio $r$ é o conjunto $B = \{z \in \mathbb{R}^{n-1} : |z| < r\}$.
	O produto cartesiano $B \times \mathbb{R}$ é definido como:
	$$B \times \mathbb{R} = \{(z, x_n) : z \in B, x_n \in \mathbb{R}\} = \{(z, x_n) : z \in \mathbb{R}^{n-1}, |z| < r, x_n \in \mathbb{R}\}$$
	Comparando as duas expressões, observamos a identidade exata entre os conjuntos. Portanto,
	$$\boxed{U_0 = B \times \mathbb{R} \quad (\ast)}$$
	
	\textbf{(b)} Seja $\lambda > 0$ e $e_n = (0, 0, \dots, 1)$ o vetor unitário na direção do último eixo. O operador translação é definido por $(\tau_{\lambda e_n} u)(x) = u(x + \lambda e_n)$.
	
	\textit{1. Bem definido}:\par
	Para que $\tau_{\lambda e_n} u$ esteja bem definida em $U$, devemos garantir que para todo $x \in U$, o ponto $x + \lambda e_n$ também pertença a $U$. 
	Se $x = (z, x_n) \in U$, então, pelas hipóteses:
	\begin{align*}
		z &\in B \ \text{e } x_n > \gamma(z)
	\end{align*}
	O ponto transladado é $x + \lambda e_n = (z, x_n + \lambda)$. Como $\lambda > 0$, temos que $x_n + \lambda > x_n > \gamma(z)$.
	Já que $z$ permanece inalterado, então $z \in B$. Logo, $(z, x_n + \lambda) \in U$. O operador $\tau_{\lambda e_n}$ mapeia $U$ em $U$.
	
	\textit{2. Lineariedade}:\par
	Para quaisquer $u, v \in L^p(U)$ e $\alpha, \beta \in \mathbb{R}$:
	\begin{align*}
		(\tau_{\lambda e_n}(\alpha u + \beta v))(x) = (\alpha u + \beta v)(x + \lambda e_n) = \alpha u(x + \lambda e_n) + \beta v(x + \lambda e_n) = \alpha (\tau_{\lambda e_n} u)(x) + \beta (\tau_{\lambda e_n} v)(x).
	\end{align*}
	Donde temos a linearidade.
	
	\textit{3. Continuidade e Norma}:\par
	Calculamos a norma de $\tau_{\lambda e_n} u$ no espaço $L^p(U)$:
	$$\|\tau_{\lambda e_n} u\|_{L^p(U)}^p = \int_U |u(x + \lambda e_n)|^p dx$$
	Efetuamos a mudança de variáveis $y = x + \lambda e_n$. O determinante do Jacobiano desta transformação é $1$. O domínio de integração $U$ transforma-se em $U + \lambda e_n$:
	$$\|\tau_{\lambda e_n} u\|_{L^p(U)}^p = \int_{U + \lambda e_n} |u(y)|^p dy$$
	Observamos que $U + \lambda e_n = \{ (z, x_n + \lambda) : z \in B, x_n > \gamma(z) \} = \{ (z, \xi) : z \in B, \xi > \gamma(z) + \lambda \}$.
	Como $\lambda > 0$, temos $\gamma(z) + \lambda > \gamma(z)$, o que implica a inclusão estrita:
	$$\boxed{U + \lambda e_n \subsetneq U \quad (\ast\ast)}$$
	Portanto, a integral sobre o subconjunto é menor ou igual à integral sobre o todo, ou seja:
	\begin{align*}
		\|\tau_{\lambda e_n} u\|_{L^p(U)}^p=\int_{U + \lambda e_n} |u(y)|^p dy \leq \int_U |u(y)|^p dy = \|u\|_{L^p(U)}^p \implies	\|\tau_{\lambda e_n} u\|_{L^p(U)} \leq \|u\|_{L^p(U)}
	\end{align*}
	Isso prova que o operador é contínuo e que $\|\tau_{\lambda e_n}\| \leq 1$.
	
	Para provar que $\|\tau_{\lambda e_n}\| = 1$, dado que já estabelecemos que $\|\tau_{\lambda e_n}\| \le 1$, resta demonstrar que o supremo não é estritamente menor que $1$. Tomemos funções $u \in C_c^\infty(U)$ cujo suporte satisfaça a inclusão $\text{supp}(u) \subset U + \lambda e_n$. Para tais funções, a translação preserva a norma integral, pois:
	$$\int_{U + \lambda e_n} |u(y)|^p dy = \int_U |u(y)|^p dy$$
	ou, equivalentemente, $\|\tau_{\lambda e_n} u\|_{L^p(U)} = \|u\|_{L^p(U)}$. Consequentemente, segue que 
	$$\boxed{\|\tau_{\lambda e_n}\| = 1.}$$
	
	\textbf{(c)} Seja $h = (h', h_n) \in \mathbb{R}^n$ com $h' \in \mathbb{R}^{n-1}$. Se $h$ não é um múltiplo positivo de $e_n$, então ou $h' \neq 0$ ou $h_n \le 0$.
	
	\textit{Caso 1: $h' \neq 0$ (Componente lateral)}\par
	Como $B$ é uma bola aberta, se $h' \neq 0$, existe um conjunto $A \subset B$ de medida positiva tal que $A + h' \not\subset B$. Para qualquer $z \in A$ tal que $z + h' \notin B$ e qualquer $x_n > \gamma(z)$, o ponto $x = (z, x_n) \in U$, mas $x + h = (z + h', x_n + h_n) \notin U_0$. Logo, $x+h \notin U$.
	
	\textit{Caso 2: $h' = 0$ e $h_n \le 0$ (Componente vertical não-positiva)}\par
	Se $h' = 0$, temos $x+h = (z, x_n + h_n)$. 
	\begin{itemize}
		\item Se $h_n = 0$, $\tau_h$ é a identidade (excluída por hipótese).
		\item Se $h_n < 0$, considere a faixa $\mathcal{F} = \{ (z, x_n) \in U : \gamma(z) < x_n < \gamma(z) - h_n \}$. Esta faixa possui medida positiva. Para qualquer $x \in \mathcal{F}$, temos $x_n + h_n < \gamma(z)$, logo $x+h \notin U$.
	\end{itemize}
	Em ambos os casos, a condição $x \in U \implies x+h \in U$ é violada em um conjunto de medida positiva. Portanto, $\tau_h$ não está bem definido em $L^p(U)$.
	
	\textbf{(d)} Seja $V_\lambda = \{ (z, x_n) \in U_0 : z \in B, x_n > \gamma(z) - \lambda \}$.
	
	\textit{1. Bem definido}:\par
	Se $x = (z, x_n) \in U$, então $z \in B$ e $x_n > \gamma(z)$. Para $\lambda > 0$, o ponto transladado $x + \lambda e_n = (z, x_n + \lambda)$ satisfaz $x_n + \lambda > \gamma(z) > \gamma(z) - \lambda$. Logo, $x + \lambda e_n \in V_\lambda$, o que torna $\tau_{\lambda e_n} : L^p(V_\lambda) \to L^p(U)$ bem definido.
	
	\textit{2. Demonstração do limite}:\par
	Seja $v \in L^p(V_\lambda)$. Definimos a extensão por zero $\tilde{v} \in L^p(\mathbb{R}^n)$ como $\tilde{v}(x) = v(x)$ se $x \in V_\lambda$ e $0$ caso contrário. Para $t > 0$ suficientemente pequeno ($t < \lambda$), temos $x \in U \implies x + t e_n \in V_\lambda$. Assim, para $x \in U$:
	$$(\tau_{t e_n} v)(x) = v(x + t e_n) = \tilde{v}(x + t e_n) = (\tau_{t e_n} \tilde{v})(x)$$
	A norma em $L^p(U)$ satisfaz:
	$$\|\tau_{t e_n} v - v\|_{L^p(U)}^p = \int_U |\tilde{v}(x + t e_n) - \tilde{v}(x)|^p dx \le \int_{\mathbb{R}^n} |\tilde{v}(x + t e_n) - \tilde{v}(x)|^p dx$$
	Pela continuidade da translação em $L^p(\mathbb{R}^n)$, $\lim_{h \to 0} \|\tau_h \tilde{v} - \tilde{v}\|_{L^p(\mathbb{R}^n)} = 0$. Aplicando para $h = t e_n$, obtemos:
	$$\boxed{\lim_{t \to 0} \|\tau_{t e_n} v - v\|_{L^p(U)} = 0}.$$
}

\sintese{
	A prova estabelece a dependência rigorosa entre a geometria do domínio e a estabilidade do operador de translação. A direção do vetor $h$ determina a boa definição do operador, enquanto a expansão do domínio para $V_\lambda$ permite recuperar a continuidade forte da translação, essencial para a regularização de funções de Sobolev.
}

\aprofundamento{
	Este resultado é a base para a construção de \textbf{Extensões de Sobolev}. Ao transladar a função para ``dentro'' do domínio, criamos a margem necessária para aplicar a convolução sem sair de $U$. O domínio $V_\lambda$ atua como um colchão de segurança que garante que a função regularizada $u_\varepsilon$ seja bem definida em $U$. Este procedimento é análogo ao método de reflexão, mas utiliza a estrutura de cilindro para simplificar a geometria.
}

%%%% QUESTÃO 96 %%%%
\questao{(Sequências regularizantes - Caso $U = \mathbb{R}^n_+$) Seja $\eta \in C_c^\infty(\mathbb{R}^n)$, $eta$ como \textbf{na questão 4}, ou seja, $\eta > 0$ em $B_1(0)$, $\text{supp } \eta = \overline{B_1(0)}$ e $\int_{\mathbb{R}^n} \eta(x) dx = 1$. Para cada $\varepsilon > 0$ defina $\eta_\varepsilon(x) = \varepsilon^{-n} \eta(\tfrac{x}{\varepsilon})$. Mostre que:
	\begin{enumerate}[label=(\alph*)]
		\item Se $\mathbb{R}^n_+ = \{x \in \mathbb{R}^n : \langle x, e_n \rangle > 0\}$ e $f \in L_{loc}^1(\mathbb{R}^n_+)$ então
		$$f_\varepsilon(x) = \int_{\mathbb{R}^n_+} \eta_\varepsilon(x-y)(\tau_{2\varepsilon e_n} f)(y) dy$$
		está bem definida em $\mathbb{R}^n_+$ e é de classe $C^\infty(\mathbb{R}^n_+)$, com $f_\varepsilon(x) = [\eta_\varepsilon * (\tau_{2\varepsilon e_n} f)](x)$;
		\item Se $f$ é contínua, $f_\varepsilon$ converge uniformemente nas partes compactas de $\mathbb{R}^n_+$ para $f$;
		\item Se $f \in C^m(\mathbb{R}^n_+)$ então
		$$D^\alpha f_\varepsilon(x) = \int_{\mathbb{R}^n_+} \eta_\varepsilon(x-y)(\tau_{2\varepsilon e_n} D^\alpha f)(y) dy, \quad \forall |\alpha| \le m;$$
		\item Se $f \in L^p(\mathbb{R}^n_+)$, então $f_\varepsilon \in L^p(\mathbb{R}^n_+)$ e $\|f_\varepsilon - f\|_{L^p(\mathbb{R}^n_+)} \to 0$ quando $\varepsilon \to 0$;
		\item Se $u \in W^{1,p}(\mathbb{R}^n_+)$, então $u_\varepsilon \in W^{1,p}(\mathbb{R}^n_+)$ e $\|u_\varepsilon - u\|_{W^{1,p}(\mathbb{R}^n_+)} \to 0$.
\end{enumerate}}

\analise{
	\textbf{Análise do Enunciado:}
	O problema propõe a construção de uma sequência de aproximações suaves para funções definidas no semi-espaço $\mathbb{R}^n_+$. A particularidade reside na definição de $f_\varepsilon$, que não utiliza a convolução simples $f * \eta_\varepsilon$, mas sim a convolução da translação $\tau_{2\varepsilon e_n} f$ com o núcleo $\eta_\varepsilon$. O objetivo é provar que essa modificação geométrica permite a regularização no bordo sem a necessidade de extensões arbitrárias da função para fora do domínio.
	
	\textbf{Fundamentação Teórica:}
	A prova fundamenta-se nas propriedades da convolução de funções $L_{loc}^1$ com núcleos $C_c^\infty$, resultando em funções $C^\infty$. A convergência em $L^p$ e a uniformidade em compactos decorrem da teoria de \textbf{Identidades Aproximadas}. Para o item (e), utilizaremos a linearidade do operador de derivação fraca e a estabilidade da norma de Sobolev sob regularização.
	
	\textbf{Estratégias:}
	\begin{enumerate}
		\item \textbf{Análise de Suporte}: No item (a), provaremos que para $x \in \mathbb{R}^n_+$, a integral de convolução $\eta_\varepsilon * (\tau_{2\varepsilon e_n} f)$ coincide com a integral sobre $\mathbb{R}^n_+$, pois o suporte de $\eta_\varepsilon(x-y)$ e a definição de $\tau_{2\varepsilon e_n}$ garantem que $y \in \mathbb{R}^n_+$.
		\item \textbf{Convergência Uniforme}: No item (b), utilizaremos a continuidade uniforme de $f$ em compactos e a propriedade $\int \eta_\varepsilon = 1$ para estimar $|f_\varepsilon(x) - f(x)|$.
		\item \textbf{Troca de Derivação}: No item (c), aplicaremos a propriedade de que a derivada da convolução é a convolução da derivada.
		\item \textbf{Estimativas de $L^p$}: No item (d), utilizaremos a desigualdade de Young para convoluções $$\|f * g\|_p \le \|f\|_1 \|g\|_p$$ e a continuidade da translação em $L^p$.
		\item \textbf{Soma de Normas}: No item (e), aplicaremos o resultado do item (d) tanto para a função $u$ quanto para suas derivadas fracas $\partial_j u$.
	\end{enumerate}
}

\solucao{
	\textbf{(a)} Seja $x \in \mathbb{R}^n_+$, ou seja, $x_n > 0$. Definimos a função $\tilde{f} = \tau_{2\varepsilon e_n} f$ estendida por zero fora de $\mathbb{R}^n_+$. O núcleo $\eta_\varepsilon(x-y)$ é não nulo apenas se $|x-y| < \varepsilon$, o que implica $y_n > x_n - \varepsilon$. 
	
	Observemos que para qualquer $y$ tal que $\eta_\varepsilon(x-y) \neq 0$, temos $y + 2\varepsilon e_n$ com componente vertical 
	$$(y_n + 2\varepsilon) > x_n - \varepsilon + 2\varepsilon = x_n + \varepsilon.$$
	Como $x_n > 0$, segue que $y_n + 2\varepsilon > \varepsilon > 0$, logo $y + 2\varepsilon e_n \in \mathbb{R}^n_+$. Isso garante que $\tau_{2\varepsilon e_n} f(y) = f(y + 2\varepsilon e_n)$ está bem definida para todo $y$ no suporte de $\eta_\varepsilon(x-\cdot)$.
	
	Assim, a integral $f_\varepsilon(x) = \int_{\mathbb{R}^n_+} \eta_\varepsilon(x-y) \tilde{f}(y) dy$ é a convolução clássica $(\eta_\varepsilon * \tilde{f})(x)$. Visto que $\eta_\varepsilon \in C_c^\infty(\mathbb{R}^n)$ e $\tilde{f} \in L_{loc}^1(\mathbb{R}^n)$, a teoria de convolução garante que $f_\varepsilon \in C^\infty(\mathbb{R}^n_+)$.
	
	\textbf{(b)} Seja $K \subset \mathbb{R}^n_+$ um compacto. Existe $\delta > 0$ tal que $\dist(K, \partial \mathbb{R}^n_+) > \delta$. Para $\varepsilon < \delta$, e qualquer $x \in K$, a bola $B_\varepsilon(x)$ está contida em $\mathbb{R}^n_+$. Assim, a integral de $f_\varepsilon$ simplifica-se para:
	$$f_\varepsilon(x) = \int_{B_\varepsilon(x)} \eta_\varepsilon(x-y) f(y + 2\varepsilon e_n) dy$$
	Utilizando o fato de que $\int_{B_\varepsilon(x)} \eta_\varepsilon(x-y) dy = 1$, podemos escrever:
	\begin{align*}
		|f_\varepsilon(x) - f(x)| &= \left| \int_{B_\varepsilon(x)} \eta_\varepsilon(x-y) f(y + 2\varepsilon e_n) dy - f(x) \int_{B_\varepsilon(x)} \eta_\varepsilon(x-y) dy \right| \\
		&\le \int_{B_\varepsilon(x)} \eta_\varepsilon(x-y) |f(y + 2\varepsilon e_n) - f(x)| dy \quad (\ast).
	\end{align*}
	Visto que $f$ é contínua, ela é uniformemente contínua em compactos.
	
	Para $\varepsilon$ suficientemente pequeno, $|f(y + 2\varepsilon e_n) - f(x)| < \varepsilon$ para todo $y \in B_\varepsilon(x)$. Esse comportamento de convergência uniforme em partes compactas para funções contínuas é análogo ao resultado demonstrado na \textbf{Questão 4}. De $(\ast)$, obtemos que: $$|f_\varepsilon(x) - f(x)| \le \int_{B_\varepsilon(x)} \eta_\varepsilon(x-y) |f(y + 2\varepsilon e_n) - f(x)| dy < \varepsilon \cdot \int_{B_\varepsilon(x)} \eta_\varepsilon(x-y) dy=\varepsilon\cdot 1=\varepsilon.$$
	
	Logo, $|f_\varepsilon(x) - f(x)| \to 0$ uniformemente em $K$.
	
	\textbf{(c)} Pela propriedade de derivação de convoluções, se $f \in C^m(\mathbb{R}^n_+)$, então $\tau_{2\varepsilon e_n} f \in C^m(\mathbb{R}^n_+)$. A derivada de ordem $\alpha$ de $f_\varepsilon = \eta_\varepsilon * \tilde{f}$ é dada por $D^\alpha f_\varepsilon = \eta_\varepsilon * (D^\alpha \tilde{f})$, resultado análogo ao obtido no item (d) da \textbf{Questão 4} para sequências regularizantes em $\mathbb{R}^n$. Substituindo a definição de $\tilde{f}$, obtemos:
	$$\boxed{D^\alpha f_\varepsilon(x) = \int_{\mathbb{R}^n_+} \eta_\varepsilon(x-y) (\tau_{2\varepsilon e_n} D^\alpha f)(y) dy}$$
	
	\textbf{(d)} Para $f \in L^p(\mathbb{R}^n_+)$, analisamos primeiramente a norma de $f_\varepsilon$. Pela desigualdade de Young para convoluções, temos que:
	$$\|f_\varepsilon\|_{L^p(\mathbb{R}^n_+)} = \|\eta_\varepsilon * (\tau_{2\varepsilon e_n} f)\|_{L^p} \le \|\eta_\varepsilon\|_{L^1} \|\tau_{2\varepsilon e_n} f\|_{L^p(\mathbb{R}^n_+)}$$
	Visto que $\|\eta_\varepsilon\|_{L^1} = 1$, resta calcular a norma da função transladada. Definimos a integral:
	$$\|\tau_{2\varepsilon e_n} f\|_{L^p(\mathbb{R}^n_+)}^p = \int_{\mathbb{R}^n_+} |f(x + 2\varepsilon e_n)|^p dx$$
	Efetuamos a mudança de variáveis $y = x + 2\varepsilon e_n$, com $dy = dx$. Analisamos a transformação do domínio de integração: a condição $x \in \mathbb{R}^n_+$ implica $x_n > 0$. Substituindo $x_n = y_n - 2\varepsilon$, obtemos a condição $y_n - 2\varepsilon > 0$, ou seja, $y_n > 2\varepsilon$. Denotando por $\mathbb{R}^n_{2\varepsilon} = \{y \in \mathbb{R}^n : y_n > 2\varepsilon\}$, temos:
	$$\|\tau_{2\varepsilon e_n} f\|_{L^p(\mathbb{R}^n_+)}^p = \int_{\mathbb{R}^n_{2\varepsilon}} |f(y)|^p dy$$
	Visto que $\mathbb{R}^n_{2\varepsilon} \subsetneq \mathbb{R}^n_+$, a integral sobre o subconjunto é limitada pela integral sobre o domínio total:
	\begin{align*}
		\int_{\mathbb{R}^n_{2\varepsilon}} |f(y)|^p dy &\le \int_{\mathbb{R}^n_+} |f(y)|^p dy \\
		\|\tau_{2\varepsilon e_n} f\|_{L^p(\mathbb{R}^n_+)}^p &\le \|f\|_{L^p(\mathbb{R}^n_+)}^p
	\end{align*}
	Portanto, $\|f_\varepsilon\|_{L^p(\mathbb{R}^n_+)} \le \|f\|_{L^p(\mathbb{R}^n_+)}$, o que prova que $f_\varepsilon \in L^p(\mathbb{R}^n_+)$.
	
	Para a convergência $\|f_\varepsilon - f\|_{L^p} \to 0$, utilizamos a desigualdade triangular:
	$$\|f_\varepsilon - f\|_{L^p(\mathbb{R}^n_+)} \le \|\eta_\varepsilon * (\tau_{2\varepsilon e_n} f) - \tau_{2\varepsilon e_n} f\|_{L^p(\mathbb{R}^n_+)} + \|\tau_{2\varepsilon e_n} f - f\|_{L^p(\mathbb{R}^n_+)}$$
	Analisamos os dois termos separadamente:
	\begin{enumerate}
		\item O primeiro termo $\lim_{\varepsilon \to 0} \|\eta_\varepsilon * (\tau_{2\varepsilon e_n} f) - \tau_{2\varepsilon e_n} f\|_{L^p} = 0$ decorre da propriedade de \textbf{identidades aproximadas}, visto que $\tau_{2\varepsilon e_n} f \in L^p(\mathbb{R}^n_+)$.
		\item O segundo termo $\lim_{\varepsilon \to 0} \|\tau_{2\varepsilon e_n} f - f\|_{L^p(\mathbb{R}^n_+)} = 0$ é garantido pela \textbf{continuidade da translação em $L^p$}, provada anteriormente na Questão 93 e 95.
	\end{enumerate}
	Dessa forma, a soma dos limites é nula, e concluímos que:
	$$\boxed{\|f_\varepsilon - f\|_{L^p(\mathbb{R}^n_+)} \to 0, \quad \text{quando } \varepsilon \to 0.}$$
	
	\textbf{(e)} Seja $u \in W^{1,p}(\mathbb{R}^n_+)$. Do item (d), sabemos que $u_\varepsilon \to u$ em $L^p(\mathbb{R}^n_+)$.
	Para as derivadas, aplicamos o item (c) com $|\alpha|=1$:
	$$\frac{\partial u_\varepsilon}{\partial x_j} = \eta_\varepsilon * \left( \tau_{2\varepsilon e_n} \frac{\partial u}{\partial x_j} \right)$$
	Como $\frac{\partial u}{\partial x_j} \in L^p(\mathbb{R}^n_+)$, aplicamos novamente o resultado do item (d) para a função $g = \frac{\partial u}{\partial x_j}$:
	$$\left\| \frac{\partial u_\varepsilon}{\partial x_j} - \frac{\partial u}{\partial x_j} \right\|_{L^p(\mathbb{R}^n_+)} \to 0, \quad \text{quando } \varepsilon \to 0$$
	Somando as convergências da função e de suas derivadas, concluímos que:
	$$\boxed{\|u_\varepsilon - u\|_{W^{1,p}(\mathbb{R}^n_+)} \to 0}.$$
}

\sintese{
	A construção de $f_\varepsilon$ resolve o problema da ``perda de informação'' na borda do semi-espaço. Ao transladar a função $f$ para o interior via $\tau_{2\varepsilon e_n}$, garantimos que o núcleo de convolução $\eta_\varepsilon$ opere sempre sobre valores definidos de $f$. O fluxo lógico foi: 
	$$\text{Translação} \to \text{Convolução} \to \text{Identidade Aproximada} \to \text{Convergência em } L^p \text{ e } W^{1,p}.$$
}

\aprofundamento{
	Esta técnica é um caso particular do teorema de \textbf{Meyers-Serrin}, que afirma que $C^\infty(\Omega) \cap W^{k,p}(\Omega)$ é denso em $W^{k,p}(\Omega)$ para qualquer aberto $\Omega$. A translação controlada $\tau_{\lambda e_n}$ é a ferramenta fundamental para provar esse resultado em domínios com fronteira regular, pois permite a regularização local sem a necessidade de extensões globais complexas, preservando a norma de Sobolev através da continuidade da translação.
}

%%%% QUESTÃO 97 %%%%
\questao{Mostre que cada função $u$ de $W^{1,\infty}(I)$ são de Lipschitz com constante $\|u'\|_{L^\infty(I)}$.}

\analise{
	\textbf{Análise do Enunciado:}
	A questão solicita a demonstração de que qualquer função pertencente ao espaço de Sobolev $W^{1,\infty}(I)$, onde $I \subset \mathbb{R}$ é um intervalo, admite um representante que é Lipschitz contínuo. A tese central é que a constante de Lipschitz $L$ é precisamente a norma $L^\infty$ da derivada fraca $u'$.
	
	\textbf{Fundamentação Teórica:}
	A prova fundamenta-se na relação entre a derivada fraca e a continuidade absoluta. De acordo com a teoria de espaços de Sobolev em dimensão 1, toda função $u \in W^{1,p}(I)$ possui um representante $\tilde{u}$ que é absolutamente contínua em $I$. Para tal representante, o Teorema Fundamental do Cálculo é válido, permitindo expressar a diferença $u(x) - u(y)$ como a integral de sua derivada fraca.
	
	\textbf{Estratégias:}
	\begin{enumerate}
		\item \textbf{Representação Integral}: Utilizar a propriedade de que funções em $W^{1, \infty}(I)$ são absolutamente contínuas, escrevendo a variação da função como a integral da derivada fraca sobre o intervalo $[y, x]$.
		\item \textbf{Estimativa de Norma}: Aplicar a definição de norma $L^\infty$ para limitar a função integranda por $\|u'\|_{L^\infty(I)}$.
		\item \textbf{Dedução da Desigualdade}: Integrar a constante resultante para obter a forma clássica da condição de Lipschitz $|u(x) - u(y)| \le L|x-y|$.
	\end{enumerate}
}

\solucao{
	Seja $u \in W^{1,\infty}(I)$. Pela teoria de imersões de Sobolev em dimensão $n=1$, sabemos que $u$ possui um representante $\tilde{u}$ que é absolutamente contínua em $I$. Para qualquer par de pontos $x, y \in I$ com $y < x$, a representação integral da função é dada por:
	$$u(x) - u(y) = \int_y^x u'(t) dt$$
	onde $u'$ denota a derivada fraca de $u$. Tomando o módulo em ambos os lados da igualdade, obtemos:
	$$|u(x) - u(y)| = \left| \int_y^x u'(t) dt \right|$$
	Pela propriedade da integral de Lebesgue, o módulo da integral é menor ou igual à integral do módulo:
	$$|u(x) - u(y)| \le \int_y^x |u'(t)| dt$$
	Visto que $u' \in L^\infty(I)$, por definição, temos que $|u'(t)| \le \|u'\|_{L^\infty(I)}$ para quase todo $t \in I$. Substituindo esse limitante superior na integral, obtemos:
	\begin{align*}
		|u(x) - u(y)| &\le \int_y^x \|u'\|_{L^\infty(I)} dt \\
		&= \|u'\|_{L^\infty(I)} \int_y^x 1 dt \\
		&= \|u'\|_{L^\infty(I)} (x - y)
	\end{align*}
	Como assumimos $y < x$, temos $(x - y) = |x - y|$. Para o caso geral $x, y \in I$, a desigualdade assume a forma:
	$$\boxed{|u(x) - u(y)| \le \|u'\|_{L^\infty(I)} |x - y|}$$
	Esta é precisamente a definição de uma função de Lipschitz com constante $L = \|u'\|_{L^\infty(I)}$.
}

\sintese{
	A demonstração revela que a norma $L^\infty$ da derivada fraca controla a taxa de variação máxima da função. O fluxo lógico foi:
	$$\text{Pertinência a } W^{1,\infty} \to \text{Absoluta Continuidade} \to \text{Representação Integral} \to \text{Limitante } L^\infty \to \text{Condição de Lipschitz}.$$
}

\aprofundamento{
	Este resultado estabelece a equivalência isométrica entre o espaço de Sobolev $W^{1,\infty}(I)$ e o espaço das funções Lipschitz $\text{Lip}(I)$. Inversamente, o \textbf{Teorema de Rademacher} garante que qualquer função de Lipschitz em um aberto de $\mathbb{R}^n$ é diferenciável quase sempre, e que sua derivada clássica coincide com a derivada fraca, pertencendo a $L^\infty$.
	
	A importância deste resultado reside no fato de que, enquanto em $W^{1,p}$ com $p < \infty$ a regularidade é medida por integrais de potência, no caso $p = \infty$ a regularidade torna-se geométrica (controle pontual da inclinação). Isso torna $W^{1,\infty}$ o espaço natural para estudar problemas de controle ótimo e equações diferenciais com coeficientes descontínuos, mas limitados.
}

%%%% QUESTÃO 98 %%%%
\questao{Suponha que $u$ é limitada e também uma função de Lipschitz em $I = (a, b)$, ou seja, $|u(x) - u(y)| \le c|x - y|$, para todo $x, y \in I$. Suponha que $I$ é um intervalo limitado e mostre que:
	\begin{enumerate}[label=(\alph*)]
		\item Para cada $\phi \in C_c^\infty(I)$ e $|h| < \min\{b - b_1, a_1 - a\}$ com $\supp \phi \subset [a_1, b_1] \subset I$, vale:
		$$\left| \int_a^b u(t)[\phi(t+h) - \phi(t)] dt \right| \le c|h| \int_I |\phi|;$$
		\item Consequentemente, para cada $1 < p < \infty$, $u \in L^p(I)$ e
		$$\left| \int_I u \phi' \right| \le c |I|^{1/p} \|\phi\|_{L^{p'}(I)}, \quad \forall \phi \in C_c^\infty(I);$$
		\item Conclua que $u \in W^{1,p}(I)$ e $\|u'\|_{L^p(I)} \le c|I|^{1/p}$;
		\item Conclua que $u \in W^{1,\infty}(I)$ e $\|u'\|_{L^\infty(I)} \le c$.
	\end{enumerate}}

\analise{
	\textbf{Análise do Enunciado:}
	O problema propõe a prova da recíproca da Questão 97: a demonstração de que a regularidade geométrica de Lipschitz implica a regularidade funcional de Sobolev $W^{1, \infty}$. O roteiro da questão é sequencial, onde cada item fornece a base para o seguinte. O ponto crítico é a transição de uma estimativa de diferenças finitas para a existência de um elemento no espaço dual $L^{p'}(I)$, o que garante a existência da derivada fraca.
	
	\textbf{Fundamentação Teórica:}
	A prova utiliza a definição de \textbf{Derivada Fraca} e a teoria de dualidade de espaços $L^p$. A passagem do item (b) para o (c) fundamenta-se no \textbf{Teorema de Representação de Riesz}, que assegura que todo funcional linear limitado em $L^{p'}(I)$ é representado por um elemento de $L^p(I)$. Vale notar que a extensão deste funcional de $C_c^\infty(I)$ para $L^{p'}(I)$ é garantida pelo \textbf{Teorema de Hahn-Banach}, dada a densidade de $C_c^\infty$ em $L^{p'}$.
	
	\textbf{Estratégias:}
	\begin{enumerate}
		\item \textbf{Mudança de Variáveis}: No item (a), transferiremos a translação do argumento de $\phi$ para $u$, permitindo a aplicação direta da constante de Lipschitz $c$.
		\item \textbf{Passagem ao Limite}: No item (b), recuperaremos a derivada $\phi'$ via limite de diferenças finitas, utilizando a Desigualdade de Hölder para introduzir a norma $L^{p'}$.
		\item \textbf{Identificação do Dual}: No item (c), definiremos um funcional linear limitado e utilizaremos a representação de Riesz para identificar a derivada fraca $u' \in L^p(I)$.
		\item \textbf{Limite de Normas}: No item (d), utilizaremos a propriedade de que a norma $L^\infty$ é o limite das normas $L^p$ quando $p \to \infty$.
	\end{enumerate}
}

\solucao{
	\textbf{(a)} Seja $\phi \in C_c^\infty(I)$ com $\text{supp } \phi \subset [a_1, b_1] \subset I$. Consideramos a integral:
	$$J = \int_a^b u(t)[\phi(t+h) - \phi(t)] dt = \int_a^b u(t) \phi(t+h) dt - \int_a^b u(t) \phi(t) dt$$
	Na primeira integral, efetuamos a mudança de variáveis $s = t+h \implies t = s-h$. Como $\text{supp } \phi \subset [a_1, b_1]$ e $|h|$ é suficientemente pequeno, a função $\phi$ permanece nula fora de $I$, permitindo que a integral seja escrita como:
	$$\int_a^b u(t) \phi(t+h) dt = \int_a^b u(s-h) \phi(s) ds$$
	Substituindo de volta em $J$:
	\begin{align*}
		J &= \int_a^b [u(s-h) - u(s)] \phi(s) ds
	\end{align*}
	Aplicando a hipótese de que $u$ é Lipschitz com constante $c$, temos $|u(s-h) - u(s)| \le c|s-h-s| = c|h|$. Assim:
	$$|J| \le \int_a^b |u(s-h) - u(s)| |\phi(s)| ds \le \int_a^b c|h| |\phi(s)| ds = c|h| \int_I |\phi|$$
	Portanto, $\boxed{\left| \int_a^b u(t)[\phi(t+h) - \phi(t)] dt \right| \le c|h| \int_I |\phi|}$.
	
	\textbf{(b)} Dividindo a desigualdade do item (a) por $h$ e fazendo $h \to 0$, a diferença finita $\tfrac{\phi(t+h) - \phi(t)}{h}$ converge uniformemente para $\phi'(t)$ em $I$. Pelo Teorema da Convergência Dominada, temos:
	$$\left| \int_I u \phi' dt \right| = \lim_{h \to 0} \left| \int_a^b u(t) \frac{\phi(t+h) - \phi(t)}{h} dt \right| \le c \int_I |\phi| dt$$
	Visto que $u$ é limitada e $I$ possui medida finita, a inclusão $L^\infty(I) \subset L^p(I)$ é imediata, logo $u \in L^p(I)$. Para $\phi \in C_c^\infty(I)$, aplicamos a \textbf{Desigualdade de Hölder}:
	$$\int_I |\phi| dt \le \|\phi\|_{L^{p'}(I)} \|1\|_{L^p(I)} = \|\phi\|_{L^{p'}(I)} |I|^{1/p}$$
	onde $\tfrac{1}{p} + \tfrac{1}{p'} = 1$. Substituindo esta estimativa, obtemos:
	$$\boxed{\left| \int_I u \phi' \right| \le c |I|^{1/p} \|\phi\|_{L^{p'}(I)}}$$
	
	\textbf{(c)} Definimos o funcional linear $L: C_c^\infty(I) \to \mathbb{R}$ por $L(\phi) = - \int_I u \phi' dt$. O resultado do item (b) mostra que $|L(\phi)| \le (c |I|^{1/p}) \|\phi\|_{L^{p'}(I)}$, logo $L$ é um funcional linear limitado no subespaço denso $C_c^\infty(I) \subset L^{p'}(I)$.
	
	Pelo \textbf{Teorema de Hahn-Banach}, $L$ estende-se unicamente a um funcional linear limitado em todo $L^{p'}(I)$. Pelo \textbf{Teorema de Representação de Riesz}, existe um único elemento $g \in L^p(I)$ tal que:
	$$L(\phi) = \int_I g \phi dt \implies - \int_I u \phi' dt = \int_I g \phi dt, \quad \forall \phi \in C_c^\infty(I)$$
	Esta é precisamente a definição de derivada fraca, identificando $g = u'$. Como $u \in L^p(I)$ e $u' = g \in L^p(I)$, concluímos que $u \in W^{1,p}(I)$. Além disso, a norma do elemento representativo coincide com a norma do funcional:
	$$\boxed{\|u'\|_{L^p(I)} = \|L\| \le c |I|^{1/p}}$$
	
	\textbf{(d)} Para concluir a pertinência a $W^{1, \infty}(I)$, analisamos o comportamento da norma de $u'$ quando $p \to \infty$. Sabemos que, para qualquer $f \in L^\infty(I)$, $\|f\|_{L^\infty} = \lim_{p \to \infty} \|f\|_{L^p}$. Aplicando isso a $u'$:
	$$\|u'\|_{L^\infty(I)} = \lim_{p \to \infty} \|u'\|_{L^p(I)} \le \lim_{p \to \infty} c |I|^{1/p}$$
	Como $|I|^{1/p} \to 1$ quando $p \to \infty$, obtemos:
	$$\|u'\|_{L^\infty(I)} \le c$$
	Sendo $u$ limitada e possuindo derivada fraca em $L^\infty(I)$, concluímos que $\boxed{u \in W^{1,\infty}(I)}$ com $\|u'\|_{L^\infty(I)} \le c$.
}

\sintese{
	A prova estabelece a equivalência entre a regularidade geométrica (Lipschitz) e a regularidade funcional (Sobolev $W^{1, \infty}$). O fluxo lógico foi:
	$$\text{Lipschitz} \to \text{Funcional Limitado em } L^{p'} \to \text{Existência de } u' \in L^p \to \text{Limite } p \to \infty \to W^{1, \infty}.$$
}

\aprofundamento{
	Este resultado é a implementação do \textbf{Teorema de Rademacher} em dimensão 1. A transição do item (b) para o (c) é o ponto mais profundo da prova, pois transforma uma estimativa de integral em uma afirmação sobre a existência de um objeto matemático (a derivada fraca).
	
	A utilização da norma de $L^p$ como ponte para chegar a $L^\infty$ demonstra que a constante de Lipschitz $c$ não é apenas um limite superior para a variação da função, mas é precisamente o supremo essencial da sua derivada fraca. Isso solidifica a ideia de que $W^{1, \infty}(I)$ é o espaço natural para funções cujas taxas de variação são globalmente limitadas.
}

%%%% QUESTÃO 99 %%%%
\questao{Mostre que $C_c^\infty(\mathbb{R}^n)$ é denso em $W^{1,p}(\mathbb{R}^n)$. Conclua que $W^{1,p}(\mathbb{R}^n) = W_0^{1,p}(\mathbb{R}^n)$.}

\analise{
	\textbf{Análise do Enunciado:}
	A questão exige a prova de que qualquer função $u \in W^{1,p}(\mathbb{R}^n)$ pode ser aproximada, na norma de Sobolev, por uma sequência de funções suaves com suporte compacto. A tese final é a identidade entre o espaço $W^{1,p}(\mathbb{R}^n)$ e o seu subespaço $W_0^{1,p}(\mathbb{R}^n)$, que é definido como o fecho de $C_c^\infty(\mathbb{R}^n)$.
	
	\textbf{Fundamentação Teórica:}
	A prova utiliza a técnica de truncação via funções de corte (\textit{cutoff functions}) e a regularização por convolução. Para funções definidas em todo o espaço $\mathbb{R}^n$, a convolução simples $u * \eta_\varepsilon$ converge para $u$ na norma $W^{1,p}$, mas não preserva a compacidade do suporte. Para resolver isso, multiplicamos $u$ por uma função $\rho_R$ que anula a função fora de uma bola de raio $2R$, permitindo a subsequente suavização.
	
	\textbf{Estratégias:}
	\begin{enumerate}
		\item \textbf{Aproximação por Suporte Compacto}: Definir $\rho \in C_c^\infty(\mathbb{R}^n)$ e a sequência $u_R = u \rho_R$. Provar que $u_R \to u$ em $W^{1,p}(\mathbb{R}^n)$ quando $R \to \infty$, detalhando a convergência da integral da cauda e o decaimento do termo do gradiente.
		\item \textbf{Aproximação por Suavização}: Para cada $u_R$ fixo, aplicar a convolução $v_{R,\varepsilon} = u_R * \eta_\varepsilon$. Provar que $v_{R,\varepsilon} \to u_R$ em $W^{1,p}(\mathbb{R}^n)$ quando $\varepsilon \to 0$, utilizando as propriedades de identidades aproximadas.
		\item \textbf{Argumento Diagonal}: Combinar os dois processos para mostrar que $\forall \varepsilon > 0$, existe $v \in C_c^\infty(\mathbb{R}^n)$ tal que $\|u - v\|_{W^{1,p}(\mathbb{R}^n)} < \varepsilon$.
		\item \textbf{Identificação de Espaços}: Aplicar a definição de $W_0^{1,p}(\Omega)$ para concluir a igualdade dos espaços.
	\end{enumerate}
}

\solucao{
	Seja $u \in W^{1,p}(\mathbb{R}^n)$. Construiremos a aproximação em duas etapas: primeiro restringindo o suporte e, depois, suavizando a função.
	
	\textit{Passo 1: Truncação via função de corte}\par
	Seja $\rho \in C_c^\infty(\mathbb{R}^n)$ uma função tal que $0 \le \rho \le 1$, $\rho(x) = 1$ para $|x| \le 1$ e $\rho(x) = 0$ para $|x| \ge 2$. Para $R > 0$, definimos a função de corte $\rho_R(x) = \rho(x/R)$. Note que $\rho_R \in C_c^\infty(\mathbb{R}^n)$ com $\text{supp}(\rho_R) \subset \overline{B_{2R}(0)}$. Definimos a sequência $u_R = u \rho_R$.
	
	Analisamos a convergência de $u_R$ para $u$ em $L^p(\mathbb{R}^n)$:
	$$\|u - u_R\|_{L^p(\mathbb{R}^n)}^p = \int_{\mathbb{R}^n} |u(x)(1 - \rho_R(x))|^p dx = \int_{\mathbb{R}^n \setminus B_R(0)} |u(x)|^p |1 - \rho_R(x)|^p dx$$
	Visto que $0 \le 1 - \rho_R \le 1$, temos:
	$$\|u - u_R\|_{L^p(\mathbb{R}^n)}^p \le \int_{|x| > R} |u(x)|^p dx$$
	Como $u \in L^p(\mathbb{R}^n)$, a integral $\int_{\mathbb{R}^n} |u|^p dx$ é finita. Pela propriedade de continuidade da medida de Lebesgue (ou pelo Teorema da Convergência Dominada), a integral sobre a região $|x| > R$ tende a zero quando $R \to \infty$. Logo, $\|u - u_R\|_{L^p(\mathbb{R}^n)} \to 0$.
	
	Para as derivadas, usamos a regra do produto: $\nabla u_R = \rho_R \nabla u + u \nabla \rho_R$. A diferença é:
	$$\|\nabla u - \nabla u_R\|_{L^p(\mathbb{R}^n)} = \|\nabla u(1 - \rho_R) - u \nabla \rho_R\|_{L^p(\mathbb{R}^n)} \le \|\nabla u(1 - \rho_R)\|_{L^p(\mathbb{R}^n)} + \|u \nabla \rho_R\|_{L^p(\mathbb{R}^n)}$$
	O primeiro termo $\int_{|x| > R} |\nabla u|^p dx$ tende a zero por razões análogas ao caso de $u$. Para o segundo termo, observe que por regra da cadeia:
	$$\nabla \rho_R(x) = \frac{1}{R} \nabla \rho\left(\frac{x}{R}\right)$$
	Definimos $C = \|\nabla \rho\|_{L^\infty(\mathbb{R}^n)}$, que é uma constante finita pois $\rho$ é suave com suporte compacto. Assim, $|\nabla \rho_R(x)| \le \frac{C}{R}$. Então:
	$$\|u \nabla \rho_R\|_{L^p(\mathbb{R}^n)}^p = \int_{R < |x| < 2R} |u(x)|^p |\nabla \rho_R(x)|^p dx \le \frac{C^p}{R^p} \int_{R < |x| < 2R} |u(x)|^p dx \le \frac{C^p}{R^p} \|u\|_{L^p(\mathbb{R}^n)}^p$$
	Como $p \ge 1$ e $R \to \infty$, o termo $\frac{C^p}{R^p} \to 0$. Concluímos que $u_R \to u$ na norma $\|\cdot\|_{W^{1,p}(\mathbb{R}^n)}$.
	
	\textit{Passo 2: Regularização por Convolução}\par
	Para cada $R$ fixo, $u_R$ possui suporte compacto em $B_{2R}(0)$. Definimos $v_{R,\varepsilon} = u_R * \eta_\varepsilon$, onde $\eta_\varepsilon$ é a função regularizante da Questão 4. Visto que $u_R \in W^{1,p}(\mathbb{R}^n)$ com suporte compacto, a função $v_{R,\varepsilon}$ é de classe $C^\infty(\mathbb{R}^n)$ e possui suporte compacto em $B_{2R+\varepsilon}(0)$.
	
	Pelas propriedades de identidades aproximadas (Questão 4), sabemos que:
	$$\|u_R - v_{R,\varepsilon}\|_{L^p(\mathbb{R}^n)} \to 0 \quad \text{e} \quad \|\nabla u_R - \nabla v_{R,\varepsilon}\|_{L^p(\mathbb{R}^n)} \to 0, \quad \text{quando } \varepsilon \to 0.$$
	Portanto, $\|u_R - v_{R,\varepsilon}\|_{W^{1,p}(\mathbb{R}^n)} \to 0$.
	
	\textit{Passo 3: Conclusão da Densidade}\par
	Dos passos 1 e 2, para qualquer $\delta > 0$, podemos escolher $R$ tal que $\|u - u_R\|_{W^{1,p}(\mathbb{R}^n)} < \delta/2$ e, para este $R$, escolhemos $\varepsilon$ tal que 
	$$\|u_R - v_{R,\varepsilon}\|_{W^{1,p}(\mathbb{R}^n)} < \delta/2.$$
	Pela desigualdade triangular:
	$$\|u - v_{R,\varepsilon}\|_{W^{1,p}(\mathbb{R}^n)} \le \|u - u_R\|_{W^{1,p}(\mathbb{R}^n)} + \|u_R - v_{R,\varepsilon}\|_{W^{1,p}(\mathbb{R}^n)} < \delta$$
	Como $v_{R,\varepsilon} \in C_c^\infty(\mathbb{R}^n)$, provamos que $C_c^\infty(\mathbb{R}^n)$ é denso em $W^{1,p}(\mathbb{R}^n)$.
	
	\textit{Passo 4: Igualdade dos Espaços}\par
	O espaço $W_0^{1,p}(\mathbb{R}^n)$ é definido como o fecho de $C_c^\infty(\mathbb{R}^n)$ na norma de $W^{1,p}(\mathbb{R}^n)$. Como acabamos de mostrar que qualquer $u \in W^{1,p}(\mathbb{R}^n)$ é o limite de uma sequência em $C_c^\infty(\mathbb{R}^n)$, segue que:
	$$\boxed{W^{1,p}(\mathbb{R}^n) = W_0^{1,p}(\mathbb{R}^n)}$$
}

\sintese{
	A prova demonstra que a ausência de uma fronteira finita em $\mathbb{R}^n$ remove a distinção entre funções que "desaparecem na borda" e funções geralmente integráveis. O fluxo lógico foi:
	$$\text{Função em } W^{1,p} \xrightarrow{\text{Truncação } \rho_R} \text{Suporte Compacto} \xrightarrow{\text{Convolução } \eta_\varepsilon} \text{Suavidade} \to \text{Densidade} \to \text{Igualdade } W_0 = W.$$
}

\aprofundamento{
	Este resultado evidencia a importância da geometria do domínio. Em um domínio limitado $\Omega$, $W_0^{1,p}(\Omega)$ é um subespaço estrito de $W^{1,p}(\Omega)$, pois as funções em $W_0^{1,p}$ devem possuir traço nulo em $\partial \Omega$. 
	
	No caso de $\mathbb{R}^n$, a \textbf{fronteira} situa-se no infinito. Como a pertinência a $L^p(\mathbb{R}^n)$ força a função a decair no infinito, a condição de traço nulo é satisfeita intrinsecamente. Este resultado justifica a utilização de funções de teste com suporte compacto para a definição de distribuições e a validade de integrações por partes em todo o espaço sem a necessidade de termos de borda.
}

%%%% QUESTÃO 100 %%%%
\questao{Considere a função $\sigma : [0, \infty) \to \mathbb{R}$ dada por $\sigma(r) = 0$ em $r \ge 2$, $\sigma(r) = 1$ em $r \le 1$ e $\sigma(r) = 2 - r$ se $1 \le r \le 2$. Mostre que:
	\begin{enumerate}[label=(\alph*)]
		\item $u(x) = \sigma(|x|)$ tem derivadas fracas dadas por $\nabla u(x) = \sigma'(|x|) \frac{x}{|x|}$ e $|\nabla u| \le 1, \forall x \in \mathbb{R}^n$ q.s.;
		\item $0 \le u \le 1$, $\text{supp } u = \overline{B_2(0)}$;
		\item $u \in W^{1,p}(\mathbb{R}^n)$;
		\item Para cada $\varepsilon > 0$, mostre que $u_\varepsilon(x) = \sigma(\frac{x}{\varepsilon})$ está em $W^{1,p}(\mathbb{R}^n)$ e
		$$\|\nabla u_\varepsilon\|_{L^p(\mathbb{R}^n)}^p \le \frac{|B_{2\varepsilon}(0) \setminus B_\varepsilon(0)|}{\varepsilon^p} = n\alpha(n)(2^n - 1)\varepsilon^{n-p}.$$
		\item Se $n > p$, então $\|\nabla u_\varepsilon\|_{L^p(\mathbb{R}^n)} \to 0$, quando $\varepsilon \to 0$, e consequentemente, $u_\varepsilon \to 0$ em $W^{1,p}(\mathbb{R}^n)$.
\end{enumerate}}

\analise{
	\textbf{Análise do Enunciado:}
	A questão propõe a análise de uma função de truncação radial. O objetivo é verificar sua pertinência aos espaços de Sobolev e analisar como a norma do gradiente escala sob contrações homotéticas. A tese central é que, para $n > p$, a energia do gradiente converge a zero mesmo que a amplitude da função permaneça constante, o que demonstra a falha da imersão de Sobolev em $L^\infty$ para $p \le n$.
	
	\textbf{Fundamentação Teórica:}
	A prova baseia-se na composição de funções Lipschitzianas. Utilizaremos o resultado da \textbf{Questão 98} para garantir que a regularidade de Lipschitz implica a existência de derivadas fracas em $L^\infty(\mathbb{R}^n)$. Para a forma do gradiente, utilizaremos a fundamentação da \textbf{Questão 91}, que provou que a derivada fraca da norma $|x|$ é $x/|x|$, tratando a singularidade na origem. Para as integrais, utilizaremos a notação $\alpha(n)$ denota o volume da bola unitária em $\mathbb{R}^n$.
	
	\textbf{Estratégias:}
	\begin{enumerate}
		\item \textbf{Composição Lipschitz}: Demonstrar que $u$ é Lipschitz por ser a composição de $\sigma$ e $|\cdot|$, invocando a Questão 98 para assegurar que $u \in W^{1, \infty}_{loc}(\mathbb{R}^n)$.
		\item \textbf{Gradiente Radial}: Aplicar a regra da cadeia para funções Lipschitz, utilizando o resultado da Questão 91 para validar a derivada da norma.
		\item \textbf{Pertinência a $L^p$}: Justificar a pertinência a $W^{1,p}$ via limitação da função e do gradiente em um domínio de suporte compacto.
		\item \textbf{Cálculo de Escala}: Efetuar a mudança de variáveis para $u_\varepsilon$, calculando a norma do gradiente no anel $B_{2\varepsilon} \setminus B_\varepsilon$ e expressando o volume via $\alpha(n)$.
		\item \textbf{Análise do Expoente}: Mostrar que a convergência para zero depende estritamente da condição $n > p$.
	\end{enumerate}
}

\solucao{
	\textbf{(a)} A função $\sigma(r)$ é Lipschitziana em $[0, \infty)$ com constante $1$, possuindo derivada clássica $\sigma'(r)$ dada quase sempre por:
	$$\sigma'(r) = \begin{cases} 0, & \text{se } r < 1 \\ -1, & \text{se } 1 < r < 2 \\ 0, & \text{se } r > 2 \end{cases}$$
	Como a norma $|x|$ também é Lipschitziana em $\mathbb{R}^n$, a composição $u(x) = \sigma(|x|)$ é Lipschitz em $\mathbb{R}^n$. Pela generalização da \textbf{Questão 98}, concluímos que $u \in W^{1, \infty}_{loc}(\mathbb{R}^n)$, o que assegura a existência de suas derivadas fracas.
	
	Para o cálculo de $\nabla u$, utilizamos a regra da cadeia. Visto que a derivada fraca da norma é $\nabla |x| = \frac{x}{|x|}$ (conforme provado na \textbf{Questão 91}, tratando a singularidade na origem), temos:
	$$\nabla u(x) = \sigma'(|x|) \nabla (|x|) = \sigma'(|x|) \frac{x}{|x|}$$
	Considerando que $|\nabla u(x)| = |\sigma'(|x|)| \cdot | \frac{x}{|x|} | = |\sigma'(|x|)|$ e que $|\sigma'(r)| \le 1$ q.s., segue que $|\nabla u(x)| \le 1$ q.s. em $\mathbb{R}^n$.
	
	\textbf{(b)} Da definição de $\sigma$, temos $\sigma(r) = 1$ para $r \in [0,1]$ e $\sigma(r) = 0$ para $r \in [2, \infty)$. Como a transição é linear, $0 \le \sigma(r) \le 1$ para todo $r \ge 0$. Logo, $0 \le u(x) \le 1$. O suporte de $u$ é o fecho do conjunto onde $u \neq 0$:
	$$\supp u = \overline{\{x \in \mathbb{R}^n : |x| < 2\}} = \overline{B_2(0)}.$$
	
	\textbf{(c)} Para $u \in W^{1,p}(\mathbb{R}^n)$, devemos verificar a integrabilidade de $u$ e $\nabla u$.
	Visto que $u$ é limitada e possui suporte compacto em $B_2(0)$, temos:
	$$\|u\|_{L^p(\mathbb{R}^n)}^p = \int_{B_2(0)} |u(x)|^p dx \le \int_{B_2(0)} 1 dx = \alpha(n) 2^n < \infty$$
	Analogamente, para o gradiente, que é nulo fora do anel $B_2(0) \setminus B_1(0)$:
	$$\|\nabla u\|_{L^p(\mathbb{R}^n)}^p = \int_{B_2(0) \setminus B_1(0)} |-1|^p dx = \alpha(n)(2^n - 1) < \infty$$
	Logo, $u \in W^{1,p}(\mathbb{R}^n)$ para todo $1 \le p \le \infty$.
	
	\textbf{(d)} Consideremos a função escalonada $u_\varepsilon(x) = \sigma(|x|/\varepsilon)$. Visto que $u_\varepsilon$ é apenas uma dilatação de $u$, ela herda as propriedades de ser limitada e Lipschitz com suporte compacto em $\overline{B_{2\varepsilon}(0)}$, o que garante automaticamente que $u_\varepsilon \in W^{1,p}(\mathbb{R}^n)$. 
	
	Pela regra da cadeia, a derivada fraca é $\nabla u_\varepsilon(x) = \frac{1}{\varepsilon} \sigma'(\frac{|x|}{\varepsilon}) \frac{x}{|x|}$. Calculamos a norma integrando sobre o anel onde $\sigma' \neq 0$, isto é, $1 < \frac{|x|}{\varepsilon} < 2 \implies \varepsilon < |x| < 2\varepsilon$:
	\begin{align*}
		\|\nabla u_\varepsilon\|_{L^p(\mathbb{R}^n)}^p &= \int_{B_{2\varepsilon}(0) \setminus B_\varepsilon(0)} \left| \frac{1}{\varepsilon} \sigma'\left(\frac{|x|}{\varepsilon}\right) \right|^p dx \\
		&= \frac{1}{\varepsilon^p} \int_{B_{2\varepsilon}(0) \setminus B_\varepsilon(0)} 1 dx = \frac{\alpha(n)((2\varepsilon)^n - \varepsilon^n)}{\varepsilon^p} = \alpha(n)(2^n - 1)\varepsilon^{n-p}.
	\end{align*}
	\textit{Note que a expressão do enunciado $n\alpha(n)(2^n-1)\varepsilon^{n-p}$ sugere a utilização da área da superfície $\omega_n = n\alpha(n)$, porém a integração direta do volume resulta em $\alpha(n)(2^n - 1)\varepsilon^{n-p}$.}
	
	\textbf{(e)} Se $n > p$, o expoente $n-p$ é positivo, implicando que:
	$$\lim_{\varepsilon \to 0} \|\nabla u_\varepsilon\|_{L^p(\mathbb{R}^n)}^p = \lim_{\varepsilon \to 0} \alpha(n)(2^n - 1)\varepsilon^{n-p} = 0$$
	Para a função, $\|u_\varepsilon\|_{L^p(\mathbb{R}^n)}^p \le \int_{B_{2\varepsilon}(0)} 1 dx = \alpha(n) 2^n \varepsilon^n \to 0$.
	Sendo a norma de Sobolev a soma das normas de $L^p$ da função e do gradiente, concluímos que $\|u_\varepsilon\|_{W^{1,p}(\mathbb{R}^n)} \to 0$, ou seja, $\boxed{u_\varepsilon \to 0 \text{ em } W^{1,p}(\mathbb{R}^n)}$.
}

\sintese{
	A construção de $u_\varepsilon$ demonstra que, em dimensões superiores ($n > p$), a energia do gradiente pode ser dissipada por contração do suporte, mesmo mantendo a amplitude da função constante. O fluxo lógico foi: $$\text{Lipschitz} \xrightarrow{\text{Questão 98}} W^{1,\infty} \xrightarrow{\text{Questão 91}} \text{Gradiente Radial} \to \text{Escalonamento } \varepsilon \to \text{Análise de } n-p.$$
}

\aprofundamento{
	Este resultado é a prova construtiva da inexistência da imersão $W^{1,p}(\mathbb{R}^n) \hookrightarrow L^\infty(\mathbb{R}^n)$ para $p \le n$. Se tal imersão existisse, teríamos $\|u_\varepsilon\|_{L^\infty} \le C \|u_\varepsilon\|_{W^{1,p}}$. Contudo, $\|u_\varepsilon\|_{L^\infty} = 1$ enquanto $\|u_\varepsilon\|_{W^{1,p}} \to 0$, o que gera a contradição $1 \le 0$. Isso fundamenta a necessidade do Teorema de Morrey para garantir a continuidade pontual, exigindo $p > n$.
}

%%%% QUESTÃO 101 %%%%
\questao{Ainda com a notação do exercício anterior, mostre que $u \in W_0^{1,p}(\mathbb{R}^n \setminus \{0\})$. Para isto basta mostrar que $v_\varepsilon = u(1 - u_\varepsilon) \in W_0^{1,p}(\mathbb{R}^n \setminus \{0\})$ e $v_\varepsilon \to u$ em $W^{1,p}(\mathbb{R}^n \setminus \{0\})$.}

\analise{
	\textbf{Análise do Enunciado:}
	A questão solicita a prova de que a função $u$ (definida na \textbf{Questão 100}) pertence ao espaço $W_0^{1,p}$ do domínio perfurado $\Omega = \mathbb{R}^n \setminus \{0\}$. A diferença crucial aqui é que, para pertencer a $W_0^{1,p}(\Omega)$, a função deve ser aproximável por funções de $C_c^\infty(\Omega)$, o que implica que a aproximação deve ter suporte compacto que \textbf{não contenha a origem}. A estratégia sugerida utiliza a função $v_\varepsilon$ para "escavar" um buraco no suporte de $u$ em torno da origem.
	
	\textbf{Fundamentação Teórica:}
	A prova repousa sobre a definição de $W_0^{1,p}$ como o fecho de $C_c^\infty$ e a propriedade de que, se $u \in W^{1,p}(\mathbb{R}^n)$ e $u$ pode ser aproximada por funções com suporte afastado da origem, então a \textbf{capacidade} do ponto $\{0\}$ é nula. Utilizaremos a regra do produto para derivadas fracas (\textbf{Questão 89}) e o resultado de convergência de $u_\varepsilon$ obtido na \textbf{Questão 100}.
	
	\textbf{Estratégias}
	\begin{enumerate}
		\item \textbf{Análise do Suporte de $v_\varepsilon$}: Demonstrar que $v_\varepsilon$ é nula em uma vizinhança da origem ($|x| < \varepsilon$) e nula fora de uma bola grande ($|x| > 2$), garantindo que $\text{supp}(v_\varepsilon)$ é um compacto contido em $\mathbb{R}^n \setminus \{0\}$.
		\item \textbf{Pertinência a $W_0^{1,p}$}: Justificar que, sendo $v_\varepsilon$ uma função Lipschitz com suporte compacto em $\Omega$, ela pertence a $W_0^{1,p}(\Omega)$ via densidade de $C_c^\infty$.
		\item \textbf{Convergência em $L^p$}: Provar que $\|v_\varepsilon - u\|_{L^p} \to 0$ utilizando o fato de que $u_\varepsilon \to 0$ pontualmente e é limitada.
		\item \textbf{Convergência do Gradiente}: Decompor $\nabla v_\varepsilon - \nabla u$ e utilizar o resultado da Questão 100(e) para mostrar que o termo envolvendo $\nabla u_\varepsilon$ converge a zero, desde que $n > p$.
	\end{enumerate}
}

\solucao{
	Iniciamos analisando a função $v_\varepsilon(x) = u(x)(1 - u_\varepsilon(x))$ e seu suporte em $\mathbb{R}^n$. Se $|x| < \varepsilon$, temos $|x|/\varepsilon < 1$, o que implica $u_\varepsilon(x) = \sigma(|x|/\varepsilon) = 1$, resultando em $v_\varepsilon(x) = u(x)(1 - 1) = 0$. Por outro lado, se $|x| > 2$, a função $u$ anula-se ($\sigma(|x|) = 0$), logo $v_\varepsilon(x) = 0 \cdot (1 - u_\varepsilon(x)) = 0$. Concluímos que o suporte de $v_\varepsilon$ está contido no anel fechado $\overline{B_2(0)} \setminus B_\varepsilon(0)$, que é um compacto contido em $\mathbb{R}^n \setminus \{0\}$. Como $u$ e $u_\varepsilon$ são funções Lipschitz, $v_\varepsilon$ também é Lipschitz com suporte compacto em $\mathbb{R}^n \setminus \{0\}$, o que garante que $v_\varepsilon \in W_0^{1,p}(\mathbb{R}^n \setminus \{0\})$.
	
	Para provar a convergência de $v_\varepsilon$ para $u$ em $L^p(\mathbb{R}^n \setminus \{0\})$, observamos que a diferença é dada por 
	$$v_\varepsilon(x) - u(x) = -u(x)u_\varepsilon(x).$$
	Estimando a norma $L^p$, temos:
	$$\|v_\varepsilon - u\|_{L^p(\mathbb{R}^n \setminus \{0\})}^p = \int_{\mathbb{R}^n \setminus \{0\}} |u(x)u_\varepsilon(x)|^p dx \le \|u\|_{L^\infty(\mathbb{R}^n)}^p \int_{B_{2\varepsilon}(0)} |u_\varepsilon(x)|^p dx$$
	Como $|u_\varepsilon| \le 1$ e a medida da bola $|B_{2\varepsilon}(0)| = \alpha(n)(2\varepsilon)^n$, obtemos:
	$$\|v_\varepsilon - u\|_{L^p(\mathbb{R}^n \setminus \{0\})}^p \le \|u\|_{L^\infty(\mathbb{R}^n)}^p \cdot \alpha(n)(2\varepsilon)^n \to 0, \quad \text{quando } \varepsilon \to 0.$$
	
	Em seguida, analisamos a convergência do gradiente. Aplicando a regra do produto para derivadas fracas (estudada na \textbf{Questão 89}), o gradiente de $v_\varepsilon$ é $\nabla v_\varepsilon = \nabla u (1 - u_\varepsilon) - u \nabla u_\varepsilon$. A diferença em relação ao gradiente de $u$ é:
	$$\nabla v_\varepsilon - \nabla u = -u_\varepsilon \nabla u - u \nabla u_\varepsilon$$
	Para o primeiro termo, notamos que $\supp(u_\varepsilon) = \overline{B_{2\varepsilon}(0)}$, logo:
	$$\int_{\mathbb{R}^n \setminus \{0\}} |u_\varepsilon \nabla u|^p dx \le \int_{B_{2\varepsilon}(0)} |\nabla u|^p dx$$
	Como $\nabla u \in L^p(\mathbb{R}^n)$, a integral sobre a bola de raio $2\varepsilon$ tende a zero por continuidade da medida de Lebesgue quando $\varepsilon \to 0$. Para o segundo termo, utilizamos a limitação de $u$:
	$$\|u \nabla u_\varepsilon\|_{L^p(\mathbb{R}^n \setminus \{0\})} \le \|u\|_{L^\infty(\mathbb{R}^n)} \|\nabla u_\varepsilon\|_{L^p(\mathbb{R}^n)}$$
	Pelo resultado da \textbf{Questão 100(e)}, sob a condição $n > p$, temos que $\|\nabla u_\varepsilon\|_{L^p(\mathbb{R}^n)} \to 0$ quando $\varepsilon \to 0$. Portanto, $\|\nabla v_\varepsilon - \nabla u\|_{L^p(\mathbb{R}^n \setminus \{0\})} \to 0$.
	
	Concluímos que $\|v_\varepsilon - u\|_{W^{1,p}(\mathbb{R}^n \setminus \{0\})} \to 0$ e como $v_\varepsilon \in \overline{C_c^\infty(\mathbb{R}^n \setminus \{0\})} = W_0^{1,p}(\mathbb{R}^n \setminus \{0\})$, a natureza fechada do espaço $W_0^{1,p}$ na topologia de $W^{1,p}(\mathbb{R}^n \setminus \{0\})$ implica que:
	$$\boxed{u \in W_0^{1,p}(\mathbb{R}^n \setminus \{0\}).}$$
}

\sintese{
A prova demonstra que a remoção de um ponto não altera a natureza do espaço de Sobolev quando a dimensão $n$ é superior ao expoente $p$. O uso da função de corte $u_\varepsilon$ permitiu a construção de uma sequência de aproximações que anulam a função na vizinhança da origem, provando que a \textbf{singularidade} do ponto $\{0\}$ é irrelevante para a norma $W^{1,p}$ sob a condição $n > p$. O fluxo lógico foi: $$\text{Construção de } v_\varepsilon \to \text{Suporte em } \Omega \to \text{Convergência de } L^p \to \text{Convergência de } \nabla \to \text{Pertinência a } W_0^{1,p}.$$
}

\aprofundamento{
Este resultado é a base para a teoria de \textbf{Capacidade em Espaços de Sobolev}. Dizemos que o ponto $\{0\}$ tem $p$-capacidade nula se a condição $W^{1,p}(\mathbb{R}^n) = W_0^{1,p}(\mathbb{R}^n \setminus \{0\})$ for satisfeita. 

Se $p \ge n$, a situação muda drasticamente. Para $p > n$, a imersão de Morrey garante que funções de $W^{1,p}$ são contínuas; assim, se $u(0) \neq 0$, é impossível aproximar $u$ por funções que são nulas em torno da origem sem explodir a norma do gradiente. Portanto, para $p \ge n$, o ponto $\{0\}$ possui capacidade positiva e a igualdade $W^{1,p} = W_0^{1,p}(\mathbb{R}^n \setminus \{0\})$ falha.
}

%%%% QUESTÃO 102 %%%%
\questao{Considere $\rho_\varepsilon : [0,1) \to \mathbb{R}$ dada por $\rho_\varepsilon(r) = 1$ em $0 \le r \le 1 - 2\varepsilon$, $\rho_\varepsilon(r) = \frac{1-r}{\varepsilon} - 1$ se $1 - 2\varepsilon \le r \le 1 - \varepsilon$ e $\rho_\varepsilon(r) = 0$ em $1 - \varepsilon \le r \le 1$. Mostre que:
	\begin{enumerate}[label=(\alph*)]
		\item $v_\varepsilon(x) = \rho_\varepsilon(|x|)$ tem derivadas fracas dadas por $\nabla v_\varepsilon(x) = \rho_\varepsilon'(|x|) \frac{x}{|x|}$ e $|\nabla v_\varepsilon| = \frac{1}{\varepsilon}$, $\forall 1 - 2\varepsilon \le |x| \le 1 - \varepsilon$ q.s.;
		\item $0 \le v_\varepsilon \le 1$, $\text{supp } v_\varepsilon = \overline{B_{1-\varepsilon}(0)}$;
		\item $v_\varepsilon \in W^{1,p}(\mathbb{R}^n)$;
		\item De fato, $v_\varepsilon(x)$ está em $W_0^{1,p}(B_1(0))$ e
		$$\|\nabla v_\varepsilon\|_{L^p(\mathbb{R}^n)}^p = \frac{|B_{1-\varepsilon}(0) \setminus B_{1-2\varepsilon}(0)|}{\varepsilon^p} \ge n\alpha(n)(1-2\varepsilon)^{n-1}\varepsilon^{1-p}.$$
		\item se $p > 1$ então $\|\nabla v_\varepsilon\|_{L^p(\mathbb{R}^n)} \to \infty$, quando $\varepsilon \to 0$.
\end{enumerate}}

\analise{
	\textbf{Análise do Enunciado:}
	A questão explora a relação entre a amplitude de uma função e a norma de seu gradiente em anéis anulares finos. A função $v_\varepsilon$ é construída para ser unitária na maior parte da bola $B_1(0)$, mas decair linearmente para zero em uma zona de transição de largura $\varepsilon$ próxima à fronteira. A tese central é que, para $p > 1$, a norma do gradiente explode quando a zona de transição encolhe, provando que a função constante $1$ não pertence a $W_0^{1,p}(B_1(0))$.
	
	\textbf{Fundamentação Teórica:}
	A prova utiliza a teoria de funções radialmente simétricas e a definição de derivada fraca para funções Lipschitz. Como $\rho_\varepsilon$ é Lipschitziana, $v_\varepsilon$ também o é, garantindo a existência de derivadas fracas via \textbf{Questão 98}. O cálculo da norma $L^p$ será realizado via integração em coordenadas esféricas, utilizando $\alpha(n)$ para denotar o volume da bola unitária e $\omega_n = n\alpha(n)$ para a área da esfera unitária.
	
	\textbf{Estratégias:}
	\begin{enumerate}
		\item \textbf{Derivação Radial}: Calcular $\rho_\varepsilon'(r)$ por partes e aplicar a regra da cadeia para obter $\nabla v_\varepsilon$, validando a norma $1/\varepsilon$ na zona de transição.
		\item \textbf{Análise de Suporte}: Verificar a continuidade de $\rho_\varepsilon$ para confirmar que o suporte é a bola de raio $1-\varepsilon$.
		\item \textbf{Pertinência a $W^{1,p}$}: Justificar a integrabilidade de $v_\varepsilon$ e $\nabla v_\varepsilon$ em $\mathbb{R}^n$ devido ao suporte compacto e à limitação do gradiente para $\varepsilon$ fixo.
		\item \textbf{Estimativa do Gradiente}: Integrar $|\nabla v_\varepsilon|^p$ sobre o anel $B_{1-\varepsilon} \setminus B_{1-2\varepsilon}$. Para a desigualdade inferior, utilizar a cota inferior do raio no anel para estimar o volume.
		\item \textbf{Análise do Limite}: Mostrar que, se $p > 1$, o expoente $1-p$ é negativo, levando a norma ao infinito quando $\varepsilon \to 0$.
	\end{enumerate}
}

\solucao{
	\textbf{(a)} A função $\rho_\varepsilon(r)$ é contínua e Lipschitziana em $[0, 1)$. Sua derivada clássica $\rho_\varepsilon'(r)$ é dada quase sempre por:
	$$\rho_\varepsilon'(r) = \begin{cases} 0, & \text{se } 0 \le r < 1 - 2\varepsilon \\ -\frac{1}{\varepsilon}, & \text{se } 1 - 2\varepsilon < r < 1 - \varepsilon \\ 0, & \text{se } 1 - \varepsilon < r < 1 \end{cases}$$
	Como $\rho_\varepsilon$ e a norma $|x|$ são Lipschitz, $v_\varepsilon(x) = \rho_\varepsilon(|x|)$ é Lipschitz em $\mathbb{R}^n$. Pela \textbf{Questão 98}, $v_\varepsilon \in W^{1, \infty}_{loc}(\mathbb{R}^n)$. Pela regra da cadeia e pelo resultado da \textbf{Questão 91}, a derivada fraca é:
	$$\nabla v_\varepsilon(x) = \rho_\varepsilon'(|x|) \frac{x}{|x|}$$
	No anel $1 - 2\varepsilon \le |x| \le 1 - \varepsilon$, temos $\rho_\varepsilon'(|x|) = -1/\varepsilon$. Assim, $|\nabla v_\varepsilon(x)| = |-\frac{1}{\varepsilon}| \cdot |\frac{x}{|x|}| = \frac{1}{\varepsilon}$ q.s.
	
	\textbf{(b)} Da definição de $\rho_\varepsilon$, observamos que a função decai linearmente de $1$ para $0$ no intervalo $[1-2\varepsilon, 1-\varepsilon]$ e é nula para $r \ge 1-\varepsilon$. Logo, $0 \le v_\varepsilon \le 1$. O suporte de $v_\varepsilon$ é o fecho do conjunto onde a função é não nula, resultando em $\text{supp } v_\varepsilon = \overline{B_{1-\varepsilon}(0)}$.
	
	\textbf{(c)} Para $v_\varepsilon \in W^{1,p}(\mathbb{R}^n)$, verificamos a integrabilidade:
	\begin{itemize}
		\item Como $|v_\varepsilon| \le 1$ e $\text{supp } v_\varepsilon$ é compacto, $\|v_\varepsilon\|_{L^p(\mathbb{R}^n)}^p \le |B_{1-\varepsilon}(0)| = \alpha(n)(1-\varepsilon)^n < \infty$.
		\item O gradiente $\nabla v_\varepsilon$ é nulo fora do anel $B_{1-\varepsilon}(0) \setminus B_{1-2\varepsilon}(0)$ e limitado por $1/\varepsilon$ dentro dele. Logo, 
		$$\|\nabla v_\varepsilon\|_{L^p(\mathbb{R}^n)}^p \le (1/\varepsilon)^p |B_{1-\varepsilon}(0)| = (1/\varepsilon)^p \alpha(n)(1-\varepsilon)^n < \infty.$$
	\end{itemize}
	Sendo a função e seu gradiente integráveis, concluímos que $v_\varepsilon \in W^{1,p}(\mathbb{R}^n)$.
	
	\textbf{(d)} Como $\text{supp}(v_\varepsilon) = \overline{B_{1-\varepsilon}(0)} \subset B_1(0)$, a função $v_\varepsilon$ pode ser estendida por zero para fora de $B_1(0)$. Sendo Lipschitz com suporte compacto em $B_1(0)$, segue que $v_\varepsilon \in W_0^{1,p}(B_1(0))$.
	Para a norma do gradiente, integramos sobre o anel de transição:
	$$\|\nabla v_\varepsilon\|_{L^p(\mathbb{R}^n)}^p = \int_{B_{1-\varepsilon}(0) \setminus B_{1-2\varepsilon}(0)} \left| \frac{1}{\varepsilon} \right|^p dx = \frac{|B_{1-\varepsilon}(0) \setminus B_{1-2\varepsilon}(0)|}{\varepsilon^p}$$
	Calculamos o volume do anel via coordenadas esféricas:
	$$|B_{1-\varepsilon}(0) \setminus B_{1-2\varepsilon}(0)| = \int_{1-2\varepsilon}^{1-\varepsilon} \omega_n r^{n-1} dr$$
	Visto que $r \ge 1-2\varepsilon$ em todo o intervalo de integração, temos:
	$$|B_{1-\varepsilon}(0) \setminus B_{1-2\varepsilon}(0)| \ge \omega_n (1-2\varepsilon)^{n-1} \int_{1-2\varepsilon}^{1-\varepsilon} dr = n\alpha(n)(1-2\varepsilon)^{n-1}\varepsilon$$
	Substituindo na norma:
	$$\boxed{\|\nabla v_\varepsilon\|_{L^p(\mathbb{R}^n)}^p \ge \frac{n\alpha(n)(1-2\varepsilon)^{n-1}\varepsilon}{\varepsilon^p} = n\alpha(n)(1-2\varepsilon)^{n-1}\varepsilon^{1-p}}$$
	
	\textbf{(e)} Se $p > 1$, então o expoente $1-p$ é estritamente negativo. Quando $\varepsilon \to 0$, o termo $\varepsilon^{1-p} = \frac{1}{\varepsilon^{p-1}} \to \infty$. Como $n\alpha(n)(1-2\varepsilon)^{n-1} \to n\alpha(n) > 0$, concluímos que:
	$$\lim_{\varepsilon \to 0} \|\nabla v_\varepsilon\|_{L^p(\mathbb{R}^n)}^p = \infty \implies \boxed{\|\nabla v_\varepsilon\|_{L^p(\mathbb{R}^n)} \to \infty.}$$
}

\sintese{
	A prova demonstra que tentar aproximar a função constante $1$ por funções de $W_0^{1,p}$ em $B_1(0)$ exige um custo de energia (norma do gradiente) que explode para $p > 1$. O fluxo lógico foi: 
	$$\text{Definição de } \rho_\varepsilon \to \text{Gradiente no Anel} \to \text{Volume do Anel} \to \text{Análise do expoente } 1-p \to \text{Explosão da Norma}.$$
}

\aprofundamento{
	Este resultado é a contrapartida da Questão 100. Enquanto na questão anterior vimos que a energia pode sumir ao concentrar a função em um ponto (se $n > p$), aqui vemos que a energia explode ao tentar forçar a função a zero na fronteira em um intervalo pequeno. 
	
	Isso prova que a constante $1$ não pertence ao espaço $W_0^{1,p}(B_1(0))$ para $p > 1$. Em termos de análise funcional, isso significa que a aplicação de traço $T: W^{1,p}(B_1) \to L^p(\partial B_1)$ é bem definida e não trivial: a condição de ``ser zero na fronteira'' impõe uma restrição real e rigorosa à função, que não pode ser ignorada através de aproximações suaves sem que a norma de Sobolev exploda.
}

%%%% QUESTÃO 103 %%%%
\questao{Seja $u \in W^{1,p}(U)$. Mostre que para toda $\phi \in C^1(\overline{U}) \cap W^{1,q}(U)$ (aqui $U$ pode ser não limitado), $\frac{1}{p} + \frac{1}{q} = 1$, tal que $\phi = 0$ sobre $\partial U$:
	$$\int_U u \frac{\partial \phi}{\partial x_j} dx = - \int_U \phi \frac{\partial u}{\partial x_j} dx.$$
}
	
\analise{
\textbf{Análise do Enunciado:} 
A questão solicita a extensão da fórmula de integração por partes (definidora da derivada fraca) para funções de teste que não possuem, necessariamente, suporte compacto, mas que pertencem ao espaço $W^{1,q}(U)$ e anulam-se na fronteira. O desafio principal reside no fato de que, se $U$ for não limitado, a integral pode divergir ou apresentar termos de borda no infinito. A tese é que a integrabilidade de $u$ em $L^p$ e de $\phi$ em $W^{1,q}$ é suficiente para garantir que a identidade de integração por partes permaneça válida.

\textbf{Fundamentação Teórica:}
A prova baseia-se na densidade de $C_c^\infty(U)$ em $W_0^{1,q}(U)$ e na técnica de truncação via funções de corte ($\rho_R$), já explorada na \textbf{Questão 99}. Utilizaremos a \textbf{Desigualdade de Hölder} para controlar os produtos $u \partial_j \phi$ e $\phi \partial_j u$, visto que a soma dos inversos dos expoentes é unitária, $1/p + 1/q = 1$.

\textbf{Estratégias:}
\begin{enumerate}
	\item \textbf{Construção da Função de Corte}: Definir uma sequência $\rho_R \in C_c^\infty(\mathbb{R}^n)$ tal que $\rho_R = 1$ em $B_R(0)$ e $\rho_R = 0$ fora de $B_{2R}(0)$.
	\item \textbf{Regularização do Teste}: Definir $\phi_R = \phi \rho_R$. Visto que $\phi = 0$ em $\partial U$, a função $\phi_R$ terá suporte compacto em $U$, permitindo o uso da definição de derivada fraca de $u$.
	\item \textbf{Expansão da Derivada}: Aplicar a regra do produto $\partial_j \phi_R = \rho_R \partial_j \phi + \phi \partial_j \rho_R$ e analisar a integral resultante.
	\item \textbf{Passagem ao Limite}: Demonstrar que, quando $R \to \infty$, o termo envolvendo $\partial_j \rho_R$ tende a zero e o termo com $\rho_R \partial_j \phi$ converge para a integral original via Teorema da Convergência Dominada.
\end{enumerate}
}

\solucao{
Seja $\rho \in C_c^\infty(\mathbb{R}^n)$ tal que $0 \le \rho \le 1$, $\rho(x) = 1$ para $|x| \le 1$ e $\rho(x) = 0$ para $|x| \ge 2$. Para cada $R > 0$, definimos $\rho_R(x) = \rho(x/R)$. 

Consideremos a função $\phi_R(x) = \phi(x) \rho_R(x)$. Como $\phi = 0$ em $\partial U$ e $\rho_R$ tem suporte compacto em $B_{2R}(0)$, a função $\phi_R$ possui suporte compacto contido em $U$. Portanto, $\phi_R$ é uma função de teste válida para a definição de derivada fraca de $u \in W^{1,p}(U)$. Assim, temos:
$$\int_U u \frac{\partial \phi_R}{\partial x_j} dx = - \int_U \phi_R \frac{\partial u}{\partial x_j} dx \quad (\ast_1)$$

Expandindo a derivada $\frac{\partial \phi_R}{\partial x_j}$ pela regra do produto, obtemos:
$$\int_U u \left( \rho_R \frac{\partial \phi}{\partial x_j} + \phi \frac{\partial \rho_R}{\partial x_j} \right) dx = - \int_U \phi \rho_R \frac{\partial u}{\partial x_j} dx$$
$$\int_U \rho_R u \frac{\partial \phi}{\partial x_j} dx + \int_U u \phi \frac{\partial \rho_R}{\partial x_j} dx = - \int_U \phi \rho_R \frac{\partial u}{\partial x_j} dx \quad (\ast_2)$$

Analisamos agora o comportamento de cada termo quando $R \to \infty$:

\textit{1. O termo $\int_U \rho_R u \frac{\partial \phi}{\partial x_j} dx$}: Como $u \in L^p(U)$ e $\frac{\partial \phi}{\partial x_j} \in L^q(U)$, o produto $u \frac{\partial \phi}{\partial x_j}$ pertence a $L^1(U)$ por Hölder. Visto que $\rho_R(x) \to 1$ pontualmente quando $R \to \infty$ e $|\rho_R| \le 1$, o Teorema da Convergência Dominada garante que:
$$\lim_{R \to \infty} \int_U \rho_R u \frac{\partial \phi}{\partial x_j} dx = \int_U u \frac{\partial \phi}{\partial x_j} dx$$

\textit{2. O termo $\int_U \phi \rho_R \frac{\partial u}{\partial x_j} dx$}:
De forma análoga, como $\phi \in L^q(U)$ e $\frac{\partial u}{\partial x_j} \in L^p(U)$, o produto $\phi \frac{\partial u}{\partial x_j}$ é integrável. Assim:
$$\lim_{R \to \infty} - \int_U \phi \rho_R \frac{\partial u}{\partial x_j} dx = - \int_U \phi \frac{\partial u}{\partial x_j} dx$$

\textit{3. O termo de erro $\int_U u \phi \frac{\partial \rho_R}{\partial x_j} dx$}:
Lembremos que $\nabla \rho_R(x) = \frac{1}{R} \nabla \rho(x/R)$. Logo, $|\frac{\partial \rho_R}{\partial x_j}| \le \frac{C}{R}$ para alguma constante $C$. Temos:
$$\left| \int_U u \phi \frac{\partial \rho_R}{\partial x_j} dx \right| \le \int_{R < |x| < 2R} |u \phi| \frac{C}{R} dx \le \frac{C}{R} \int_{R < |x| < 2R} |u \phi| dx$$
Como $u \in L^p(U)$ e $\phi \in L^q(U)$, a integral $\int_U |u \phi| dx$ é finita (Hölder). Portanto, a integral sobre a casca anular tende a zero quando $R \to \infty$. Além disso, o fator $1/R$ acelera essa convergência:
$$\lim_{R \to \infty} \int_U u \phi \frac{\partial \rho_R}{\partial x_j} dx = 0$$

Substituindo esses limites na identidade $(\ast_2)$, obtemos finalmente:
$$\boxed{\int_U u \frac{\partial \phi}{\partial x_j} dx = - \int_U \phi \frac{\partial u}{\partial x_j} dx}$$
}

\sintese{
A prova demonstra que a integrabilidade global de $u$ e $\phi$ em espaços conjugados ($L^p$ e $L^q$) é suficiente para anular a contribuição da fronteira no infinito. O fluxo lógico foi: 
$$\text{Construção de } \phi_R \in C_c^\infty \to \text{Controle do termo de erro via } \nabla \rho_R \to \text{Limite por Convergência Dominada}.$$
}

\aprofundamento{
Este resultado é fundamental para a definição de operadores diferenciais em domínios não limitados. Ele prova que a integração por partes não requer suporte compacto, desde que a função e sua derivada \textbf{decaiam} suficientemente rápido no infinito para pertencerem a $L^p$ e $L^q$. 

Geometricamente, isso significa que a \textbf{fronteira no infinito} de um espaço de Sobolev $W^{1,p}(\mathbb{R}^n)$ comporta-se como se a função fosse nula lá. Esse princípio é a base para a aplicação de transformadas de Fourier e para a prova de que $W^{1,p}(\mathbb{R}^n) = W_0^{1,p}(\mathbb{R}^n)$, como visto na Questão 99.
}

%%%% QUESTÃO 104 %%%%
\questao{Considere $r > 0$, $B = B_r(x_0) \subset \mathbb{R}^{n-1}$, $x_0 \in \mathbb{R}^{n-1}$, e uma função $\psi : B \to \mathbb{R}$ de classe $C^1(\overline{B})$. Defina
	\begin{align*}
		Q_+ &= \{z = (x, x_n) : \psi(x) < x_n < \psi(x) + r\}, \\
		Q_- &= \{w = (x, y_n) : \psi(x) - r < y_n < \psi(x)\}, \\
		G &= \{z = (x, x_n) : x_n = \psi(x)\}, \\
		H &: B \times \mathbb{R} \to B \times \mathbb{R} \text{ dado por } H(x, x_n) = (x, 2\psi(x) - x_n) \\
		\text{e } Q &= \{z = (x, x_n) : \psi(x) - r < x_n < \psi(x) + r\}.
	\end{align*}
	Mostre que:
	\begin{enumerate}[label=(\alph*)]
		\item $Q = Q_- \cup G \cup Q_+$;
		\item $H(Q_+) = Q_-$, $H(Q_-) = Q_+$ e $H \circ H = I$, ou seja, $H(H(x, x_n)) = (x, x_n)$;
		\item $\det H'(z) = -1$ e para $f \in L^p(Q_+)$
		$$\int_{Q_-} f(H(w))g(w) dw = \int_{Q_+} f(z)g(H(z)) dz;$$
		\item Se $u \in W^{1,p}(Q_+)$ é tal que $\supp u \subset Q_+ \cup G$, mostre que
		$$u^*(z) = \begin{cases} u(z), & z \in Q_+ \\ u(H(z)), & z = (x, x_n) \in Q_- \end{cases}$$
		está em $W^{1,p}(Q)$;
		\item Para $j = 1, 2, \dots, n-1$:
		$$\frac{\partial u^*}{\partial x_j}(z) = \begin{cases} \tfrac{\partial u}{\partial x_j}(z), & z \in Q_+ \\ \tfrac{\partial u}{\partial x_j}(H(z)) + 2 \tfrac{\partial u}{\partial x_n}(H(z)) \tfrac{\partial \psi}{\partial x_j}(x), & z = (x, x_n) \in Q_- \end{cases}$$
		\item Quando $j = n$:
		$$\frac{\partial u^*}{\partial x_n}(z) = \begin{cases} \tfrac{\partial u}{\partial x_n}(z), & z \in Q_+ \\ 	
		- \tfrac{\partial u}{\partial x_n}(H(z)), & z = (x, x_n) \in Q_- \end{cases}$$
		\item Mostre uma desigualdade de imersão assim: Existe uma constante $C = C(\|\psi\|_\infty, \|\nabla \psi\|_\infty)$ tal que:
		$$\|u^*\|_{W^{1,p}(Q)} \le C \|u\|_{W^{1,p}(Q_+)}$$
\end{enumerate}}

\analise{
	\textbf{Análise do Enunciado:}
	O problema propõe a construção de uma extensão por reflexão de uma função de Sobolev definida em um domínio limitado superiormente por um grafo. O ponto central é a definição do mapa $H$, que reflete pontos de $Q_+$ em $Q_-$ em relação à superfície $G$. A principal dificuldade técnica reside no fato de que $H$ não é uma isometria linear, mas uma transformação dependente da geometria de $\psi$, o que altera a expressão das derivadas via regra da cadeia e requer o uso do determinante do Jacobiano para a mudança de variáveis em $L^p$.
	
	\textbf{Fundamentação Teórica:}
	A prova baseia-se na definição de derivada fraca e no teorema de mudança de variáveis para integrais múltiplas. Para provar que $u^* \in W^{1,p}(Q)$, utilizaremos a propriedade de que a função $u$ anula-se no traço sobre $G$ ($\supp u \subset Q_+ \cup G$), eliminando a contribuição de saltos na interface durante a integração por partes.
	
	\textbf{Estratégias:}
	\begin{enumerate}
		\item \textbf{Verificação Geométrica}: Provar a partição de $Q$ e as propriedades do mapa $H$ via substituição direta nas definições de $Q_+$ e $Q_-$.
		\item \textbf{Cálculo do Jacobiano}: Calcular a matriz derivada $H'$ e seu determinante para validar a mudança de variáveis em $L^p$.
		\item \textbf{Validação do Espaço de Sobolev}: Usar a definição de derivada fraca e a hipótese de suporte para mostrar que a junção de $u$ e $u \circ H$ não gera distribuições singulares (como a Delta de Dirac) na superfície $G$.
		\item \textbf{Regra da Cadeia}: Aplicar a diferenciação de funções compostas para derivar $u^*(z) = u(H(z))$, considerando a dependência de $H$ em relação a $x$ e $x_n$.
		\item \textbf{Estimativa de Norma}: Utilizar a limitação de $\nabla \psi$ e a mudança de variáveis para cotar a norma de $u^*$ em termos da norma de $u$.
	\end{enumerate}
}

\solucao{
	\textbf{(a)} O conjunto $Q$ é definido como $Q = \{ (x, x_n) \in B \times \mathbb{R} : \psi(x) - r < x_n < \psi(x) + r \}$.
	Para qualquer ponto $z = (x, x_n) \in Q$, a componente $x_n$ deve satisfazer a desigualdade $\psi(x) - r < x_n < \psi(x) + r$. Pela lei da ordem em $\mathbb{R}$, existem três possibilidades para $x_n$ em relação a $\psi(x)$:
	\begin{enumerate}
		\item $x_n > \psi(x)$: Neste caso, combinando com a definição de $Q$, temos $\psi(x) < x_n < \psi(x) + r$, o que implica $z \in Q_+$.
		\item $x_n < \psi(x)$: Neste caso, temos $\psi(x) - r < x_n < \psi(x)$, o que implica $z \in Q_-$.
		\item $x_n = \psi(x)$: Neste caso, o ponto satisfaz a condição de igualdade, logo $z \in G$.
	\end{enumerate}
	Como cada ponto de $Q$ cai necessariamente em um desses três casos, segue que $\boxed{Q = Q_- \cup G \cup Q_+}$.
	
	\textbf{(b)} Seja $z = (x, x_n)$. O mapa de reflexão é $H(z) = (x, 2\psi(x) - x_n)$.
	\begin{itemize}
		\item Se $z \in Q_+$, então $\psi(x) < x_n < \psi(x) + r$. Multiplicando por $-1$, temos $-\psi(x) - r < -x_n < -\psi(x)$. Somando $2\psi(x)$ em todos os termos:
		$$\psi(x) - r < 2\psi(x) - x_n < \psi(x)$$
		Isso prova que $H(z) \in Q_-$.
		\item Se $z \in Q_-$, então $\psi(x) - r < x_n < \psi(x)$. Multiplicando por $-1$, temos $-\psi(x) < -x_n < -\psi(x) + r$. Somando $2\psi(x)$:
		$$\psi(x) < 2\psi(x) - x_n < \psi(x) + r$$
		Isso prova que $H(z) \in Q_+$.
	\end{itemize}
	Finalmente, para a involução $H \circ H$:
	$$H(H(x, x_n)) = H(x, 2\psi(x) - x_n) = (x, 2\psi(x) - (2\psi(x) - x_n)) = (x, x_n)$$
	Logo, $\boxed{H \circ H = I}$.
	
	\textbf{(c)} A matriz Jacobiana de $H$ é obtida derivando $H_k$ em relação a $z_j$. Para $k < n$, $H_k(x, x_n) = x_k$, logo $\partial_j H_k = \delta_{jk}$. Para $k = n$, $H_n(x, x_n) = 2\psi(x) - x_n$, logo $\partial_j H_n = 2\partial_j \psi(x)$ para $j < n$ e $\partial_n H_n = -1$. A matriz é:
	$$H'(z) = \begin{pmatrix} 
		1 & \dots & 0 & 0 \\ 
		\vdots & \ddots & \vdots & \vdots \\ 
		0 & \dots & 1 & 0 \\ 
		2\partial_1 \psi(x) & \dots & 2\partial_{n-1} \psi(x) & -1 
	\end{pmatrix}$$
	O determinante de uma matriz triangular inferior é o produto dos elementos da diagonal: $\det H' = (1)^{n-1} \cdot (-1) = -1$.
	Para a integral, usamos a mudança de variáveis $z = H(w)$. O Jacobiano é $|\det H'| = 1$. Como $H(Q_-) = Q_+$, temos:
	$$\int_{Q_-} f(H(w))g(w) dw = \int_{Q_+} f(z)g(H(z)) \frac{1}{|\det H'|} dz = \int_{Q_+} f(z)g(H(z)) dz$$
	
	\textbf{(d)} Para provar que $u^* \in W^{1,p}(Q)$, devemos mostrar que existe um vetor $\mathbf{g} \in L^p(Q; \mathbb{R}^n)$ tal que, para toda função de teste $\phi \in C_c^\infty(Q)$, a identidade $\int_Q u^* \partial_j \phi \, dx = - \int_Q g_j \phi \, dx$ seja satisfeita para cada $j=1, \dots, n$.
	
	Definimos o candidato a derivada fraca $g_j$ como a função que coincide com $\partial_j u$ em $Q_+$ e com as expressões deduzidas nos itens \textbf{(e)} e \textbf{(f)} em $Q_-$. Consideremos a integral:
	$$\int_Q u^* \partial_j \phi \, dx = \int_{Q_+} u \partial_j \phi \, dx + \int_{Q_-} u(H(z)) \partial_j \phi(z) dz \quad (\ast_1)$$
	No primeiro termo, aplicamos a fórmula de integração por partes para funções de Sobolev. Como $u \in W^{1,p}(Q_+)$, temos:
	$$\int_{Q_+} u \partial_j \phi \, dx = - \int_{Q_+} \partial_j u \phi \, dx + \int_{\partial Q_+} u \phi \nu_j d\sigma$$
	Onde $\nu$ é a normal exterior a $Q_+$. A fronteira $\partial Q_+$ contém a superfície $G$. A hipótese $\text{supp } u \subset Q_+ \cup G$ implica que $u$ anula-se no traço sobre a fronteira $\partial Q_+ \setminus G$ (pois $\phi$ tem suporte compacto em $Q$) e, crucialmente, que o traço de $u$ sobre $G$ é nulo ($u|_G = 0$). Portanto, a integral de superfície $\int_{\partial Q_+} u \phi \nu_j d\sigma$ é nula.
	
	Para o segundo termo em $(\ast_1)$, efetuamos a mudança de variáveis $z = H(w)$. Como $\det H' = -1$, temos $dz = dw$ e $H(Q_-) = Q_+$. A integral torna-se:
	$$\int_{Q_-} u(H(z)) \partial_j \phi(z) dz = \int_{Q_+} u(w) \partial_j \phi(H(w)) dw$$
	Agora, aplicamos novamente a integração por partes em $Q_+$ para a função $u(w)$:
	$$\int_{Q_+} u(w) \partial_j (\phi \circ H)(w) dw = - \int_{Q_+} \partial_j u(w) \phi(H(w)) dw + \int_{\partial Q_+} u(w) \phi(H(w)) \nu_j d\sigma$$
	Novamente, a integral de superfície anula-se pois $u|_G = 0$. Retornando para a variável $z$ via $w = H(z)$, obtemos:
	$$\int_{Q_-} u(H(z)) \partial_j \phi(z) dz = - \int_{Q_-} \left[ \nabla u(H(z)) \cdot \partial_j H(z) \right] \phi(z) dz$$
	Somando as duas contribuições, a integral total sobre $Q$ é:
	$$\int_Q u^* \partial_j \phi \, dx = - \left( \int_{Q_+} \partial_j u \phi \, dx + \int_{Q_-} \left[ \nabla u(H(z)) \cdot \partial_j H(z) \right] \phi(z) dz \right)$$
	Esta expressão coincide exatamente com a definição de derivada fraca $\int_Q u^* \partial_j \phi \, dx = - \int_Q \partial_j u^* \phi \, dx$, onde $\partial_j u^*$ é a função definida por partes nos itens seguintes. Como $u \in W^{1,p}(Q_+)$ e $|\det H'|=1$, a função resultante pertence a $L^p(Q)$. Logo, $\boxed{u^* \in W^{1,p}(Q)}$.
	
	\textbf{(e)} Para $z \in Q_-$, temos a definição $u^*(z) = u(H(z))$. Para calcular a derivada fraca em relação a $x_j$ ($j < n$), aplicamos a regra da cadeia para funções em espaços de Sobolev:
	$$\frac{\partial u^*}{\partial x_j}(z) = \sum_{k=1}^{n} \frac{\partial u}{\partial x_k}(H(z)) \frac{\partial H_k}{\partial x_j}(z) \quad (\ast_2)$$
	Analisamos as derivadas parciais do mapa de reflexão $H(x, x_n) = (x_1, \dots, x_{n-1}, 2\psi(x) - x_n)$:
	\begin{itemize}
		\item Para $k < n$, $H_k = x_k$, logo $\frac{\partial H_k}{\partial x_j} = \delta_{kj}$.
		\item Para $k = n$, $H_n = 2\psi(x) - x_n$, logo $\frac{\partial H_n}{\partial x_j} = 2\partial_j \psi(x)$.
	\end{itemize}
	Substituindo esses valores em $(\ast_2)$, obtemos:
	$$\frac{\partial u^*}{\partial x_j}(z) = \frac{\partial u}{\partial x_j}(H(z)) \cdot 1 + \frac{\partial u}{\partial x_n}(H(z)) \cdot 2\partial_j \psi(x)$$
	$$\boxed{\frac{\partial u^*}{\partial x_j}(z) = \frac{\partial u}{\partial x_j}(H(z)) + 2 \frac{\partial u}{\partial x_n}(H(z)) \frac{\partial \psi}{\partial x_j}(x), \quad z \in Q_-}$$
	
	\textbf{(f)} Para $j = n$, aplicamos a mesma regra da cadeia $(\ast_2)$. Observamos que:
	\begin{itemize}
		\item Para $k < n$, $\frac{\partial H_k}{\partial x_n} = \frac{\partial x_k}{\partial x_n} = 0$.
		\item Para $k = n$, $\frac{\partial H_n}{\partial x_n} = \frac{\partial}{\partial x_n}(2\psi(x) - x_n) = -1$.
	\end{itemize}
	Substituindo na soma:
	$$\frac{\partial u^*}{\partial x_n}(z) = \sum_{k=1}^{n-1} \frac{\partial u}{\partial x_k}(H(z)) \cdot 0 + \frac{\partial u}{\partial x_n}(H(z)) \cdot (-1)$$
	$$\boxed{\frac{\partial u^*}{\partial x_n}(z) = -\frac{\partial u}{\partial x_n}(H(z)), \quad z \in Q_-}$$
	
	\textbf{(g)} Estimamos a norma $\|u^*\|_{W^{1,p}(Q)}^p = \|u^*\|_{L^p(Q)}^p + \|\nabla u^*\|_{L^p(Q)}^p$.
	Para a norma $L^p$, utilizamos a mudança de variáveis $z = H(w)$ em $Q_-$, com $|\det H'| = 1$:
	$$\|u^*\|_{L^p(Q)}^p = \int_{Q_+} |u(z)|^p dz + \int_{Q_-} |u(H(z))|^p dz = 2 \int_{Q_+} |u(z)|^p dz = 2 \|u\|_{L^p(Q_+)}^p \quad (\ast_3)$$
	Para o gradiente em $Q_-$, utilizamos as fórmulas dos itens **(e)** e **(f)**. Pela desigualdade triangular:
	$$|\nabla u^*(z)| \le |\nabla u(H(z))| + 2 |\nabla u(H(z))| |\nabla \psi(x)| = |\nabla u(H(z))| (1 + 2\|\nabla \psi\|_\infty)$$
	Elevando à potência $p$ e integrando em $Q_-$:
	$$\int_{Q_-} |\nabla u^*(z)|^p dz \le (1 + 2\|\nabla \psi\|_\infty)^p \int_{Q_-} |\nabla u(H(z))|^p dz$$
	Novamente, via mudança de variáveis $z = H(w)$, a integral da direita é precisamente $\|\nabla u\|_{L^p(Q_+)}^p$. Somando as contribuições de $Q_+$ e $Q_-$, temos:
	$$\|\nabla u^*\|_{L^p(Q)}^p \le \|\nabla u\|_{L^p(Q_+)}^p + (1 + 2\|\nabla \psi\|_\infty)^p \|\nabla u\|_{L^p(Q_+)}^p = [1 + (1 + 2\|\nabla \psi\|_\infty)^p] \|\nabla u\|_{L^p(Q_+)}^p$$
	Combinando a estimativa de $(\ast_3)$ e a do gradiente, existe uma constante $C$ dependendo exclusivamente de $\|\psi\|_\infty$ e $\|\nabla \psi\|_\infty$ tal que:
	$$\boxed{\|u^*\|_{W^{1,p}(Q)} \le C \|u\|_{W^{1,p}(Q_+)}}$$
}

\sintese{
	A reflexão através de um grafo generaliza a reflexão especular. A chave da prova é que a transformação $H$ preserva a medida (determinante $-1$), mas distorce a direção do gradiente dependendo da inclinação da superfície $\psi$. A condição de suporte $\supp u \subset Q_+ \cup G$ é o que garante que a função estendida não possua uma \textbf{quebra} (descontinuidade) na interface $G$, permitindo que ela pertença ao espaço de Sobolev.
}

\aprofundamento{
	Este procedimento é a base para a prova do \textbf{Teorema de Extensão de Sobolev}, que afirma que para domínios com fronteira Lipschitz, existe um operador de extensão $E: W^{1,p}(\Omega) \to W^{1,p}(\mathbb{R}^n)$. A reflexão através do grafo é a forma mais simples de construir tal operador localmente. A desigualdade de imersão no item (g) mostra que a extensão é um operador contínuo entre espaços de Sobolev, preservando a regularidade da função original.
}

%%%% QUESTÃO 105 %%%%
\questao{
	Prove diretamente que se $u \in W^{1,p}(J), J = (0, 1) \subset \mathbb{R}, 1 < p < \infty$, então $u$ satisfaz
	$$|u(x) - u(y)| \le \|u'\|_{L^p(J)} |x - y|^{1 - \frac{1}{p}},$$
	para quase todos $x, y \in J$. Conclua que toda função de $W^{1,p}(J)$ é contínua.
}

\analise{
	\textbf{Análise do Enunciado}: A questão solicita a demonstração direta de que funções em $W^{1,p}$ em dimensão um possuem um representante Hölder contínuo. A tese é a obtenção de uma estimativa a priori que controle a variação pontual de $u$ através da norma $L^p$ de sua derivada fraca.
	
	\textbf{Fundamentação Teórica}: A prova baseia-se no Teorema Fundamental do Cálculo para funções de Sobolev e na \textbf{Desigualdade de Hölder}. De acordo com \textbf{Brezis} (Capítulo 8), toda função em $W^{1,p}(J)$ com $p > 1$ admite um representante contínuo em $\overline{J}$.
	
	\textbf{Estratégias}:
	\begin{enumerate}
		\item Utilizar a representação integral $u(x) - u(y) = \int_y^x u'(t) \, dt$ para funções em $W^{1,p}(J)$.
		\item Aplicar a Desigualdade de Hölder para isolar a norma da derivada $\|u'\|_{L^p}$ e a medida do intervalo $|x-y|$.
		\item Formalizar a passagem do \textbf{quase todo} para a continuidade uniforme, definindo explicitamente a dependência de $\delta$ em função de $\varepsilon$ e da norma de Sobolev.
	\end{enumerate}
}

\solucao{
	Seja $u \in W^{1,p}(J)$. Sabemos que $u$ possui um representante contínuo em $\overline{J}$ tal que, para quase todos $x, y \in J$, a seguinte representação integral é válida:
	$$u(x) - u(y) = \int_y^x u'(t) \, dt$$
	Assumindo, sem perda de generalidade, que $y < x$, aplicamos a \textbf{Desigualdade de Hölder} para o par de expoentes conjugados $p$ e $q$, onde $\tfrac{1}{p} + \tfrac{1}{q} = 1$:
	\begin{align*}
		|u(x) - u(y)| &= \left| \int_y^x u'(t) \cdot 1 \, dt \right| \\
		&\le \left( \int_y^x |u'(t)|^p \, dt \right)^{\frac{1}{p}} \left( \int_y^x 1^q \, dt \right)^{\frac{1}{q}} \\
		&\le \left( \int_J |u'(t)|^p \, dt \right)^{\frac{1}{p}} (x - y)^{\frac{1}{q}}
	\end{align*}
	
	Como $\tfrac{1}{q} = 1 - \tfrac{1}{p}$ e reconhecendo a norma $\|u'\|_{L^p(J)}$, obtemos a estimativa a priori:
	$$\boxed{|u(x) - u(y)| \le \|u'\|_{L^p(J)} |x - y|^{1 - \frac{1}{p}}\quad (\ast)}$$
	
	Para provar que $u$ é contínua, devemos mostrar que satisfaz a definição de continuidade uniforme em $\overline{J}$. Seja $\varepsilon > 0$ um valor arbitrário. Precisamos encontrar um $\delta > 0$ tal que, para quaisquer $x, y \in \overline{J}$, a condição $|x - y| < \delta$ implique $|u(x) - u(y)| < \varepsilon$.
	
	Denotemos a constante $M = \|u'\|_{L^p(J)}$ e o expoente $\gamma = 1 - \frac{1}{p}$. Pela desigualdade $(\ast)$, temos:
	$$|u(x) - u(y)| \le M |x - y|^\gamma$$
	Para que $|u(x) - u(y)| < \varepsilon$, basta que:
	$$M |x - y|^\gamma < \varepsilon \implies |x - y|^\gamma < \frac{\varepsilon}{M} \implies |x - y| < \left( \frac{\varepsilon}{M} \right)^{\frac{1}{\gamma}}$$
	Portanto, definindo $\delta = \left( \frac{\varepsilon}{\|u'\|_{L^p(J)}} \right)^{\frac{1}{1 - 1/p}}$, verificamos que:
	$$|x - y| < \delta \implies |u(x) - u(y)| \le M \left( \left( \frac{\varepsilon}{M} \right)^{\frac{1}{\gamma}} \right)^\gamma = M \left( \frac{\varepsilon}{M} \right) = \varepsilon$$
	Como $\delta$ depende exclusivamente de $\varepsilon$ e da norma de $u$ no espaço de Sobolev (e não da escolha de $x$ ou $y$), a função $u$ é \textbf{uniformemente contínua} em $\overline{J}$. Consequentemente, toda função de $W^{1,p}(J)$ possui um representante contínuo.
}

\sintese{
	A prova demonstra que a integrabilidade da derivada em $L^p$ com $p > 1$ é forte o suficiente para impedir saltos ou oscilações infinitas. A transição da norma global para a regularidade pontual é feita através da desigualdade de Hölder, que fornece a "taxa" de variação da função.
	
	O fluxo lógico foi:
	$$\boxed{\text{Repres. Integral} \to \text{Desig. de Hölder} \to \text{Estimativa de Hölder} \to \text{Definição de } \delta(\varepsilon) \to \text{Cont. Uniforme}}$$
}

\aprofundamento{
	Este resultado é a base da imersão de Sobolev $W^{1,p}(J) \hookrightarrow C^{0, \gamma}(\overline{J})$ com $\gamma = 1 - 1/p$. Note que se $p=1$, a estimativa resultaria em $|u(x)-u(y)| \le \|u'\|_{L^1} \cdot 1$, o que garante apenas que a função é de variação limitada ($\text{BV}$), mas não necessariamente contínua. A condição $p > 1$ é, portanto, a fronteira analítica que garante a existência de um representante contínuo.
}

%%%% QUESTÃO 106 %%%%
\questao{
	Suponha que $u \in L^p(U)$, $1 < p < \infty$, e que existe $C > 0$ tal que: para todo $\Omega \subset\subset U$ e $|h| < \text{dist}(\Omega, \partial U)$, então
	$$\|\tau_h u - u\|_{L^p(\Omega)} \le C|h|.$$
	Mostre que: $u \in W^{1,p}(U)$ e que $\|\nabla u\|_{L^p(U)}$ é a menor constante $C$ na desigualdade acima.
}
\analise{
	\textbf{Análise do Enunciado}: A questão propõe a equivalência entre a regularidade de Sobolev $W^{1,p}$ e a estabilidade da função sob translações no espaço $L^p$. O objetivo é provar que a limitação do quociente de diferença implica a existência de derivadas fracas e que a norma do gradiente é a constante ótima.
	
	\textbf{Fundamentação Teórica}: A demonstração baseia-se na \textbf{Proposição 9.3 de Haim Brezis}. O pilar central é a representação de funcionais lineares via dualidade: o dual do espaço $L^{p'}(\Omega)$ é isomorfo ao espaço $L^p(\Omega)$, onde $p'$ é o expoente conjugado $\left(\frac{1}{p} + \frac{1}{p'} = 1\right)$. Se um funcional linear é limitado em $L^{p'}$, ele é representado por uma função em $L^p$.
	
	\textbf{Estratégias}:
	\begin{enumerate}
		\item \textbf{Construção do Funcional}: Definir o funcional linear $L(\phi) = - \int u \partial_j \phi$ e provar que ele é limitado em $L^{p'}$ usando a hipótese de translação e a desigualdade de Hölder.
		\item \textbf{Aplicação da Dualidade}: Usar o Teorema de Representação de Riesz para garantir que $L$ é representado por uma função $g_j \in L^p(\Omega)$, identificando-a como a derivada fraca.
		\item \textbf{Prova da Minimalidade}: Utilizar a representação integral da translação e a Desigualdade de Minkowski para integrais para provar que $\|\nabla u\|_{L^p}$ é a constante ótima.
	\end{enumerate}
}

\solucao{
	Seja $\Omega \subset\subset U$ e $\phi \in C_c^\infty(\Omega)$. Para qualquer $h \neq 0$ tal que $|h| < \text{dist}(\Omega, \partial U)$, observamos que:
	$$\int_\Omega u(x) \frac{\partial \phi}{\partial x_j}(x) \, dx = \lim_{h \to 0} \int_\Omega u(x) \frac{\phi(x + h e_j) - \phi(x)}{h} \, dx$$
	Efetuando a mudança de variável $y = x + h e_j$ no primeiro termo da direita, temos:
	\begin{align*}
		\int_\Omega u(x) \frac{\partial \phi}{\partial x_j}(x) \, dx &= \lim_{h \to 0} \frac{1}{h} \left( \int_{\Omega + h e_j} u(y - h e_j) \phi(y) \, dy - \int_\Omega u(x) \phi(x) \, dx \right) \\
		&= \lim_{h \to 0} \int_\Omega \frac{(\tau_{-h e_j} u)(x) - u(x)}{h} \phi(x) \, dx
	\end{align*}
	
	Definimos agora o funcional linear $L : C_c^\infty(\Omega) \to \mathbb{R}$ por:
	$$L(\phi) = - \int_\Omega u \frac{\partial \phi}{\partial x_j} \, dx = \lim_{h \to 0} \int_\Omega \frac{u(x) - (\tau_{-h e_j} u)(x)}{h} \phi(x) \, dx$$
	Aplicando a \textbf{Desigualdade de Hölder} para o par $(p, p')$:
	\begin{align*}
		|L(\phi)| &= \lim_{h \to 0} \left| \int_\Omega \frac{u - \tau_{-h e_j} u}{h} \phi \, dx \right| \\
		&\le \lim_{h \to 0} \left\| \frac{u - \tau_{-h e_j} u}{h} \right\|_{L^p(\Omega)} \|\phi\|_{L^{p'}(\Omega)} \\
		&\le \lim_{h \to 0} \frac{1}{|h|} (C|h|) \|\phi\|_{L^{p'}(\Omega)} = C \|\phi\|_{L^{p'}(\Omega)}
	\end{align*}
	A desigualdade $|L(\phi)| \le C \|\phi\|_{L^{p'}(\Omega)}$ prova que $L$ é um funcional linear limitado sobre $C_c^\infty(\Omega)$, que é denso em $L^{p'}(\Omega)$. Pela teoria de espaços duais, como o dual de $L^{p'}$ é $L^p$, existe um único elemento $g_j \in L^p(\Omega)$ tal que:
	$$L(\phi) = \int_\Omega g_j(x) \phi(x) \, dx \implies - \int_\Omega u \frac{\partial \phi}{\partial x_j} \, dx = \int_\Omega g_j \phi \, dx$$
	Isso é precisamente a definição de que $g_j$ é a derivada fraca $\partial_j u$. Como isso vale para todo $j=1, \dots, n$ e todo $\Omega \subset\subset U$, concluímos que $u \in W^{1,p}(U)$.
	
	Para provar que $\|\nabla u\|_{L^p(U)}$ é a \textbf{menor} constante $C$, consideremos a representação integral da translação para $u \in W^{1,p}(U)$:
	$$(\tau_{he_j} u)(x) - u(x) = \int_0^h \partial_j u(x + s e_j) \, ds$$
	Aplicando a \textbf{Desigualdade de Minkowski para integrais}:
	\begin{align*}
		\|\tau_{he_j} u - u\|_{L^p(\Omega)} &= \left\| \int_0^h \partial_j u(\cdot + s e_j) \, ds \right\|_{L^p(\Omega)} \\
		&\le \int_0^h \|\partial_j u(\cdot + s e_j)\|_{L^p(\Omega)} \, ds \\
		&\le \int_0^h \|\partial_j u\|_{L^p(U)} \, ds = |h| \cdot \|\partial_j u\|_{L^p(U)}
	\end{align*}
	Como vimos que qualquer constante $C$ deve satisfazer $C \ge \|g_j\|_{L^p} = \|\partial_j u\|_{L^p}$ para que o funcional seja limitado, e provamos que $C = \|\partial_j u\|_{L^p}$ é suficiente para a desigualdade, concluímos que:
	$$\boxed{C_{min} = \|\nabla u\|_{L^p(U)}}$$
}

\solucaoalt{
	Seja $\Omega \subset\subset U$ um aberto limitado e $e_j$ o vetor da base canônica. Definimos o quociente de diferença na direção $j$ para $h \in \mathbb{R} \setminus \{0\}$:
	$$v_{h,j}(x) = \frac{u(x + h e_j) - u(x)}{h} = \frac{(\tau_{h e_j} u)(x) - u(x)}{h}$$
	Pela hipótese do enunciado, para $|h| < \text{dist}(\Omega, \partial U)$, temos:
	$$\|v_{h,j}\|_{L^p(\Omega)} = \frac{1}{|h|} \|\tau_{h e_j} u - u\|_{L^p(\Omega)} \le \frac{1}{|h|} C|h| = C$$
	Assim, para cada $j \in \{1, \dots, n\}$, a sequência $\{v_{h,j}\}_{h \neq 0}$ é limitada na norma de $L^p(\Omega)$. Como $1 < p < \infty$, o espaço $L^p(\Omega)$ é reflexivo. Pelo Teorema de Banach-Alaoglu, existe uma subsequência $h_k \to 0$ e um elemento $g_j \in L^p(\Omega)$ tal que:
	$$v_{h_k,j} \rightharpoonup g_j \quad \text{fracamente em } L^p(\Omega), \quad \text{quando } k \to \infty.$$
	
	Para provar que $g_j$ é a derivada fraca de $u$, tomemos qualquer função teste $\phi \in C_c^\infty(\Omega)$. Pela definição de $v_{h,j}$:
	$$\int_\Omega v_{h,j}(x) \phi(x) dx = \int_\Omega \frac{u(x + h e_j) - u(x)}{h} \phi(x) dx$$
	Abrindo a integral e efetuando a mudança de variável $y = x + h e_j$ no primeiro termo:
	\begin{align*}
		\int_\Omega v_{h,j}(x) \phi(x) dx &= \frac{1}{h} \int_{\Omega + h e_j} u(y) \phi(y - h e_j) dy - \frac{1}{h} \int_\Omega u(x) \phi(x) dx \\
		&= \int_\Omega u(x) \left[ \frac{\phi(x - h e_j) - \phi(x)}{h} \right] dx
	\end{align*}
	Quando $h \to 0$, a diferença finita $\frac{\phi(x - h e_j) - \phi(x)}{h}$ converge uniformemente para $-\frac{\partial \phi}{\partial x_j}(x)$ no suporte de $\phi$. Assim:
	$$\lim_{h \to 0} \int_\Omega v_{h,j}(x) \phi(x) dx = \int_\Omega u(x) \left( -\frac{\partial \phi}{\partial x_j}(x) \right) dx$$
	Por outro lado, pela convergência fraca $v_{h_k,j} \rightharpoonup g_j$:
	$$\lim_{k \to \infty} \int_\Omega v_{h_k,j}(x) \phi(x) dx = \int_\Omega g_j(x) \phi(x) dx$$
	Igualando os dois resultados, obtemos a identidade fundamental:
	$$\boxed{\int_\Omega u \frac{\partial \phi}{\partial x_j} dx = - \int_\Omega g_j \phi dx}$$
	Isso prova que $\partial_j u = g_j$. 
	Visto que a identidade fundamental é válida para qualquer $\Omega \subset\subset U$, definimos a derivada fraca global $\partial_j u = g_j$. Para analisar a norma global, consideramos uma exaustão de abertos $\{\Omega_k\}_{k\in\mathbb{N}}$ tal que $\Omega_1 \subset\subset \Omega_2 \subset\subset \dots \subset U$ com $\bigcup_{k=1}^\infty \Omega_k = U$. Pela semicontinuidade inferior da norma sob convergência fraca, temos para cada $j$ e para cada $k$:
	$$\|\partial_j u\|_{L^p(\Omega_k)} = \|g_j\|_{L^p(\Omega_k)} \le \liminf_{m \to \infty} \|v_{h_m,j}\|_{L^p(\Omega_k)} \le C$$
	Tomando o limite $k \to \infty$, a continuidade da norma $L^p$ para sequências crescentes de conjuntos garante que:
	$$\|\partial_j u\|_{L^p(U)} = \lim_{k \to \infty} \|\partial_j u\|_{L^p(\Omega_k)} \le C$$
	Isso prova que as derivadas fracas pertencem a $L^p(U)$, logo, $u \in W^{1,p}(U)$.
	
	Para provar que $\|\nabla u\|_{L^p(U)}$ é a \textbf{menor} constante $C$, devemos mostrar que a desigualdade original é satisfeita quando escolhemos $C = \|\partial_j u\|_{L^p(U)}$. Para $u \in W^{1,p}(U)$, a diferença de translação pode ser representada pela integral de sua derivada fraca:
	$$(\tau_{he_j} u)(x) - u(x) = \int_0^h \partial_j u(x + s e_j) \, ds$$
	Aplicando a \textbf{Desigualdade de Minkowski para integrais}, obtemos:
	\begin{align*}
		\|\tau_{he_j} u - u\|_{L^p(\Omega)} &= \left\| \int_0^h \partial_j u(\cdot + s e_j) \, ds \right\|_{L^p(\Omega)} \\
		&\le \int_0^h \|\partial_j u(\cdot + s e_j)\|_{L^p(\Omega)} \, ds \\
		&\le \int_0^h \|\partial_j u\|_{L^p(U)} \, ds = |h| \cdot \|\partial_j u\|_{L^p(U)}
	\end{align*}
	Esta dedução prova que a norma da derivada fraca é uma constante válida para a desigualdade e, combinada com a semicontinuidade inferior vista anteriormente, conclui-se que ela é a menor possível:
	$$\boxed{C_{min} = \|\nabla u\|_{L^p(U)}}$$
}

\sintese{
	A prova substitui a busca direta por um limite de quocientes pela análise da norma de um funcional linear. Ao provar que o funcional $\phi \mapsto - \int u \partial_j \phi$ é limitado em $L^{p'}$, a estrutura geométrica do espaço $L^p$ (reflexividade e dualidade) garante a existência da derivada fraca.
	
	O fluxo lógico foi:
	$$\text{Translação} \to \text{Funcional Limitado em } L^{p'} \to \text{Representação de Riesz } (L^{p'})^{'} \cong L^p \to \text{Derivada Fraca} \to \text{Minkowski}$$
}

\aprofundamento{
	A elegância desta prova reside no fato de que não precisamos de regulizadores para "adivinhar" a derivada; a própria estrutura do espaço $L^p$ a fornece. A restrição $1 < p < \infty$ é vital: se $p=1$, o dual de $L^\infty$ não é $L^1$ (contém medidas singulares), o que explica por que a regularidade de Sobolev falha em $p=1$ e somos forçados a usar o espaço $BV(\Omega)$.
}

%%%% QUESTÃO 107 %%%%
\questao{
	Prove que $W^{1,p}(J), J = (0, 1) \subset \mathbb{R}, 1 < p < \infty$, está imerso compactamente em $C(\overline{J})$ (use o teorema de Arzelá-Àscoli).
}

\analise{
	\textbf{Análise do Enunciado}: A tese é a prova de que a imersão $W^{1,p}(J) \hookrightarrow C(\overline{J})$ é compacta. Isso exige mostrar que qualquer conjunto limitado em $W^{1,p}(J)$ é relativamente compacto em $C(\overline{J})$, ou seja, que qualquer sequência $\{u_n\}$ tal que $\|u_n\|_{W^{1,p}} \le M$ admite uma subsequência que converge uniformemente para algum $u \in C(\overline{J})$.
	
	\textbf{Fundamentação Teórica}: O pilar desta prova é o \textbf{Teorema de Arzelá-Ascoli}, que estabelece que um subconjunto de $C(\overline{J})$ é relativamente compacto se, e somente se, for uniformemente limitado e equicontínuo. Utilizaremos a estimativa de Hölder rigorosamente deduzida na \textbf{Questão 105} para garantir a equicontinuidade e a técnica de ancoragem via média integral para a limitação uniforme.
	
	\textbf{Estratégias}:
	\begin{enumerate}
		\item Considerar um conjunto $\mathcal{B}$ limitado em $W^{1,p}(J)$, tal que $\sup_{u \in \mathcal{B}} \|u\|_{W^{1,p}} \le M$.
		\item Demonstrar a \textbf{Limitação Uniforme}: Definir a média $\overline{u}$, utilizar o Teorema do Valor Intermediário para encontrar um ponto de ancoragem $c$ e estimar $\|u\|_\infty$ através da norma de Sobolev.
		\item Demonstrar a \textbf{Equicontinuidade}: Aplicar a estimativa $|u(x) - u(y)| \le \|u'\|_{L^p} |x - y|^{1 - \frac{1}{p}}$ e definir explicitamente o $\delta(\varepsilon)$ dependente apenas de $M$ e $\varepsilon$.
		\item Aplicar o Teorema de Arzelá-Ascoli para concluir a compacidade da imersão.
	\end{enumerate}
}

\solucao{
	Seja $\mathcal{B}$ um conjunto limitado em $W^{1,p}(J)$. Existe $M > 0$ tal que $\|u\|_{W^{1,p}} \le M$ para todo $u \in \mathcal{B}$. Lembremos que $\|u\|_{W^{1,p}}^p = \|u\|_{L^p}^p + \|u'\|_{L^p}^p$, o que implica que tanto $\|u\|_{L^p} \le M$ quanto $\|u'\|_{L^p} \le M$.
	
	\textbf{1. Prova da Limitação Uniforme}:
	Seja $\overline{u} = \frac{1}{|J|} \int_J u(s) \, ds$ a média dos valores de $u$ em $\overline{J}$. Pela desigualdade de Hölder:
	$$|\overline{u}| \le \frac{1}{|J|} \int_J |u(s)| \cdot 1 \, ds \le \frac{1}{|J|} \|u\|_{L^p(J)} \|1\|_{L^{p'}(J)} = \|u\|_{L^p(J)} \le M$$
	Como $u$ é contínua em $\overline{J}$, pelo \textbf{Teorema do Valor Intermediário}, existe um ponto $c \in \overline{J}$ tal que $u(c) = \overline{u}$. Para qualquer $x \in \overline{J}$, utilizamos a representação integral e a estimativa de Hölder da \textbf{Questão 105}:
	\begin{align*}
		|u(x)| &= |u(c) + \int_c^x u'(t) \, dt| \\
		&\le |u(c)| + \|u'\|_{L^p(J)} |x - c|^{1 - \frac{1}{p}} \\
		&\le M + M \cdot (1)^{1 - \frac{1}{p}} = 2M
	\end{align*}
	Logo, $\sup_{u \in \mathcal{B}} \|u\|_\infty \le 2M$, provando que o conjunto é \textbf{uniformemente limitado}.
	
	\textbf{2. Prova da Equicontinuidade}:
	Dada a Questão 105, sabemos que para quaisquer $x, y \in J$ e para todo $u \in \mathcal{B}$:
	$$|u(x) - u(y)| \le \|u'\|_{L^p(J)} |x - y|^{1 - \frac{1}{p}} \le M |x - y|^{1 - \frac{1}{p}}$$
	Seja $\varepsilon > 0$. Definimos $\delta = \left( \frac{\varepsilon}{M} \right)^{\frac{1}{1 - 1/p}}$. Assim, para todo $u \in \mathcal{B}$, se $|x - y| < \delta$, segue que:
	$$|u(x) - u(y)| \le M \left( \left( \frac{\varepsilon}{M} \right)^{\frac{1}{1 - 1/p}} \right)^{1 - \frac{1}{p}} = M \left( \frac{\varepsilon}{M} \right) = \varepsilon$$
	Visto que $\delta$ depende exclusivamente de $\varepsilon$ e $M$, o conjunto $\mathcal{B}$ é \textbf{equicontínuo}.
	
	\textbf{3. Conclusão}:
	O conjunto $\mathcal{B}$ é um subconjunto de $C(\overline{J})$ que é uniformemente limitado e equicontínuo. Pelo \textbf{Teorema de Arzelá-Ascoli}, $\mathcal{B}$ é relativamente compacto em $C(\overline{J})$. Isso implica que qualquer sequência $\{u_n\} \subset \mathcal{B}$ admite uma subsequência que converge uniformemente para alguma função $u \in C(\overline{J})$. Logo, a imersão $W^{1,p}(J) \hookrightarrow C(\overline{J})$ é compacta.
}

\sintese{
	A prova revela a \textbf{troca} de regularidade: a integrabilidade da derivada ($\|u'\|_{L^p}$) é convertida em controle pontual da variação da função e em limitação da amplitude. A compacidade surge porque a norma de Sobolev impede oscilações violentas e a fuga para o infinito.
	
	O fluxo lógico foi:
	$$\boxed{\text{Média } \overline{u} \to \text{TVI (Ponto } c\text{)} \to \text{Lim. Uniforme} \to \text{Eqüicontinuidade} \to \text{Arzelá-Ascoli}}$$
}

\aprofundamento{
	Este resultado é um caso particular do \textbf{Teorema de Rellich-Kondrachov}. Em dimensões maiores ($n > 1$), a imersão $W^{1,p} \hookrightarrow L^q$ é compacta para $q < p^* = \tfrac{np}{n-p}$. No caso $n=1$, como $p > 1$, temos $p^* = \infty$, o que explica por que a imersão ocorre no espaço de funções contínuas $C(\overline{J})$. A compacidade é a ferramenta fundamental para a existência de soluções em EDPs, pois permite extrair limites de sequências de aproximações.
}

%%%% QUESTÃO 108 %%%%
\questao{
	Seja $G : \mathbb{R} \to \mathbb{R}$ uma função de classe $C^1$ tal que $G(0) = 0, \ |G'(t)| \le M$, para todo $t \in \mathbb{R}$. Prove que $v(x) = G \circ u(x)$ é uma função em $W^{1,p}(U)$, sempre que $u \in W^{1,p}(U)$. Mais ainda
	$$\frac{\partial v}{\partial x_j} = G'(u) \frac{\partial u}{\partial x_j}.$$
}

\analise{
	\textbf{Análise do Enunciado}: A questão solicita a prova de que a composição de uma função de Sobolev $u$ com uma função $C^1$ com derivada limitada preserva a pertinência ao espaço $W^{1,p}$. A tese divide-se em duas partes: a prova de que $v \in L^p(U)$ e a prova de que a derivada fraca de $v$ existe e é dada pela regra da cadeia clássica.
	
	\textbf{Fundamentação Teórica}: A prova baseia-se na teoria de composição em espaços de Sobolev. O fato de $G$ possuir derivada limitada implica que $G$ é uma função \textbf{Lipschitziana}. Funções Lipschitz preservam a regularidade de Sobolev. Utilizaremos a técnica de regularização via regularização para justificar a passagem ao limite.
	
	\textbf{Estratégias}:
	\begin{enumerate}
		\item \textbf{Integrabilidade de $v$}: Utilizar a condição $G(0)=0$ e $|G'| \le M$ para provar que $|G(t)| \le M|t|$, garantindo que $v \in L^p(U)$.
		\item \textbf{Regularização}: Introduzir uma sequência de funções suaves $u_k = u * \eta_k$ que converge para $u$ em $W^{1,p}(U)$.
		\item \textbf{Regra da Cadeia Clássica}: Definir $v_k = G(u_k)$, que por ser a composição de funções $C^\infty$ e $C^1$, é derivável no sentido clássico, satisfazendo $\partial_j v_k = G'(u_k) \partial_j u_k$.
		\item \textbf{Convergência em $L^p$}: Provar que $v_k \to v$ e $\nabla v_k \to G'(u) \nabla u$ em $L^p(U)$, utilizando a continuidade de $G'$ e a convergência de $u_k$.
		\item \textbf{Identificação da Derivada Fraca}: Passar ao limite na definição de derivada fraca para concluir que $v \in W^{1,p}(U)$.
	\end{enumerate}
}

\solucao{
	Primeiramente, provamos que $v \in L^p(U)$. Como $G$ é de classe $C^1$ e $G(0)=0$, pelo Teorema do Valor Médio, para qualquer $t \in \mathbb{R}$, existe $\xi$ entre $0$ e $t$ tal que $G(t) - G(0) = G'(\xi)(t - 0)$. Logo:
	$$|G(t)| = |G'(\xi)| |t| \le M|t|$$
	Aplicando isso a $v(x) = G(u(x))$, temos:
	$$\int_U |v(x)|^p \, dx = \int_U |G(u(x))|^p \, dx \le \int_U (M|u(x)|)^p \, dx = M^p \|u\|_{L^p(U)}^p < \infty$$
	Portanto, $v \in L^p(U)$.
	
	Para a derivada, seja $u_k = u * \eta_k$ a regularizante de $u$ em $U$. Sabemos que $u_k \in C^\infty(U)$ e $u_k \to u$ em $W^{1,p}(U)$. Definimos $v_k = G(u_k)$. Como $u_k$ é suave e $G \in C^1$, $v_k$ é derivável no sentido clássico e:
	$$\frac{\partial v_k}{\partial x_j}(x) = G'(u_k(x)) \frac{\partial u_k}{\partial x_j}(x)$$
	Analisamos a convergência de $v_k$ para $v$ em $L^p(U)$. Pela propriedade Lipschitz de $G$:
	$$\|v_k - v\|_{L^p(U)} = \|G(u_k) - G(u)\|_{L^p(U)} \le M \|u_k - u\|_{L^p(U)} \to 0$$
	Agora, analisamos a convergência de $\nabla v_k$. Consideremos o termo:
	$$\| \nabla v_k - G'(u) \nabla u \|_{L^p(U)} = \| G'(u_k) \nabla u_k - G'(u) \nabla u \|_{L^p(U)}$$
	Adicionamos e subtraímos o termo $G'(u_k) \nabla u$:
	\begin{align*}
		\| \nabla v_k - G'(u) \nabla u \|_{L^p(U)} &\le \| G'(u_k) \nabla u_k - G'(u_k) \nabla u \|_{L^p(U)} + \| G'(u_k) \nabla u - G'(u) \nabla u \|_{L^p(U)} \\
		&\le \| G'(u_k) \|_{L^\infty} \| \nabla u_k - \nabla u \|_{L^p(U)} + \| G'(u_k) - G'(u) \|_{L^p(U, |\nabla u|^p)}
	\end{align*}
	No primeiro termo, $\|G'(u_k)\|_{L^\infty} \le M$. No segundo termo, como $G'$ é contínua e $u_k \to u$ em $L^p$, segue que $G'(u_k) \to G'(u)$ em medida e, por estar limitado por $M$, converge também em $L^p$ (Teorema da Convergência Dominada). Assim:
	$$\| \nabla v_k - G'(u) \nabla u \|_{L^p(U)} \to 0 \quad \text{quando } k \to \infty$$
	
	Finalmente, tomemos qualquer $\phi \in C_c^\infty(U)$. Como $v_k$ é suave, temos:
	$$\int_U v_k \frac{\partial \phi}{\partial x_j} \, dx = - \int_U \frac{\partial v_k}{\partial x_j} \phi \, dx = - \int_U G'(u_k) \frac{\partial u_k}{\partial x_j} \phi \, dx$$
	Passando ao limite $k \to \infty$, e utilizando a convergência forte de $v_k$ e $\nabla v_k$ em $L^p$:
	$$\int_U v \frac{\partial \phi}{\partial x_j} \, dx = - \int_U G'(u) \frac{\partial u}{\partial x_j} \phi \, dx$$
	Esta identidade prova que $v \in W^{1,p}(U)$ e que sua derivada fraca é:
	$$\boxed{\frac{\partial v}{\partial x_j} = G'(u) \frac{\partial u}{\partial x_j}}$$
}

\sintese{
	A prova demonstra que a regularidade de Sobolev é preservada sob composições com funções Lipschitz. O contorno analítico ocorre na regularização: transformamos a função de Sobolev em uma sequência de funções suaves onde a regra da cadeia é trivial, e provamos que a estrutura da derivada é preservada no limite $L^p$.
	
	O fluxo lógico foi:
	$$\boxed{\text{Lipschitz de } G \to \text{Regularização} \to \text{ Regra da Cadeia Clássica} \to \text{TCDL em } L^p \to \text{Derivada Fraca}}$$
}

\aprofundamento{
	Um resultado mais geral é o \textbf{Teorema de Stampacchia}, que afirma que se $u \in W^{1,p}(U)$ e $G$ é Lipschitz, então $G \circ u \in W^{1,p}(U)$. No caso onde $G$ não é $C^1$, mas apenas Lipschitz, a derivada $G'$ existe quase em toda parte (\textbf{Teorema de Rademacher}), e a fórmula da regra da cadeia ainda é válida quase em toda parte. Isso é fundamental para tratar funções como $u^+ = \max(u, 0)$, onde $G(t) = \max(t, 0)$ é Lipschitz mas não $C^1$.
}

%%%% QUESTÃO 109 %%%%
\questao{
	Sejam $c \in \mathbb{R}$ e $G: \mathbb{R} \to \mathbb{R}$ uma função Lipschitziana, derivável em $\mathbb{R} \setminus \{c\}$ tal que $G(0) = 0$. Prove que $v(x) = G \circ u(x)$ é uma função em $W^{1,p}(U)$, sempre que $u \in W^{1,p}(U)$. Mais ainda,
	$$ \frac{\partial v}{\partial x_j} = G'(u)\frac{\partial u}{\partial x_j} \text{ quando } u(x) \neq c \quad \text{e} \quad \frac{\partial v}{\partial x_j} = 0 \text{ em } u(x) = c $$
	Para isto basta verificar que:
	\begin{enumerate}[label=(\alph*)]
		\item $|G(t)| \le M|t|$ e $|G'(t)| \le M$, para todo $t \in \mathbb{R} \setminus \{c\}$, onde $M$ é a constante de Lipschitz para a função $G$;
		\item Existe uma família $G_m: \mathbb{R} \to \mathbb{R}$ de funções deriváveis tais que $|G_m(t)| \le 2M|t|$ e $|G'_m(t)| \le 2M$, para todo $t \in \mathbb{R}$, com $G_m(0) = 0$, $G_m(t) \to G(t)$, para todo $t \in \mathbb{R}$ e $G'_m(t) \to G'(t)$, para todo $t \in \mathbb{R} \setminus \{c\}$;
		\item $\nabla u = 0$ quase sempre sobre o conjunto $U_c = \{x \in U : u(x) = c\}$;
		\item Use o teorema da convergência dominada para concluir;
		\item Para construir a família $G_m$, sem perda de generalidade suponha que $c = 0$ e $G(0) = 0$. Defina em $\left[0, \tfrac{1}{m}\right]$: $g_m(t) = a_m t^2 + b_m t^3$, com $a_m$ e $b_m$ tais que: $g_m\left(\tfrac{1}{m}\right) = G\left(\tfrac{1}{m}\right)$ e $g'_m\left(\tfrac{1}{m}\right) = G'\left(\tfrac{1}{m}\right)$. Mostre que $g_m(0) = g'_m(0) = 0$, $|g_m(t)| \le \tfrac{7M}{m}$ e $|g'_m(t)| \le 17M$. Em $\left[-\tfrac{1}{m}, 0\right]$: $h_m(t) = \tilde{a}_m t^2 + \tilde{b}_m t^3$, com $\tilde{a}_m$ e $\tilde{b}_m$ tais que: $h_m\left(-\tfrac{1}{m}\right) = G\left(-\tfrac{1}{m}\right)$ e $h'_m\left(-\tfrac{1}{m}\right) = G'\left(-\tfrac{1}{m}\right)$. Finalmente, defina
		$$ G_m(t) = \begin{cases} 
			G(t), & \text{em } |t| \ge \tfrac{1}{m} \\ 
			h_m(t), & \text{em } \left[-\tfrac{1}{m}, 0\right] \\ 
			g_m(t), & \text{em } \left[0, \tfrac{1}{m}\right] \end{cases} $$
		mostre que $|G_m(t) - G(t)| \le \tfrac{8M}{m}$ e $|G'_m(t)| \le 2M$, para todo $t \in \mathbb{R}$ para $m$ grande. Veja que as desigualdades do item (b) seguem naturalmente.
\end{enumerate}}

\analise{
	\textbf{Análise do Enunciado}: A questão propõe a extensão da regra da cadeia para funções de Sobolev quando a função externa $G$ é apenas Lipschitz. O ponto crítico é a possível não-derivabilidade de $G$ no ponto $c$. A tese exige provar que a função composta $v$ permanece no espaço $W^{1,p}(U)$ e que sua derivada fraca é nula onde $u$ assume o valor crítico $c$.
	
	\textbf{Fundamentação Teórica}: A demonstração mobiliza a teoria de funções Lipschitz e a propriedade de anulação do gradiente em conjuntos de nível, fundamental nos espaços de Sobolev. A convergência final é estabelecida via Teorema da Convergência Dominada de Lebesgue, utilizando a sequência de aproximações suaves $G_m$ para evitar a singularidade em $c$.
	
	\textbf{Estratégias}:
	\begin{enumerate}
		\item \textbf{Análise de Cotas}: Estabelecer as limitações de $G$ e $G'$ a partir da constante de Lipschitz $M$.
		\item \textbf{Construção Algébrica de $G_m$}: Resolver o sistema linear para determinar os coeficientes dos polinômios de colagem e provar que as derivadas permanecem limitadas.
		\item \textbf{Análise de Medida}: Justificar por que $\nabla u = 0$ q.s. em $U_c$.
		\item \textbf{Passagem ao Limite}: Aplicar a regra da cadeia para $v_m = G_m \circ u$ e usar o TCD para recuperar a derivada de $v$.
	\end{enumerate}
}

\solucao{
	Procedemos à demonstração seguindo o roteiro proposto, iniciando pela base algébrica da suavização:
	
	\textbf{(e) Construção e Cotas de $G_m$}:
	Sem perda de generalidade, assumimos $c = 0$ e $G(0) = 0$. Para $t \in \left[ 0, \frac{1}{m} \right]$, definimos $g_m(t) = a_m t^2 + b_m t^3$. As condições de colagem em $t = \frac{1}{m}$ impõem o sistema:
	\begin{align*}
		g_m\left(\frac{1}{m}\right) &= \frac{a_m}{m^2} + \frac{b_m}{m^3} = G\left(\frac{1}{m}\right) \\
		g_m'\left(\frac{1}{m}\right) &= \frac{2a_m}{m} + \frac{3b_m}{m^2} = G'\left(\frac{1}{m}\right)
	\end{align*}
	Multiplicando a primeira equação por $3m$, obtemos $\frac{3a_m}{m} + \frac{3b_m}{m^2} = 3mG\left(\frac{1}{m}\right)$. Subtraindo a segunda equação, isolamos o coeficiente quadrático:
	$$ \frac{a_m}{m} = 3mG\left(\frac{1}{m}\right) - G'\left(\frac{1}{m}\right) \implies a_m = 3m^2G\left(\frac{1}{m}\right) - mG'\left(\frac{1}{m}\right) $$
	Substituindo $a_m$ na equação de continuidade:
	$$ \frac{b_m}{m^3} = G\left(\frac{1}{m}\right) - \left( \frac{3m^2G\left(\frac{1}{m}\right) - mG'\left(\frac{1}{m}\right)}{m^2} \right) = -2G\left(\frac{1}{m}\right) + \frac{1}{m}G'\left(\frac{1}{m}\right) $$
	Portanto, $b_m = -2m^3G\left(\frac{1}{m}\right) + m^2G'\left(\frac{1}{m}\right)$. Notamos que $g_m(0) = 0$ e $g_m'(0) = 0$. Usando $|G\left(\frac{1}{m}\right)| \le \frac{M}{m}$ e $|G'\left(\frac{1}{m}\right)| \le M$, estimamos os coeficientes:
	$$ |a_m| \le 3m^2\left(\frac{M}{m}\right) + m(M) = 4mM \quad \text{e} \quad |b_m| \le 2m^3\left(\frac{M}{m}\right) + m^2(M) = 3m^2M $$
	Para $t \in \left[ 0, \frac{1}{m} \right]$, temos:
	$$ |g_m(t)| \le |a_m|t^2 + |b_m|t^3 \le (4mM)\left(\frac{1}{m^2}\right) + (3m^2M)\left(\frac{1}{m^3}\right) = \frac{7M}{m} $$
	A derivada é $g_m'(t) = 2a_m t + 3b_m t^2$. Assim:
	$$ |g_m'(t)| \le 2(4mM)\left(\frac{1}{m}\right) + 3(3m^2M)\left(\frac{1}{m^2}\right) = 8M + 9M = 17M $$
	A construção de $h_m$ em $\left[ -\frac{1}{m}, 0 \right]$ é análoga. Para $|t| < \frac{1}{m}$:
	$$ |G_m(t) - G(t)| \le |g_m(t)| + |G(t)| \le \frac{7M}{m} + \frac{M}{m} = \frac{8M}{m} $$
	Para $m$ suficientemente grande, as cotas $|G_m(t)| \le 2M|t|$ e $|G_m'(t)| \le 2M$ são satisfeitas.
	
	\textbf{(a) Estimativas de $G$}:
	Como $G$ é Lipschitziana com constante $M$, temos $|G(t) - G(s)| \le M|t - s|$. Fixando $s=0$ e $G(0)=0$, obtemos $|G(t)| \le M|t|$. Nos pontos onde $G$ é derivável, a magnitude da derivada é limitada por $M$ pois:
	$$ |G'(t)| = \lim_{h \to 0} \frac{|G(t+h) - G(t)|}{|h|} \le \lim_{h \to 0} \frac{M|h|}{|h|} = M $$
	
	\textbf{(b) Convergência de $G_m$}:
	A construção anterior garante que $G_m \to G$ uniformemente em $\mathbb{R}$ e $G_m' \to G'$ pontualmente em $\mathbb{R} \setminus \{c\}$, validando as premissas do item (b).
	
	\textbf{(c) Gradiente no Conjunto de Nível}:
	Seja $U_c = \{x \in U : u(x) = c\}$. Por propriedade fundamental dos espaços de Sobolev, se uma função $u \in W^{1,p}(U)$ é constante em um conjunto mensurável, sua derivada fraca anula-se quase sempre nesse conjunto. Portanto:
	$$ \nabla u = 0 \quad \text{quase sempre em } U_c $$
	
	\textbf{(d) Conclusão via TCD}:
	Definimos $v_m(x) = G_m(u(x))$. Como $G_m \in C^1$, temos $\frac{\partial v_m}{\partial x_j} = G_m'(u) \frac{\partial u}{\partial x_j}$. Analisamos o limite $m \to \infty$:
	\begin{itemize}
		\item Se $u(x) \neq c$, então $G_m'(u(x)) \to G'(u(x))$.
		\item Se $u(x) = c$, então $G_m'(u(x)) \frac{\partial u}{\partial x_j} \to G_m'(c) \cdot 0 = 0$.
	\end{itemize}
	Logo, $\frac{\partial v_m}{\partial x_j} \to G'(u) \frac{\partial u}{\partial x_j} \chi_{\{u \neq c\}}$ quase sempre. Como $|G_m'(u) \frac{\partial u}{\partial x_j}| \le 17M |\frac{\partial u}{\partial x_j}| \in L^p(U)$, o \textbf{Teorema da Convergência Dominada} garante a convergência forte em $L^p(U)$. Concluímos que $v \in W^{1,p}(U)$ e:
	$$\boxed{\frac{\partial v}{\partial x_j} = \begin{cases} G'(u) \frac{\partial u}{\partial x_j}, & \text{se } u(x) \neq c \\ 0, & \text{se } u(x) = c \end{cases}}$$
}

\sintese{
	A demonstração consolida o princípio de que transformações Lipschitzianas preservam a regularidade em espaços de Sobolev, estendendo a validade da diferenciação fraca através do anulamento geométrico do gradiente em patamares de descontinuidade da derivada.
	
	O fluxo lógico foi:
	$$\boxed{\text{Interpolação de } G_m \implies \text{Estimativas Uniformes de } G_m' \implies \text{Anulamento de } \nabla u \text{ em } U_c \implies \text{TCD } L^p(U).}$$
}

\aprofundamento{
	Este resultado provê a infraestrutura teórica necessária para o emprego de métodos variacionais em Equações Diferenciais Parciais não lineares. A possibilidade de compor funções de Sobolev com mapeamentos Lipschitzianos cujas derivadas possuem descontinuidades isoladas fundamenta a aplicação de operadores de truncamento, como a função parte positiva $u^+ = \max\{u, 0\}$. A propriedade $\nabla u = 0$ q.s. em $\{u=c\}$ é a razão pela qual a derivada de $u^+$ é simplesmente $\nabla u \cdot \chi_{\{u > 0\}}$, permitindo a análise de subsoluções elípticas.
}

%%%% QUESTÃO 110 %%%%
\questao{
	Seja $u \in W^{1,p}(\mathbb{R}^n)$. Defina
	$$w(x) = \begin{cases} u^3, & \text{se } |u| \le 1 \\ u, & \text{se } |u| > 1. \end{cases}$$
	Mostre que $w \in W^{1,p}(\mathbb{R}^n)$, e calcule $\int_{\mathbb{R}^n} \nabla w \cdot \nabla u \, dx$.
}

\analise{
	\textbf{Análise do Enunciado}: A questão solicita a prova de que a função $w$, definida por uma colagem de potências de $u$, pertence ao espaço de Sobolev $W^{1,p}(\mathbb{R}^n)$. Em seguida, pede-se o cálculo de uma integral que representa o produto interno dos gradientes de $w$ e $u$. A tese repousa na identificação de $w$ como a composição $G \circ u$, onde $G$ é uma função definida por partes.
	
	\textbf{Fundamentação Teórica}: Utilizaremos o resultado da Questão 109, que estabelece que se $u \in W^{1,p}(U)$ e $G: \mathbb{R} \to \mathbb{R}$ é uma função Lipschitziana com $G(0)=0$, então $G \circ u \in W^{1,p}(U)$. Para isso, devemos provar que a função $G$ definida no problema é globalmente Lipschitziana.
	
	\textbf{Estratégias}:
	\begin{enumerate}
		\item \textbf{Definição da Função Externa}: Isolar a função $G(t)$ tal que $w = G(u)$.
		\item \textbf{Prova de Lipschitz}: Verificar a continuidade de $G$ nos pontos de colagem e a limitação de sua derivada em cada ramo para determinar a constante de Lipschitz $M$.
		\item \textbf{Aplicação da Regra da Cadeia}: Utilizar a fórmula da derivada fraca para composições Lipschitzianas para expressar $\nabla w$.
		\item \textbf{Decomposição da Integral}: Particionar o domínio $\mathbb{R}^n$ nos conjuntos $\{|u| \le 1\}$ e $\{|u| > 1\}$ para calcular a integral solicitada.
	\end{enumerate}
}

\solucao{
	Definimos a função externa $G: \mathbb{R} \to \mathbb{R}$ por
	$$G(t) = \begin{cases} t^3, & \text{se } |t| \le 1 \\ t, & \text{se } |t| > 1 \end{cases}$$
	de modo que $w(x) = (G \circ u)(x)$. Para provar que $w \in W^{1,p}(\mathbb{R}^n)$, devemos primeiro garantir que $G$ seja globalmente Lipschitziana. 
	
	\textbf{1. Prova da Continuidade de $G$}:
	A função $G$ é claramente contínua nos intervalos abertos $(-1, 1)$ e $\mathbb{R} \setminus [-1, 1]$. Analisamos a continuidade nos pontos de transição (colagem) $t = 1$ e $t = -1$:
	\begin{itemize}
		\item Para $t = 1$:
		$$\lim_{t \to 1^-} G(t) = \lim_{t \to 1^-} (t^3) = 1^3 = 1 \quad \text{e} \quad \lim_{t \to 1^+} G(t) = \lim_{t \to 1^+} (t) = 1$$
		\item Para $t = -1$:
		$$\lim_{t \to -1^-} G(t) = \lim_{t \to -1^-} (t) = -1 \quad \text{e} \quad \lim_{t \to -1^+} G(t) = \lim_{t \to -1^+} (t^3) = (-1)^3 = -1$$
	\end{itemize}
	Visto que os limites laterais coincidem com os valores da função em ambos os pontos, $G$ é contínua em todo $\mathbb{R}$.
	
	\textbf{2. Propriedade Lipschitziana}:
	A função $G$ é derivável em $\mathbb{R} \setminus \{-1, 1\}$, com a seguinte derivada:
	$$G'(t) = \begin{cases} 3t^2, & \text{se } |t| < 1 \\ 1, & \text{se } |t| > 1 \end{cases}$$
	Analisamos a magnitude de $G'(t)$:
	\begin{itemize}
		\item Se $|t| < 1$, então $0 \le 3t^2 < 3$.
		\item Se $|t| > 1$, então $|G'(t)| = 1$.
	\end{itemize}
	Assim, $|G'(t)| \le 3$ para quase todo $t \in \mathbb{R}$. Pelo teorema que relaciona a limitação da derivada com a constante de Lipschitz para funções contínuas, $G$ é \textbf{Lipschitziana} com constante $M = 3$. Além disso, nota-se que $G(0) = 0$.
	
	\textbf{3. Pertencimento ao Espaço de Sobolev}:
	Pelo resultado estabelecido na Questão 109, como $u \in W^{1,p}(\mathbb{R}^n)$ e $G$ é uma função Lipschitziana com $G(0)=0$, a composição $w = G \circ u$ pertence necessariamente ao espaço $W^{1,p}(\mathbb{R}^n)$. A derivada fraca de $w$ é dada pela regra da cadeia para funções Lipschitz:
	$$\nabla w(x) = G'(u(x)) \nabla u(x)$$
	Substituindo a expressão de $G'(u)$, obtemos a caracterização do gradiente:
	$$\nabla w(x) = \begin{cases} 3u^2(x) \nabla u(x), & \text{se } |u(x)| \le 1 \\ \nabla u(x), & \text{se } |u(x)| > 1 \end{cases}$$
	
	\textbf{4. Cálculo da Integral de Energia}:
	Para calcular a integral $\int_{\mathbb{R}^n} \nabla w \cdot \nabla u \, dx$, decompomos o domínio $\mathbb{R}^n$ nos conjuntos mensuráveis $A = \{x \in \mathbb{R}^n : |u(x)| \le 1\}$ e $B = \{x \in \mathbb{R}^n : |u(x)| > 1\}$. Temos:
	\begin{align*}
		\int_{\mathbb{R}^n} \nabla w \cdot \nabla u \, dx &= \int_A \nabla w \cdot \nabla u \, dx + \int_B \nabla w \cdot \nabla u \, dx \\
		&= \int_A \left( 3u^2 \nabla u \right) \cdot \nabla u \, dx + \int_B (\nabla u) \cdot \nabla u \, dx \\
		&= \int_{\{|u| \le 1\}} 3u^2 |\nabla u|^2 \, dx + \int_{\{|u| > 1\}} |\nabla u|^2 \, dx
	\end{align*}
	Utilizando a função característica $\chi$, podemos expressar o resultado de forma compacta e rigorosa:
	$$\boxed{\int_{\mathbb{R}^n} \nabla w \cdot \nabla u \, dx = \int_{\mathbb{R}^n} \left( 3u^2 \chi_{\{|u| \le 1\}} + \chi_{\{|u| > 1\}} \right) |\nabla u|^2 \, dx}$$
}

\sintese{
	A prova demonstra que a regularidade de Sobolev é robusta sob a composição com funções Lipschitzianas, mesmo que estas apresentem descontinuidades na derivada. O cálculo da integral revela que a "energia" do gradiente de $w$ é uma versão ponderada da energia de $u$, onde o peso é determinado pela magnitude de $u$.
	
	O fluxo lógico foi:
	$$\boxed{\text{Definição de } G \to \text{Prova de Lipschitz} \to \text{Regra da Cadeia Sobolev} \to \text{Decomposição de Domínio}}$$
}

\aprofundamento{
	A função $G(t)$ utilizada nesta questão é um exemplo de função de "truncamento" ou "saturação". Em problemas de EDPs não lineares, como na equação de p-Laplaciano ou em modelos de plasticidade, é comum encontrar operadores que alteram a difusividade da função dependendo de sua amplitude. O fato de $w \in W^{1,p}$ garante que podemos aplicar métodos variacionais a este novo funcional, assegurando que a \textbf{quebra} de regularidade de $G'$ em $t=\pm 1$ não prejudique a integrabilidade global do gradiente.
}

%%%% QUESTÃO 111 %%%%
\questao{
	Seja $u \in W^{1,p}(U)$.
	\begin{enumerate}[label=(\alph*)]
		\item Prove que $v(x) = \frac{u(x)^3}{u(x)^2 + 1}$ é uma função em $W^{1,p}(U)$. Quem é a derivada fraca $\frac{\partial v}{\partial x_j}$?
		\item Faça o mesmo para a função $w(x) = 1 - \cos u(x)$. Mostre que é uma função em $W^{1,p}(U)$ e determine a derivada fraca $\frac{\partial w}{\partial x_j}$.
	\end{enumerate}
}

\analise{
	\textbf{Análise do Enunciado}: A questão solicita a prova de que duas funções compostas, $v = G(u)$ e $w = H(u)$, pertencem ao espaço de Sobolev $W^{1,p}(U)$, onde $u$ já é conhecido pertencer a esse espaço. A tese exige a identificação explícita das derivadas fracas de cada composição.
	
	\textbf{Fundamentação Teórica}: Utilizaremos o resultado estabelecido na Questão 108, que afirma: se $u \in W^{1,p}(U)$ e $G: \mathbb{R} \to \mathbb{R}$ é uma função de classe $C^1$ com derivada limitada ($|G'| \le M$) e $G(0)=0$, então $G \circ u \in W^{1,p}(U)$ e $\partial_j (G \circ u) = G'(u) \partial_j u$. A verificação de que $G(0)=0$ é essencial para garantir que a função resultante pertença a $L^p(U)$, especialmente se o domínio $U$ for ilimitado.
	
	\textbf{Estratégias}:
	\begin{enumerate}
		\item \textbf{Análise de $G(t)$}: Para o item (a), definir $G(t) = \frac{t^3}{t^2+1}$, provar que $G(0)=0$ e demonstrar que $G'$ é limitada em $\mathbb{R}$.
		\item \textbf{Cálculo da Derivada Fraca de $v$}: Aplicar a regra da cadeia para funções de Sobolev utilizando o $G'(t)$ calculado.
		\item \textbf{Análise de $H(t)$}: Para o item (b), definir $H(t) = 1 - \cos t$, provar que $H(0)=0$ e que $|H'| \le 1$.
		\item \textbf{Cálculo da Derivada Fraca de $w$}: Aplicar a regra da cadeia para funções de Sobolev utilizando o $H'(t)$ calculado.
	\end{enumerate}
}

\solucao{
	\textbf{(a) Análise da função $v(x)$}:
	Definimos a função externa $G: \mathbb{R} \to \mathbb{R}$ por $G(t) = \frac{t^3}{t^2 + 1}$. Note que $G(0) = \frac{0}{1} = 0$. Como $G$ é a razão de dois polinômios onde o denominador nunca se anula ($\forall t \in \mathbb{R}, t^2+1 \ge 1$), $G$ é de classe $C^\infty$. Calculamos sua derivada via regra do quociente:
	\begin{align*}
		G'(t) &= \frac{(3t^2)(t^2 + 1) - (t^3)(2t)}{(t^2 + 1)^2} \\
		&= \frac{3t^4 + 3t^2 - 2t^4}{(t^2 + 1)^2} = \frac{t^4 + 3t^2}{(t^2 + 1)^2}
	\end{align*}
	Para provar que $G$ é Lipschitziana, devemos mostrar que $G'(t)$ é limitada. Analisamos o comportamento de $G'(t)$:
	\begin{itemize}
		\item Para $t \to \pm \infty$, temos $\lim_{t \to \pm \infty} \frac{t^4 + 3t^2}{t^4 + 2t^2 + 1} = 1$.
		\item Como $G'$ é contínua em $\mathbb{R}$ e possui limites finitos no infinito, ela é necessariamente limitada. 
	\end{itemize}
	De acordo com o resultado da Questão 108, como $u \in W^{1,p}(U)$, $G(0)=0$ e $|G'| \le M$, a composição $v = G \circ u$ pertence a $W^{1,p}(U)$. A derivada fraca é dada por:
	$$\boxed{\frac{\partial v}{\partial x_j} = \frac{u^2(u^2 + 3)}{(u^2 + 1)^2} \frac{\partial u}{\partial x_j}}$$
	
	\textbf{(b) Análise da função $w(x)$}:
	Definimos a função externa $H: \mathbb{R} \to \mathbb{R}$ por $H(t) = 1 - \cos t$. Verificamos que $H(0) = 1 - \cos(0) = 1 - 1 = 0$. A função $H$ é de classe $C^\infty$ e sua derivada é:
	$$H'(t) = \sin t$$
	Observamos que $|H'(t)| = |\sin t| \le 1$ para todo $t \in \mathbb{R}$. Portanto, $H$ é uma função Lipschitziana com constante $M=1$. Novamente, aplicando o resultado da Questão 108, como $u \in W^{1,p}(U)$, $H(0)=0$ e $H'$ é limitada, a composição $w = H \circ u$ pertence a $W^{1,p}(U)$. A derivada fraca é dada por:
	$$\boxed{\frac{\partial w}{\partial x_j} = \sin(u) \frac{\partial u}{\partial x_j}}$$
}

\sintese{
	A resolução de ambos os itens baseou-se na verificação das hipóteses do teorema de composição em espaços de Sobolev. O ponto central foi a demonstração de que as funções externas $G$ e $H$ possuem derivadas limitadas, o que as torna globalmente Lipschitzianas. Uma vez estabelecida essa propriedade e a anulação na origem, a pertinência ao espaço $W^{1,p}$ e a forma da derivada fraca seguem automaticamente da regra da cadeia para funções de Sobolev.
	
	O fluxo lógico foi:
	$$\boxed{\text{Definição de } G, H \to \text{Cálculo de } G', H' \to \text{Prova de Limitação} \to \text{Aplicação de } W^{1,p} \text{ Comp.}}$$
}

\aprofundamento{
	A exigência $G(0)=0$ é fundamental para garantir que $v \in L^p(U)$. Se $G(0) \neq 0$, a função constante $G(0)$ teria que pertencer a $L^p(U)$, o que só ocorreria se o domínio $U$ fosse de medida finita. No caso de domínios ilimitados, a condição $G(0)=0$ assegura a integrabilidade de $v$ através da estimativa $|G(u)| \le M|u|$. Este resultado é um exemplo de como a regularidade de funções em espaços de Sobolev é preservada por transformações não lineares, desde que estas não \textbf{estiquem} a função de forma a romper a integrabilidade da derivada.
}

%%%% QUESTÃO 112 %%%%
\questao{
	Suponha que $u \in W^{1,1}(\mathbb{R}^n)$ e $c > 0$ são tais que $|\nabla u(x)| \le c|u(x)|$, quase todo ponto em $\mathbb{R}^n$. Mostre que $u \in W^{1,q}(\mathbb{R}^n)$, para todo $q > 1$.
}

\analise{
	\textbf{Análise do Enunciado}: A questão solicita a prova de que a condição $|\nabla u| \le c|u|$ força a função $u$ a pertencer a todos os espaços de Sobolev $W^{1,q}(\mathbb{R}^n)$ para $q > 1$, partindo da hipótese mínima de que $u \in W^{1,1}(\mathbb{R}^n)$. Trata-se de um problema de ganho de integrabilidade iterativo.
	
	\textbf{Fundamentação Teórica}: A base da prova é a \textbf{Desigualdade de Sobolev} (Gagliardo-Nirenberg-Sobolev). Para $1 \le p < n$, se $u \in W^{1,p}(\mathbb{R}^n)$, então $u \in L^{p^*}(\mathbb{R}^n)$ com $p^* = \frac{np}{n-p}$. A estratégia consiste em utilizar a hipótese de controle do gradiente para elevar o índice de integrabilidade $p$ em passos sucessivos. 
	
	\textbf{Estratégias}:
	\begin{enumerate}
		\item \textbf{Passo Inicial}: Utilizar a imersão $W^{1,1}(\mathbb{R}^n) \hookrightarrow L^{p_1}(\mathbb{R}^n)$ para encontrar o primeiro expoente de integrabilidade $p_1 = \frac{n}{n-1}$.
		\item \textbf{Iteração (Bootstrapping)}: Mostrar que se $u \in L^p(\mathbb{R}^n)$, então a condição $|\nabla u| \le c|u|$ implica que $\nabla u \in L^p(\mathbb{R}^n)$, logo $u \in W^{1,p}(\mathbb{R}^n)$.
		\item \textbf{Progressão do Expoente}: Aplicar a desigualdade de Sobolev novamente para mostrar que $u \in L^{p^*}$, onde $p^* > p$.
		\item \textbf{Convergência para $q$}: Demonstrar que a sequência de expoentes $p_k$ cresce indefinidamente, cobrindo qualquer $q \in (1, \infty)$.
	\end{enumerate}
}

\solucao{
	Iniciamos a análise partindo da hipótese $u \in W^{1,1}(\mathbb{R}^n)$. Pela \textbf{Desigualdade de Gagliardo-Nirenberg-Sobolev} para $p=1$ e $n > 1$, sabemos que existe uma constante $C_S$ tal que:
	$$\|u\|_{L^{p_1}(\mathbb{R}^n)} \le C_S \|\nabla u\|_{L^1(\mathbb{R}^n)}, \quad \text{onde } p_1 = \frac{n}{n-1}$$
	Como $u \in W^{1,1}(\mathbb{R}^n)$, temos que $u \in L^{p_1}(\mathbb{R}^n)$.
	
	Agora, utilizamos a hipótese fundamental $|\nabla u(x)| \le c|u(x)|$. Elevando ambos os lados à potência $p_1$, obtemos:
	$$|\nabla u(x)|^{p_1} \le c^{p_1} |u(x)|^{p_1}$$
	Integrando sobre $\mathbb{R}^n$:
	$$\int_{\mathbb{R}^n} |\nabla u|^{p_1} \, dx \le c^{p_1} \int_{\mathbb{R}^n} |u|^{p_1} \, dx \implies \|\nabla u\|_{L^{p_1}(\mathbb{R}^n)} \le c \|u\|_{L^{p_1}(\mathbb{R}^n)}$$
	Visto que $u \in L^{p_1}(\mathbb{R}^n)$ e $\nabla u \in L^{p_1}(\mathbb{R}^n)$, concluímos que $u \in W^{1,p_1}(\mathbb{R}^n)$.
	
	Podemos agora definir um processo iterativo para a sequência de expoentes $\{p_k\}$. Seja $p_k$ o expoente de integrabilidade alcançado no passo $k$, com $p_1 = \frac{n}{n-1}$. Se $p_k < n$, a imersão de Sobolev garante que:
	$$u \in W^{1,p_k}(\mathbb{R}^n) \implies u \in L^{p_{k+1}}(\mathbb{R}^n), \quad \text{com } p_{k+1} = \frac{np_k}{n-p_k}$$
	Novamente, a condição $|\nabla u| \le c|u|$ implica que:
	$$\|\nabla u\|_{L^{p_{k+1}}(\mathbb{R}^n)} \le c \|u\|_{L^{p_{k+1}}(\mathbb{R}^n)} < \infty$$
	Logo, $u \in W^{1, p_{k+1}}(\mathbb{R}^n)$. Analisando a sequência de expoentes:
	$$p_{k+1} = \frac{n p_k}{n - p_k} = \frac{p_k}{1 - \frac{p_k}{n}}$$
	Como $p_k < n$, temos $1 - \frac{p_k}{n} < 1$, o que implica $p_{k+1} > p_k$. A sequência $\{p_k\}$ é estritamente crescente. Se $p_k$ convergisse para um limite $L$, teríamos $L = \frac{nL}{n-L}$, o que implicaria $n-L = n \implies L=0$ (impossível, pois $p_1 > 1$) ou $L = \infty$. Portanto, $p_k \to \infty$ quando $k \to \infty$.
	
	Se em algum passo $p_k \ge n$, a imersão de Sobolev garante que $u \in L^q(\mathbb{R}^n)$ para todo $q \in [p_k, \infty)$. Independentemente do caso, para qualquer $q > 1$ dado, existe um $k$ tal que $p_k \ge q$. Como $u \in W^{1,p_k}(\mathbb{R}^n)$, segue que $u \in L^q(\mathbb{R}^n)$ e, pela condição do gradiente, $\nabla u \in L^q(\mathbb{R}^n)$.
	
	Concluímos, portanto, que:
	$$\boxed{u \in W^{1,q}(\mathbb{R}^n), \quad \forall q \in (1, \infty)}$$
}

\sintese{
	A prova demonstra o fenômeno de \textbf{bootstrapping}: a integrabilidade da função e da sua derivada estão acopladas. A imersão de Sobolev eleva a integrabilidade da função, e a condição $|\nabla u| \le c|u|$ \textbf{arrasta} a derivada para esse novo espaço, criando um ciclo de ganho de regularidade que converge para qualquer $q < \infty$.
	
	O fluxo lógico foi:
	$$\boxed{u \in W^{1,1} \to L^{p_1} \to W^{1,p_1} \to L^{p_2} \to \dots \to W^{1,q} \text{ para todo } q > 1}$$
}

\aprofundamento{
	Este resultado é intimamente ligado às estimativas de decaimento exponencial de funções em $L^p$. Na verdade, a condição $|\nabla u| \le c|u|$ sugere que $u$ se comporta, no máximo, como uma exponencial $e^{cx}$. Em domínios limitados, isso é trivial, mas em $\mathbb{R}^n$, a exigência de que $u \in L^1$ inicialmente força a função a decair rapidamente no infinito, permitindo que o processo de bootstrapping funcione. Se $u$ não estivesse em $L^1$, a função $u(x) = e^{cx}$ satisfaria a condição do gradiente, mas não pertenceria a nenhum $L^q(\mathbb{R}^n)$.
}


\end{document}
	
	\newpage
	% =============================================================================
	% BIBLIOGRAFIA
	% =============================================================================
	% Esta linha adiciona a Bibliografia ao Sumário (ToC)
	\addcontentsline{toc}{section}{Referências Bibliográficas} 
	\printbibliography[title={Referências Bibliográficas}]
	
\end{document}